% Generated mechanically from frozen input p07/ko/curves.tex.
% Do not edit this generated file; edit the bound locale source instead.

% OK, start here.
%

\setcounter{chapter}{52}
\chapter{대수곡선}



\phantomsection
\label{curves-section-phantom}


\section{서론}
\label{curves-section-introduction}

\noindent
이 장에서는 대수곡선 이론의 일부를 전개한다. 복소수체 위의 대수곡선을
다루는 참고문헌으로는 \cite{ACGH}가 있다.

\medskip\noindent
이미 알고 있는 내용. 일반적인 대수기하학뿐 아니라 대수곡선에 관한 몇 가지
구체적인 결과도 이미 증명하였다. 그 목록은 다음과 같다.
\begin{enumerate}
\item 다양체, 절 \ref{varieties-section-dimension-one}에서
$1$차원 Noether 스킴의 아핀 열린집합과 풍부한 가역층을 논하였다.
\item 다양체, 절 \ref{varieties-section-curves}에서 곡선은 아핀이거나
사영임을 보았다.
\item 다양체, 절 \ref{varieties-section-divisors-curves}에서
고유 곡선 위의 국소 자유 가군의 차수를 논하였다.
\item 곡선의 Picard 스킴, 절 \ref{pic-section-introduction}에서
대수적으로 닫힌 체 위의 비특이 사영곡선의 Picard 스킴을 논하였다.
\end{enumerate}





\section{곡선과 함수체}
\label{curves-section-curves-function-fields}

\noindent
이 절에서는 곡선의 경우에 다양체, 절
\ref{varieties-section-varieties-rational-maps}의 결과를 더 자세히 전개한다.

\begin{lemma}
\label{curves-lemma-extend-over-dvr}
$k$를 체라 하자. $X$를 곡선, $Y$를 고유 다양체라 하자.
$U \subset X$를 공집합이 아닌 열린집합으로 두고 $f : U \to Y$를 사상이라 하자.
$x \in X$가 닫힌점이고 $\mathcal{O}_{X, x}$가 이산값매김환이면, $x$를 포함하는
열린집합 $U \subset U' \subset X$와 $f$를 연장하는 다양체의 사상
$f' : U' \to Y$가 존재한다.
\end{lemma}

\begin{proof}
이는 스킴의 사상, 보조정리 \ref{morphisms-lemma-extend-across}의 특수한 경우이다.
\end{proof}

\begin{lemma}
\label{curves-lemma-extend-over-normal-curve}
$k$를 체라 하자. $X$를 정규 곡선, $Y$를 고유 다양체라 하자.
$X$에서 $Y$로 가는 유리사상들의 집합은 사상 $X \to Y$들의 집합과 같다.
\end{lemma}

\begin{proof}
모든 국소환이 이산값매김환이므로(예를 들어 다양체, 보조정리
\ref{varieties-lemma-regular-point-on-curve} 참조), $X$에서 $Y$로 가는 유리사상은
보조정리 \ref{curves-lemma-extend-over-dvr}에 의해 사상 $X \to Y$로 연장된다.
반대로 두 사상 $f,g: X \to Y$가 유리사상으로서 동치이면,
스킴의 사상, 보조정리 \ref{morphisms-lemma-equality-of-morphisms}에 의해 $f=g$이다.
\end{proof}

\begin{lemma}
\label{curves-lemma-flat}
$k$를 체라 하자. $f : X \to Y$를 $k$ 위 곡선들의 상수가 아닌 사상이라 하자.
$Y$가 정규이면 $f$는 평탄하다.
\end{lemma}

\begin{proof}
$y \in Y$로 가는 $x \in X$를 택하자. 그러면 $\mathcal{O}_{Y,y}$는 체이거나
이산값매김환이다(다양체, 보조정리
\ref{varieties-lemma-regular-point-on-curve}). $f$는 상수가 아니므로 우세이다
($X$의 일반점을 $Y$의 일반점으로 보내야 하기 때문이다). 따라서
$\mathcal{O}_{Y,y}\to\mathcal{O}_{X,x}$는 단사이다
(스킴의 사상, 보조정리 \ref{morphisms-lemma-dominant-between-integral}).
이에 따라 $\mathcal{O}_{X,x}$는 $\mathcal{O}_{Y,y}$-가군으로서 꼬임이 없고,
따라서 대수학 심화, 보조정리
\ref{more-algebra-lemma-valuation-ring-torsion-free-flat}에 의해
$\mathcal{O}_{X,x}$는 $\mathcal{O}_{Y,y}$-가군으로서 평탄하다.
\end{proof}

\begin{lemma}
\label{curves-lemma-finite}
$k$를 체라 하자. $f:X\to Y$를 $k$ 위 스킴들의 사상이라 하자. 다음을 가정하자.
\begin{enumerate}
\item $Y$는 $k$ 위에서 분리되어 있다.
\item $X$는 $k$ 위에서 고유이고 차원이 $\leq 1$이다.
\item 차원이 $1$인 모든 기약성분 $Z\subset X$에 대하여 $f(Z)$는 적어도 두 점을 가진다.
\end{enumerate}
그러면 $f$는 유한이다.
\end{lemma}

\begin{proof}
스킴의 사상, 보조정리 \ref{morphisms-lemma-image-proper-scheme-closed}에 의해
사상 $f$는 고유하다. 따라서 $f(X)$는 닫혀 있고 닫힌점의 상도 닫힌점이다.
$X$의 한 닫힌점의 상인 $y\in Y$를 잡자. 조건 (3)에 의해
$f^{-1}(\{y\})$는 $X$의 차원 $1$인 기약성분들의 어느 일반점도 포함하지 않는
닫힌 부분집합이다. 따라서 $f^{-1}(\{y\})$는 유한하다. 그러므로 스킴의 사상 심화, 보조정리
\ref{more-morphisms-lemma-proper-finite-fibre-finite-in-neighbourhood}에 의해
$f$는 $y$의 어떤 열린 근방 위에서 유한이다($Y$가 Noether이면 더 쉬운
스킴의 코호몰로지, 보조정리
\ref{coherent-lemma-proper-finite-fibre-finite-in-neighbourhood}를 쓸 수 있다).
위에서 이러한 점 $y$가 충분히 많음을 보았으므로 증명이 끝난다.
\end{proof}

\begin{lemma}
\label{curves-lemma-extend-to-completion}
$k$를 체라 하자. $X\to Y$를 다양체의 사상이라 하고, $Y$는 고유이며 $X$는
곡선이라 하자. $X\to\overline{X}$가 열린몰입이고 $\overline{X}$가 사영곡선인
분해 $X\to\overline{X}\to Y$가 존재한다.
\end{lemma}

\begin{proof}
보조정리 \ref{curves-lemma-extend-over-dvr}와 다양체, 보조정리
\ref{varieties-lemma-reduced-dim-1-projective-completion}에서 바로 따른다.
\end{proof}

\noindent
다음은 이 절의 주정리이다. 다양체의 사상 $f:X\to Y$의 상 $f(X)$가 $Y$의
한 점 $y$로만 이루어지면 이 사상을 {\it 상수}라고 하겠다. 이때 $y$는 $Y$의
닫힌점이다($X$의 닫힌점의 상이 $Y$의 닫힌점이기 때문이다).

\begin{theorem}
\label{curves-theorem-curves-rational-maps}
$k$를 체라 하자. 다음 범주들은 표준적으로 동치이다.
\begin{enumerate}
\item 초월차수가 $1$인 유한 생성 체확장 $K/k$들의 범주.
\item 곡선들과 우세 유리사상들의 범주.
\item 정규 사영곡선들과 상수가 아닌 사상들의 범주.
\item 비특이 사영곡선들과 상수가 아닌 사상들의 범주.
\item 정칙 사영곡선들과 상수가 아닌 사상들의 범주.
\item 정규 고유 곡선들과 상수가 아닌 사상들의 범주.
\end{enumerate}
\end{theorem}

\begin{proof}
범주 (1)과 (2)의 동치는 다양체, 정리
\ref{varieties-theorem-varieties-rational-maps}의 동치를 제한한 것이다.
구체적으로 다양체가 곡선일 필요충분조건은 그 함수체의 초월차수가 $1$인 것이다.
예를 들어 다양체, 보조정리
\ref{varieties-lemma-dimension-locally-algebraic}를 보라.

\medskip\noindent
(3), (4), (5), (6)의 범주들은 같다. 먼저 ``정칙''과 ``비특이''는 동의어이다.
스킴의 성질, 정의 \ref{properties-definition-regular}을 보라. Noether인
$1$차원 스킴에서는 정규성과 정칙성이 같은 조건이다(스킴의 성질, 보조정리
\ref{properties-lemma-regular-normal} 및
\ref{properties-lemma-normal-dimension-1-regular}). 곡선의 경우에는
다양체, 보조정리 \ref{varieties-lemma-regular-point-on-curve}를 보라.
따라서 (3)은 (5)와 같다. 마지막으로 다양체, 보조정리
\ref{varieties-lemma-dim-1-proper-projective}에 의해 (6)은 (3)과 같다.

\medskip\noindent
$f:X\to Y$가 비특이 사영곡선들의 상수가 아닌 사상이면, $f$는 $X$의 일반점
$\eta$를 $Y$의 일반점 $\xi$로 보낸다. 따라서 범주 (1)에서의 사상
$k(Y)=\mathcal{O}_{Y,\xi}\to\mathcal{O}_{X,\eta}=k(X)$를 얻는다. 두 사상
$f,g:X\to Y$가 같은 사상 $k(Y)\to k(X)$를 준다면, (1)과 (2)의 동치에 의해
$f$와 $g$는 유리사상으로서 동치이고, 따라서 보조정리
\ref{curves-lemma-extend-over-normal-curve}에 의해 $f=g$이다. 반대로 범주 (1)에서
사상 $k(Y)\to k(X)$가 주어졌다고 하자. 그러면 어떤 공집합이 아닌 열린집합
$U\subset X$에 대하여 사상 $U\to Y$를 얻는다. 보조정리
\ref{curves-lemma-extend-over-dvr}에 의해 이는 $X$ 전체로 연장되며, 범주 (5)의 사상을
얻는다. 따라서 완전충실 함자 (5)$\to$(1)이 있음을 알 수 있다.

\medskip\noindent
증명을 마치려면 (1)의 모든 $K/k$가 어떤 정규 사영곡선의 함수체임을 보여야 한다.
이미 어떤 곡선 $X$에 대하여 $K=k(X)$임을 안다. $X$를 그 정규화로 바꾸면
($X$와 쌍유리인 다양체이다) $X$가 정규라고 가정할 수 있다
(다양체, 보조정리
\ref{varieties-lemma-normalization-locally-algebraic}). 이제 다양체, 보조정리
\ref{varieties-lemma-reduced-dim-1-projective-completion}에서와 같이
$\overline{X}\setminus X=\{x_1,\ldots,x_n\}$인 $X\to\overline{X}$를 택한다.
$X$가 정규이고 각 국소환 $\mathcal{O}_{\overline{X},x_i}$도 정규이므로
$\overline{X}$는 원하는 정규 사영곡선이다. (비고: 먼저 다양체, 보조정리
\ref{varieties-lemma-dim-1-projective-completion}을 써서 콤팩트화한 뒤
다양체, 보조정리 \ref{varieties-lemma-normalization-locally-algebraic}를 써서
정규화해도 된다. 이렇게 하면 다소 까다로운 스킴의 사상, 보조정리
\ref{morphisms-lemma-relative-normalization-normal-codim-1}을 피할 수 있다.)
\end{proof}

\begin{definition}
\label{curves-definition-normal-projective-model}
$k$를 체, $X$를 곡선이라 하자. $X$의 {\it 비특이 사영모형}이란 $Y$가
비특이 사영곡선이고 $\varphi:k(X)\to k(Y)$가 함수체의 동형인 순서쌍
$(Y,\varphi)$를 말한다.
\end{definition}

\noindent
정리 \ref{curves-theorem-curves-rational-maps}에 의해 비특이 사영모형은 유일한
동형을 제외하고 결정된다. 따라서 흔히 ``그 비특이 사영모형''이라고 말한다.
보통 표기에서 $\varphi$를 생략한다. 주의: $Y$가 $k$ 위에서 매끄럽지 않을 수도
있지만, 보조정리 \ref{curves-lemma-nonsingular-model-smooth}가 보여 주듯 이는 양의 표수에서만
일어날 수 있다.

\begin{lemma}
\label{curves-lemma-nonsingular-model-smooth}
$k$를 체, $X$를 곡선, $Y$를 $X$의 비특이 사영모형이라 하자. $k$가 완전체이면
$Y$는 매끄러운 사영곡선이다.
\end{lemma}

\begin{proof}
예를 들어 다양체, 보조정리 \ref{varieties-lemma-regular-point-on-curve}를 보라.
\end{proof}

\begin{lemma}
\label{curves-lemma-smooth-models}
$k$를 체라 하고 $X$를 $k$ 위의 기하적으로 기약인 곡선이라 하자. 체확장
$K/k$에 대하여 $Y_K$로 $(X_K)_{red}$의 비특이 사영모형을 나타내자.
\begin{enumerate}
\item $X$가 고유이면 $Y_K$는 $X_K$의 정규화이다.
\item $Y_K$가 매끄럽게 되는 유한 순수 비분리 확장 $K/k$가 존재한다.
\item $Y_K$가 매끄러울 때마다\footnote{사실 기하적으로 환원이기만 해도 된다.}
$H^0(Y_K,\mathcal{O}_{Y_K})=K$이다.
\item 다음 가환 그림을 주자.
$$
\xymatrix{
\Omega & K' \ar[l] \\
K \ar[u] & k \ar[l] \ar[u]
}
$$
$Y_K$와 $Y_{K'}$가 매끄러우면
$Y_\Omega=(Y_K)_\Omega=(Y_{K'})_\Omega$이다.
\end{enumerate}
\end{lemma}

\begin{proof}
$X'$을 $X$의 비특이 사영모형이라 하자. 그러면 $X'$과 $X$는 서로 동형인
공집합이 아닌 열린 부분스킴들을 가진다. 특히 $X$가 기하적으로 기약이므로
$X'$도 그러하다(일부 세부사항은 생략한다). 따라서 $X$가 사영이라고 가정해도 된다.

\medskip\noindent
$X$가 고유라고 하자. 그러면 $X_K$가 고유이고, 그 정규화 $(X_K)^\nu$는
고유 스킴 위의 유한 스킴이므로 고유이다(다양체, 보조정리
\ref{varieties-lemma-normalization-locally-algebraic} 및 스킴의 사상, 보조정리
\ref{morphisms-lemma-finite-proper},
\ref{morphisms-lemma-composition-proper}). 한편 $X$가 기하적으로 기약이므로
$X_K$는 기약이다. 따라서 $X_K^\nu$는 고유이고 정규이며 기약이고
$(X_K)_{red}$와 쌍유리이다. 고유 곡선은 사영이므로(다양체, 보조정리
\ref{varieties-lemma-dim-1-proper-projective}) 이는 (1)을 증명한다.

\medskip\noindent
(2)의 증명. $X$가 고유이고 (1)이 성립하므로 다양체, 보조정리
\ref{varieties-lemma-finite-extension-geometrically-normal}을 적용하여 $Y_K$가
기하적으로 정규가 되는 유한 순수 비분리 확장 $K/k$를 찾을 수 있다.
곡선에서는 정규성과 정칙성이 같으므로(스킴의 성질, 보조정리
\ref{properties-lemma-normal-dimension-1-regular}) $Y_K$는 기하적으로 정칙이다.
따라서 다양체, 보조정리
\ref{varieties-lemma-geometrically-regular-smooth}에 의해 $Y$는 매끄러운 다양체이다.

\medskip\noindent
$Y_K$가 기하적으로 환원이면 $Y_K$는 기하적으로 정역이다(다양체, 보조정리
\ref{varieties-lemma-geometrically-integral}). 그러면 다양체, 보조정리
\ref{varieties-lemma-regular-functions-proper-variety}에 의해
$H^0(Y_K,\mathcal{O}_{Y_K})=K$임을 알 수 있다. 매끄러운 다양체는 기하적으로
환원이므로(실제로 기하적으로 정칙이다. 다양체, 보조정리
\ref{varieties-lemma-geometrically-regular-smooth} 참조) 이는 (3)을 증명한다.

\medskip\noindent
$Y_K$가 매끄러우면 모든 확장 $\Omega/K$에 대하여 밑변환 $(Y_K)_\Omega$가
$\Omega$ 위에서 매끄럽다(스킴의 사상, 보조정리
\ref{morphisms-lemma-base-change-smooth}). 따라서
$Y_\Omega=(Y_K)_\Omega$임이 분명하다. 이는 (4)를 증명한다.
\end{proof}






\section{선형계}
\label{curves-section-linear-series}

\noindent
고전적인 서술(주의 \ref{curves-remark-classical-linear-series} 참조)과 달리
여기서는 선형계를 다음과 같이 정의한다.

\begin{definition}
\label{curves-definition-linear-series}
$k$를 체라 하고 $X$를 $k$ 위의 차원 $\leq1$인 고유 스킴이라 하자.
$d\geq0$, $r\geq0$라 하자. {\it 차수 $d$, 차원 $r$인 선형계}란
$\mathcal{L}$이 차수 $d$인 가역 $\mathcal{O}_X$-가군이고
(다양체, 정의 \ref{varieties-definition-degree-invertible-sheaf})
$V\subset H^0(X,\mathcal{L})$가 차원 $r+1$인 $k$-부분벡터공간인 순서쌍
$(\mathcal{L},V)$를 말한다. 이를 줄여서 $(\mathcal{L},V)$는 $X$ 위의
{\it $\mathfrak g^r_d$}라고 하겠다.
\end{definition}

\noindent
주로 $X$가 비특이 고유 곡선일 때 이를 사용한다. 사실 위 정의는 고전적인
$\mathfrak g^r_d$의 정의를 일반화하는 한 가지 방법일 뿐이다. 예를 들어
$X$가 고유 곡선이면 $\mathcal{L}$이 계수 $1$인 꼬임 없는 연접
$\mathcal{O}_X$-가군인 경우까지 허용하여 선형계를 일반화할 수 있다.
비특이 곡선에서는 모든 꼬임 없는 연접 가군이 국소 자유이므로, 비특이
고유 곡선의 경우에는 이 정의가 우리의 개념과 일치한다.

\medskip\noindent
다음 보조정리는 고유 비특이 곡선 위의 선형계가 갖는 기하적 의미를 설명한다.

\begin{lemma}
\label{curves-lemma-linear-series}
$k$를 체라 하고 $X$를 $k$ 위의 비특이 고유 곡선이라 하자.
$(\mathcal{L},V)$를 $X$ 위의 $\mathfrak g^r_d$라 하자. 그러면 $k$ 위 다양체의 사상
$$
\varphi:X\longrightarrow\mathbf{P}^r_k=\text{Proj}(k[T_0,\ldots,T_r])
$$
과 사상 $\alpha:\varphi^*\mathcal{O}_{\mathbf{P}^r_k}(1)\to\mathcal{L}$가
존재하여, $\alpha$는 $\varphi^*T_0,\ldots,\varphi^*T_r$을 $V$의 한 기저로 보낸다.
\end{lemma}

\begin{proof}
$s_0,\ldots,s_r\in V$를 $k$-기저라 하자. $X$가 비특이이므로 사상
$s_0,\ldots,s_r:\mathcal{O}_X^{\oplus r+1}\to\mathcal{L}$의 상
$\mathcal{L}'\subset\mathcal{L}$은 가역 $\mathcal{O}_X$-가군이다. 예를 들어
인자, 보조정리 \ref{divisors-lemma-torsion-free-over-regular-dim-1}에서 따른다.
이제 구성, 보조정리 \ref{constructions-lemma-projective-space}를 적용하면
보조정리의 명제에 나오는 사상
$$
\varphi=\varphi_{(\mathcal{L}',(s_0,\ldots,s_r))}:X\longrightarrow\mathbf{P}^r_k
$$
를 얻는다.
\end{proof}

\begin{lemma}
\label{curves-lemma-linear-series-trivial-existence}
$k$를 체라 하고 $X$를 $k$ 위의 차원 $\leq1$인 고유 스킴이라 하자.
$X$가 $\mathfrak g^r_d$를 가지면 모든 $0\leq s\leq r$에 대하여
$X$는 $\mathfrak g^s_d$도 가진다.
\end{lemma}

\begin{proof}
$k$ 위의 차원 $r+1$인 벡터공간 $V$는 모든 $0\leq s\leq r$에 대하여
차원 $s+1$인 선형부분공간을 가지므로 성립한다.
\end{proof}

\begin{lemma}
\label{curves-lemma-g1d}
$k$를 체라 하고 $X$를 $k$ 위의 비특이 고유 곡선이라 하자.
$(\mathcal{L},V)$를 $X$ 위의 $\mathfrak g^1_d$라 하자. 보조정리
\ref{curves-lemma-linear-series}의 사상 $\varphi:X\to\mathbf{P}^1_k$는 다음 중 하나이다.
\begin{enumerate}
\item 상수가 아니고 차수가 $\leq d$이다.
\item $\mathbf{P}^1_k$의 한 닫힌점을 통하여 분해되며, 이 경우
$H^0(X,\mathcal{O}_X)\not = k$이다.
\end{enumerate}
\end{lemma}

\begin{proof}
보조정리 \ref{curves-lemma-linear-series}에 의해
$\mathcal{L}'=\varphi^*\mathcal{O}_{\mathbf{P}^1_k}(1)$에 대하여 영이 아닌 사상
$\mathcal{L}'\to\mathcal{L}$가 있다. 따라서 다양체, 보조정리
\ref{varieties-lemma-check-invertible-sheaf-trivial}에 의해
$0\leq\deg(\mathcal{L}')\leq d$이다. $\deg(\mathcal{L}')=0$이면 같은 보조정리에서
$\mathcal{L}'\cong\mathcal{O}_X$를 얻고, 선형독립인 단면이 두 개 있으므로 (2)의
경우이다. $\deg(\mathcal{L}')>0$이면 $\varphi$는 상수가 아니다(상수사상에 의한
가역 가군의 당김은 자명하기 때문이다). 따라서 다양체, 보조정리
\ref{varieties-lemma-degree-pullback-map-proper-curves}에 의해
$$
\deg(\mathcal{L}')=\deg(X/\mathbf{P}^1_k)\deg(\mathcal{O}_{\mathbf{P}^1_k}(1)).
$$
$\mathcal{O}_{\mathbf{P}^1_k}(1)$의 차수는 $1$이므로 증명이 끝난다.
\end{proof}

\begin{lemma}
\label{curves-lemma-grd-inequalities}
$k$를 체라 하고 $X$를 $H^0(X,\mathcal{O}_X)=k$인 $k$ 위의 고유 곡선이라 하자.
$X$가 $\mathfrak g^r_d$를 가지면 $r\leq d$이다. 등호가 성립하면
$H^1(X,\mathcal{O}_X)=0$, 즉 $X$의 종수(정의 \ref{curves-definition-genus})는 $0$이다.
\end{lemma}

\begin{proof}
$(\mathcal{L},V)$를 $\mathfrak g^r_d$라 하자. 이렇게 해도 $r$만 커질 수 있으므로
$V=H^0(X,\mathcal{L})$이라고 가정해도 된다. 영이 아닌 $s\in V$를 택하자.
그러면 $s$의 영점 스킴은 유효 Cartier 인자 $D\subset X$이고,
$\mathcal{L}=\mathcal{O}_X(D)$이며 짧은 완전열
$$
0\to\mathcal{O}_X\to\mathcal{L}\to\mathcal{L}|_D\to0
$$
을 얻는다. 인자, 보조정리 \ref{divisors-lemma-characterize-OD}와 주의
\ref{divisors-remark-ses-regular-section}를 보라. 다양체, 보조정리
\ref{varieties-lemma-degree-effective-Cartier-divisor}에 의해
$\deg(D)=\deg(\mathcal{L})=d$이다. $D$는 Artin 스킴이므로
$\mathcal{L}|_D\cong\mathcal{O}_D$이다.\footnote{우리의 경우에는
$D\to\Spec(k)$가 유한이므로 인자, 보조정리
\ref{divisors-lemma-finite-trivialize-invertible-upstairs}에서 따른다.}
따라서
$$
\dim_kH^0(D,\mathcal{L}|_D)=\dim_kH^0(D,\mathcal{O}_D)=\deg(D)=d.
$$
한편 가정에 의해 $\dim_kH^0(X,\mathcal{O}_X)=1$이고
$\dim H^0(X,\mathcal{L})=r+1$이다. 따라서 $r+1\leq1+d$, 즉 $r\leq d$이다.

\medskip\noindent
등호가 성립한다고 하자. 그러면
$H^0(X,\mathcal{L})\to H^0(X,\mathcal{L}|_D)$는 전사이다.
$H^1(X,\mathcal{L})$가 영임을 안다면 긴 완전 코호몰로지 열로부터
$H^1(X,\mathcal{O}_X)$가 영이라고 결론 내려 증명이 끝난다. 우리의 전략은
$\mathcal{L}$을 소멸성이 있는 큰 거듭제곱으로 바꾸는 것이다.
$\mathcal{L}|_D$는 자명한 가역 가군이므로(위 참조), $D$에 제한했을 때
$\mathcal{L}|_D$를 생성하는 $\mathcal{L}$의 단면 $t$를 찾을 수 있다. 곱셈사상
$$
\mu:H^0(X,\mathcal{L})\otimes_kH^0(X,\mathcal{L})
\longrightarrow H^0(X,\mathcal{L}^{\otimes2})
$$
과 짧은 완전열
$$
0\to\mathcal{L}\xrightarrow{s}\mathcal{L}^{\otimes2}
\to\mathcal{L}^{\otimes2}|_D\to0
$$
을 생각하자. $H^0(\mathcal{L})\to H^0(\mathcal{L}|_D)$가 전사이고 $t$가
$\mathcal{L}|_D$의 한 자명화로 가므로,
$\mu(H^0(X,\mathcal{L})\otimes t)$는 $H^0(X,\mathcal{L}^{\otimes2})$의
부분공간으로서 $\mathcal{L}^{\otimes2}|_D$의 전역단면 위로 전사한다. 따라서
$$
\dim H^0(X,\mathcal{L}^{\otimes2})=r+1+d=2r+1
=\deg(\mathcal{L}^{\otimes2})+1.
$$
즉 $\mathcal{L}^{\otimes2}$도 $\mathcal{L}$과 같은 성질, 곧 전역단면 공간의
차원이 차수보다 하나 큰 성질을 가진다. $\mathcal{L}$은 풍부하므로
(다양체, 보조정리 \ref{varieties-lemma-ample-curve}), 어떤 $n_0$가 존재하여
모든 $n\geq n_0$에 대해 $\mathcal{L}^{\otimes n}$의 $H^1$은 소멸한다
(스킴의 코호몰로지, 보조정리
\ref{coherent-lemma-coherent-proper-ample}). 따라서 충분히 큰 $m$에 대하여
$n=2^m$인 $\mathcal{L}^{\otimes n}$에 위 논증을 적용하면 보조정리가 성립한다.
\end{proof}

\begin{remark}[고전적 정의]
\label{curves-remark-classical-linear-series}
$X$를 대수적으로 닫힌 체 $k$ 위의 매끄러운 사영곡선이라 하자. 두 유효
Cartier 인자 $D,D'\subset X$가 $\mathcal{O}_X$-가군으로서
$\mathcal{O}_X(D)\cong\mathcal{O}_X(D')$일 필요충분조건으로 이들을
{\it 선형동치}라고 한다. $\Pic(X)=\text{Cl}(X)$이므로(인자, 보조정리
\ref{divisors-lemma-local-rings-UFD-c1-bijective}), $D$와 $D'$가 선형동치일
필요충분조건은 $D$와 $D'$에 대응하는 Weil 인자가 $\text{Cl}(X)$의 같은 원소를
정의하는 것이다. 차수 $d$인 유효 Cartier 인자 $D\subset X$가 주어졌을 때
$D$의 {\it 완전 선형계} $|D|$는 $D$와 선형동치인 유효 Cartier 인자
$E\subset X$들의 집합이다. 달리 말하면 $|D|$는 $\mathcal{O}_X(D)$에 대응하는
닫힌점 위에서 사상
$$
\gamma_d:\underline{\Hilbfunctor}^d_{X/k}
\longrightarrow\underline{\Picardfunctor}^d_{X/k}
$$
의 올에 있는 닫힌점들의 집합이다(곡선의 Picard 스킴, 보조정리
\ref{pic-lemma-picard-pieces}). 이로써 $|D|$에는 자연스러운 스킴 구조가
주어지며, $m+1=h^0(\mathcal{O}_X(D))$일 때
$|D|\cong\mathbf{P}^m_k$임을 알 수 있다. 더 표준적으로는
$$
|D|=\mathbf{P}(H^0(X,\mathcal{O}_X(D))^\vee)
$$
이다. 여기서 $(-)^\vee$는 $k$-선형 쌍대를 뜻하고 $\mathbf{P}$는
구성, 예 \ref{constructions-example-projective-space}에서와 같다.
이 용어로 $X$ 위의 {\it 선형계}란 선형방정식들로 잘라낼 수 있는 닫힌
부분다양체 $L\subset|D|$를 말한다. $L$의 차원이 $r$이면
$L=\mathbf{P}(V^\vee)$이고, 여기서
$V\subset H^0(X,\mathcal{O}_X(D))$는 차원 $r+1$인 선형부분공간이다.
따라서 고전적 선형계 $L\subset|D|$는 위에서 정의한 선형계
$(\mathcal{O}_X(D),V)$에 대응한다.
\end{remark}






\section{쌍대성}
\label{curves-section-duality}

\noindent
이 절에서는 체 위의 고유 $1$차원 스킴에 관하여, 쌍대화 복합체와 쌍대성에
대한 매우 일반적인 내용이 함의하는 바를 구체화한다. 체 위의 더 높은 차원
고유 스킴에 대한 유사한 논의는 스킴의 쌍대성, 절
\ref{duality-section-duality-proper-over-field}을 보라.

\begin{lemma}
\label{curves-lemma-duality-dim-1}
$X$를 체 $k$ 위의 차원 $\leq1$인 고유 스킴이라 하자. 다음 성질을 갖는
쌍대화 복합체 $\omega_X^\bullet$이 존재한다.
\begin{enumerate}
\item $H^i(\omega_X^\bullet)$는 $i=-1,0$일 때만 영이 아닐 수 있다.
\item $\omega_X=H^{-1}(\omega_X^\bullet)$는 지지가 차원 $1$인 기약성분들인
연접 Cohen--Macaulay 가군이다.
\item 닫힌점 $x\in X$에 대하여 가군 $H^0(\omega_{X,x}^\bullet)$가 영이 아닐
필요충분조건은 다음 중 하나가 성립하는 것이다.
\begin{enumerate}
\item $\dim(\mathcal{O}_{X,x})=0$이다.
\item $\dim(\mathcal{O}_{X,x})=1$이고 $\mathcal{O}_{X,x}$가 Cohen--Macaulay가 아니다.
\end{enumerate}
\item $K\in D_\QCoh(\mathcal{O}_X)$에 대하여 완전삼각형 및 이동과 양립하는
함자적 동형\footnote{Yoneda 보조정리에 의해 이 성질은
$D_\QCoh(\mathcal{O}_X)$에서 $\omega_X^\bullet$을 유일한 동형을 제외하고
특징짓는다. 실제로 $\omega_X^\bullet$은 $D^b_{\textit{Coh}}(\mathcal{O}_X)$의
대상이므로
$K\in D^b_{\textit{Coh}}(\mathcal{O}_X)$만 생각해도 충분하다.}
$$
\Ext^i_X(K,\omega_X^\bullet)=\Hom_k(H^{-i}(X,K),k)
$$
이 있다.
\item $X$ 위의 준연접층 $\mathcal{F}$에 대하여 함자적 동형
$\Hom(\mathcal{F},\omega_X)=\Hom_k(H^1(X,\mathcal{F}),k)$가 있다.
\item $X\to\Spec(k)$가 상대차원 $1$인 매끄러운 사상이면
$\omega_X\cong\Omega_{X/k}$이다.
\end{enumerate}
\end{lemma}

\begin{proof}
$f:X\to\Spec(k)$를 구조사상이라 하자. 스킴의 쌍대성, 주의
\ref{duality-remark-relative-dualizing-complex}에 기술된 상대 쌍대화 복합체
$$
\omega_X^\bullet=\omega_{X/k}^\bullet=a(\mathcal{O}_{\Spec(k)})
$$
로 시작한다. $a$는
$f_*:D_\QCoh(\mathcal{O}_X)\to D(\mathcal{O}_{\Spec(k)})$의 오른쪽 수반이므로
(4)는 구성상 성립한다. $f$가 고유이므로 정의에 의해
$f^!(\mathcal{O}_{\Spec(k)})=a(\mathcal{O}_{\Spec(k)})$이다. 스킴의 쌍대성,
절 \ref{duality-section-upper-shriek}을 보라. 따라서 $\omega_X^\bullet$과
$\omega_X$는 스킴의 쌍대성, 예
\ref{duality-example-proper-over-local} 및 예
\ref{duality-example-equidimensional-over-field}에서와 같다. (1), (2)는
스킴의 쌍대성, 보조정리 \ref{duality-lemma-vanishing-good-dualizing}에서 따른다.
닫힌점 $x\in X$에 대하여 $\omega_{X,x}^\bullet$은 $\mathcal{O}_{X,x}$ 위의
정규화된 쌍대화 복합체이다. 스킴의 쌍대성, 보조정리
\ref{duality-lemma-good-dualizing-normalized}를 보라. 그러면 (3)은 쌍대화 복합체, 보조정리 \ref{dualizing-lemma-apply-CM}에서 따른다. 연접 $\mathcal{F}$에
대해서 (5)는 스킴의 쌍대성, 보조정리
\ref{duality-lemma-dualizing-module-proper-over-A}에서 따르고, 일반적인 경우에는
$K=\mathcal{F}[0]$, $i=-1$로 두고 (4)를 풀어 쓰면 된다. (6)은 스킴의 쌍대성, 보조정리 \ref{duality-lemma-smooth-proper}에서 따른다.
\end{proof}

\begin{lemma}
\label{curves-lemma-duality-dim-1-CM}
$X$를 체 $k$ 위의 Cohen--Macaulay이고 순수차원 $1$인 고유 스킴이라 하자.
보조정리 \ref{curves-lemma-duality-dim-1}의 가군 $\omega_X$는 다음 성질을 가진다.
\begin{enumerate}
\item $\omega_X$는 $X$ 위의 쌍대화 가군이다(스킴의 쌍대성, 절
\ref{duality-section-dualizing-module}).
\item $\omega_X$는 지지가 $X$인 연접 Cohen--Macaulay 가군이다.
\item $K\in D_\QCoh(X)$에 대하여 이동과 양립하는 함자적 동형
$\Ext^i_X(K,\omega_X[1])=\Hom_k(H^{-i}(X,K),k)$가 있다.
\item $X$ 위의 준연접층 $\mathcal{F}$에 대하여 함자적 동형
$\Ext^{1+i}(\mathcal{F},\omega_X)=\Hom_k(H^{-i}(X,\mathcal{F}),k)$가 있다.
\end{enumerate}
\end{lemma}

\begin{proof}
보조정리 \ref{curves-lemma-duality-dim-1}의 증명에서 $\omega_X$는 스킴의 쌍대성,
예 \ref{duality-example-proper-over-local}의 것이므로 쌍대화 가군임을
상기하자. 다른 명제들은 보조정리 \ref{curves-lemma-duality-dim-1}과, $X$가
Cohen--Macaulay이므로 $\omega_X^\bullet=\omega_X[1]$이라는 사실(스킴의 쌍대성, 보조정리 \ref{duality-lemma-dualizing-module-CM-scheme})에서 따른다.
\end{proof}

\begin{remark}
\label{curves-remark-rework-duality-locally-free}
$X$를 체 $k$ 위의 차원 $\leq1$인 고유 스킴이라 하고 $\omega_X^\bullet$과
$\omega_X$를 보조정리 \ref{curves-lemma-duality-dim-1}에서와 같이 두자. $\mathcal{E}$가
쌍대 $\mathcal{E}^\vee$를 갖는 유한 국소 자유 $\mathcal{O}_X$-가군이면 표준동형
$$
\Hom_k(H^{-i}(X,\mathcal{E}),k)=
H^i(X,\mathcal{E}^\vee\otimes_{\mathcal{O}_X}^{\mathbf{L}}\omega_X^\bullet)
$$
이 있다. 이는 보조정리와 코호몰로지, 보조정리
\ref{cohomology-lemma-dual-perfect-complex}에서 따른다. $X$가 Cohen--Macaulay이고
순수차원 $1$이면 보조정리 \ref{curves-lemma-duality-dim-1-CM}에 의해 표준동형
$$
\Hom_k(H^{-i}(X,\mathcal{E}),k)=
H^{1+i}(X,\mathcal{E}^\vee\otimes_{\mathcal{O}_X}\omega_X)
$$
이 있다. 특히 $\mathcal{L}$이 가역 $\mathcal{O}_X$-가군이면
$$
\dim_kH^0(X,\mathcal{L})=
\dim_kH^1(X,\mathcal{L}^{\otimes-1}\otimes_{\mathcal{O}_X}\omega_X)
$$
이고
$$
\dim_kH^1(X,\mathcal{L})=
\dim_kH^0(X,\mathcal{L}^{\otimes-1}\otimes_{\mathcal{O}_X}\omega_X)
$$
이다.
\end{remark}

\noindent
다음은 쌍대화 복합체에 대한 건전성 점검이다.

\begin{lemma}
\label{curves-lemma-sanity-check-duality}
$X$를 체 $k$ 위의 차원 $\leq1$인 고유 스킴이라 하고 $\omega_X^\bullet$과
$\omega_X$를 보조정리 \ref{curves-lemma-duality-dim-1}에서와 같이 두자.
\begin{enumerate}
\item 어떤 체 $k'$에 대하여 $X\to\Spec(k)$가
$X\to\Spec(k')\to\Spec(k)$로 분해되면, $\omega_X^\bullet$과 $\omega_X$는
$k$를 $k'$로 바꾼 성질 (4), (5), (6)을 만족한다.
\item $K/k$가 체확장이면 $\omega_X^\bullet$과 $\omega_X$를 밑변환 $X_K$로
당긴 것들은 사상 $X_K\to\Spec(K)$에 대한 보조정리
\ref{curves-lemma-duality-dim-1}의 것들이다.
\end{enumerate}
\end{lemma}

\begin{proof}
$f:X\to\Spec(k)$를 구조사상이라 하자. 명제 (1)의 실제 의미는
$\omega_X^\bullet$과 $\omega_X$가 사상 $f':X\to\Spec(k')$에 대하여 보조정리
\ref{curves-lemma-duality-dim-1}에 나오는 것들이라는 뜻이다. 보조정리
\ref{curves-lemma-duality-dim-1}의 증명에서 $\omega_X^\bullet=a(\mathcal{O}_{\Spec(k)})$로
두었고, 여기서 $a$는 $f$에 대한 스킴의 쌍대성, 보조정리
\ref{duality-lemma-twisted-inverse-image}의 오른쪽 수반이다. 따라서
$a(\mathcal{O}_{\Spec(k)})\cong a'(\mathcal{O}_{\Spec(k)})$임을 보여야 하는데,
여기서 $a'$은 $f'$에 대한 스킴의 쌍대성, 보조정리
\ref{duality-lemma-twisted-inverse-image}의 오른쪽 수반이다.
$k'\subset H^0(X,\mathcal{O}_X)$이므로 $k'/k$는 유한확장이다(스킴의 코호몰로지, 보조정리 \ref{coherent-lemma-proper-over-affine-cohomology-finite}).
수반의 유일성에 의해 $a=a'\circ b$이고, 여기서 $b$는
$g:\Spec(k')\to\Spec(k)$에 대한 스킴의 쌍대성, 보조정리
\ref{duality-lemma-twisted-inverse-image}의 오른쪽 수반이다. 달리 말하면
$f^!=(f')^!\circ g^!$이다. 따라서 스킴의 쌍대성, 예
\ref{duality-example-affine-twisted-inverse-image}에 의해 $k'$-가군으로서
$\Hom_k(k',k)\cong k'$임을 보이면 충분하다. 이들은 같은 차원(즉 차원 $1$)의
$k'$-벡터공간이므로 성립한다.

\medskip\noindent
(2)의 증명. 스킴의 쌍대성, 보조정리
\ref{duality-lemma-more-base-change}에 의해 $a$에 대한 밑변환이 성립하기
때문이다. 스킴의 쌍대성, 주의
\ref{duality-remark-relative-dualizing-complex}의 논의도 보라.
\end{proof}

\begin{lemma}
\label{curves-lemma-closed-immersion-dim-1-CM}
$X$를 체 $k$ 위의 차원 $\leq1$인 고유 스킴이라 하고 $i:Y\to X$를
닫힌몰입이라 하자. $\omega_X^\bullet$, $\omega_X$, $\omega_Y^\bullet$,
$\omega_Y$를 보조정리 \ref{curves-lemma-duality-dim-1}에서와 같이 두자. 그러면
\begin{enumerate}
\item $\omega_Y^\bullet=R\SheafHom(\mathcal{O}_Y,\omega_X^\bullet)$이다.
\item $\omega_Y=\SheafHom(\mathcal{O}_Y,\omega_X)$이고
$i_*\omega_Y=\SheafHom_{\mathcal{O}_X}(i_*\mathcal{O}_Y,\omega_X)$이다.
\end{enumerate}
\end{lemma}

\begin{proof}
$g:Y\to\Spec(k)$와 $f:X\to\Spec(k)$를 구조사상이라 하자. 그러면
$g=f\circ i$이다. $f,g,i$에 대하여 스킴의 쌍대성, 보조정리
\ref{duality-lemma-twisted-inverse-image}의 오른쪽 수반을 각각 $a,b,c$라 하자.
수반의 유일성과 $Rg_*=Rf_*\circ Ri_*$에 의해 $b=c\circ a$이다.
보조정리 \ref{curves-lemma-duality-dim-1}의 증명에서
$\omega_X^\bullet=a(\mathcal{O}_{\Spec(k)})$,
$\omega_Y^\bullet=b(\mathcal{O}_{\Spec(k)})$로 두었다. 따라서
$\omega_Y^\bullet=c(\omega_X^\bullet)$이고, 스킴의 쌍대성, 보조정리
\ref{duality-lemma-twisted-inverse-image-closed}에 의해 (1)을 얻는다.
$\omega_X=H^{-1}(\omega_X^\bullet)$이고
$\omega_Y=H^{-1}(\omega_Y^\bullet)$이므로
$\omega_Y=\SheafHom(\mathcal{O}_Y,\omega_X)$이다. 그러면 스킴의 쌍대성,
보조정리 \ref{duality-lemma-sheaf-with-exact-support-ext}에 의해
$i_*\omega_Y=\SheafHom_{\mathcal{O}_X}(i_*\mathcal{O}_Y,\omega_X)$이다.
\end{proof}

\begin{lemma}
\label{curves-lemma-closed-subscheme-reduced-gorenstein}
$X$를 체 $k$ 위의 Gorenstein이고 환원이며 순수차원 $1$인 고유 스킴이라 하자.
$i:Y\to X$를 순수차원 $1$인 환원 닫힌 부분스킴이라 하고,
$j:Z\to X$를 $X\setminus Y$의 스킴론적 폐포라 하자. 그러면
\begin{enumerate}
\item $Y$와 $Z$는 Cohen--Macaulay이다.
\item $\mathcal{I}\subset\mathcal{O}_X$, 각각 $\mathcal{J}\subset\mathcal{O}_X$가
$X$ 안에서 $Y$, 각각 $Z$의 아이디얼층이면
$$
\mathcal{I}=i_*\mathcal{I}'\quad\text{and}\quad
\mathcal{J}=j_*\mathcal{J}'
$$
이다. 여기서 $\mathcal{I}'\subset\mathcal{O}_Z$, 각각
$\mathcal{J}'\subset\mathcal{O}_Y$는 $Z$, 각각 $Y$ 안에서 $Y\cap Z$의 아이디얼층이다.
\item $\omega_Y=\mathcal{J}'(i^*\omega_X)$이고
$i_*(\omega_Y)=\mathcal{J}\omega_X$이다.
\item $\omega_Z=\mathcal{I}'(i^*\omega_X)$이고
$i_*(\omega_Z)=\mathcal{I}\omega_X$이다.
\item 다음 짧은 완전열들이 있다.
\begin{align*}
0 \to \omega_X \to i_*i^*\omega_X \oplus j_*j^*\omega_X \to
\mathcal{O}_{Y \cap Z} \to 0 \\
0 \to i_*\omega_Y \to \omega_X \to j_*j^*\omega_X \to 0 \\
0 \to j_*\omega_Z \to \omega_X \to i_*i^*\omega_X \to 0 \\
0 \to i_*\omega_Y \oplus j_*\omega_Z \to \omega_X \to
\mathcal{O}_{Y \cap Z} \to 0 \\
0 \to \omega_Y \to i^*\omega_X \to \mathcal{O}_{Y \cap Z} \to 0 \\
0 \to \omega_Z \to j^*\omega_X \to \mathcal{O}_{Y \cap Z} \to 0
\end{align*}
\end{enumerate}
여기서 $\omega_X$, $\omega_Y$, $\omega_Z$는 보조정리
\ref{curves-lemma-duality-dim-1}에서와 같다.
\end{lemma}

\begin{proof}
환원 $1$차원 Noether 스킴은 Cohen--Macaulay이므로 (1)은 참이다. $X$가
환원이므로 스킴론적으로 $X=Y\cup Z$이다. 스킴의 사상, 보조정리
\ref{morphisms-lemma-scheme-theoretic-union}의 표기를 쓰면 그 보조정리에 의해
짧은 완전열
$$
0\to\mathcal{O}_X\to\mathcal{O}_Y\oplus\mathcal{O}_Z
\to\mathcal{O}_{Y\cap Z}\to0
$$
이 있다. 이제
$\mathcal{J}=\Ker(\mathcal{O}_X\to\mathcal{O}_Z)$,
$\mathcal{J}'=\Ker(\mathcal{O}_Y\to\mathcal{O}_{Y\cap Z})$,
$\mathcal{I}=\Ker(\mathcal{O}_X\to\mathcal{O}_Y)$,
$\mathcal{I}'=\Ker(\mathcal{O}_Z\to\mathcal{O}_{Y\cap Z})$이므로 그림쫓기를 하면
(2)를 얻는다. $\mathcal{I}+\mathcal{J}$는 $Y\cap Z$의 아이디얼층이고
$\mathcal{I}\cap\mathcal{J}=0$임을 주목하자. 따라서 다음 완전열들이 있다.
\begin{align*}
0 \to \mathcal{O}_X \to \mathcal{O}_Y \oplus \mathcal{O}_Z \to
\mathcal{O}_{Y \cap Z} \to 0 \\
0 \to \mathcal{J} \to \mathcal{O}_X \to \mathcal{O}_Z \to 0 \\
0 \to \mathcal{I} \to \mathcal{O}_X \to \mathcal{O}_Y \to 0 \\
0 \to \mathcal{J} \oplus \mathcal{I} \to \mathcal{O}_X \to
\mathcal{O}_{Y \cap Z} \to 0 \\
0 \to \mathcal{J}' \to \mathcal{O}_Y \to \mathcal{O}_{Y \cap Z} \to 0 \\
0 \to \mathcal{I}' \to \mathcal{O}_Z \to \mathcal{O}_{Y \cap Z} \to 0.
\end{align*}
$X$가 Gorenstein이므로 $\omega_X$는 가역 $\mathcal{O}_X$-가군이다(스킴의 쌍대성, 보조정리 \ref{duality-lemma-gorenstein}). $Y\cap Z$의 차원이 $0$이므로
$\omega_X|_{Y\cap Z}\cong\mathcal{O}_{Y\cap Z}$이다. 따라서 (3), (4)를
증명하면 위의 짧은 완전열들에 가역 가군 $\omega_X$를 텐서하여 명제의 짧은
완전열들을 얻는다. 대칭성에 의해 (3)만 증명하면 충분하고, (2)에 의해
$i_*(\omega_Y)=\mathcal{J}\omega_X$만 증명하면 충분하다.

\medskip\noindent
보조정리 \ref{curves-lemma-closed-immersion-dim-1-CM}에 의해
$i_*\omega_Y=\SheafHom_{\mathcal{O}_X}(i_*\mathcal{O}_Y,\omega_X)$이다.
$\omega_X$가 가역임을 다시 이용하면, $1$에서의 값매김으로 주어진 사상이
$\SheafHom_{\mathcal{O}_X}(\mathcal{O}_X/\mathcal{I},\mathcal{O}_X)$에서
$\mathcal{J}$로 가는 동형임을 보이는 것으로 충분하다. 즉 $\mathcal{J}$가
$\mathcal{I}$의 소멸아이디얼임을 보이면 되는데, 이는 위에서 따른다.
\end{proof}






\section{Riemann--Roch}
\label{curves-section-Riemann-Roch}

\noindent
$k$를 체라 하고 $X$를 $k$ 위의 차원 $\leq1$인 고유 스킴이라 하자.
다양체, 절 \ref{varieties-section-divisors-curves}에서 상수 계수의 국소
자유 $\mathcal{O}_X$-가군 $\mathcal{E}$의 차수를 공식
\begin{equation}
\label{curves-equation-degree}
\deg(\mathcal{E})=\chi(X,\mathcal{E})-
\text{rank}(\mathcal{E})\chi(X,\mathcal{O}_X)
\end{equation}
으로 정의하였다. 다양체, 정의
\ref{varieties-definition-degree-invertible-sheaf}를 보라. Chow 호몰로지 장에서는
$\mathcal{E}$의 첫째 Chern 류를 순환 위의 연산으로 정의하였고(Chow 호몰로지,
절 \ref{chow-section-intersecting-chern-classes}),
\begin{equation}
\label{curves-equation-degree-c1}
\deg(\mathcal{E})=\deg(c_1(\mathcal{E})\cap[X]_1)
\end{equation}
임을 증명하였다. Chow 호몰로지, 보조정리
\ref{chow-lemma-degree-vector-bundle}를 보라. (\ref{curves-equation-degree})와
(\ref{curves-equation-degree-c1})를 결합하면 Riemann--Roch 공식의 첫 형태
\begin{equation}
\label{curves-equation-rr}
\chi(X,\mathcal{E})=\deg(c_1(\mathcal{E})\cap[X]_1)+
\text{rank}(\mathcal{E})\chi(X,\mathcal{O}_X)
\end{equation}
를 얻는다. $\mathcal{L}$이 가역 $\mathcal{O}_X$-가군이면 다양체, 정의
\ref{varieties-definition-intersection-number}에서 정의한 수치적 교차수
$(\mathcal{L}\cdot X)$도 생각할 수 있다. 그러나 다양체, 보조정리
\ref{varieties-lemma-intersection-numbers-and-degrees-on-curves}에 의해
\begin{equation}
\label{curves-equation-numerical-degree}
(\mathcal{L}\cdot X)=\deg(\mathcal{L})
\end{equation}
이므로 새로운 내용은 없다. $\mathcal{L}$이 풍부하면 이 정수는 양수이고,
다양체, 정의 \ref{varieties-definition-degree}에서와 같이
\begin{equation}
\label{curves-equation-degree-X}
\deg_\mathcal{L}(X)=(\mathcal{L}\cdot X)=\deg(\mathcal{L})
\end{equation}
를 $\mathcal{L}$에 관한 $X$의 차수라고 한다.

\medskip\noindent
진정한 Riemann--Roch 정리를 얻으려면 $\chi(X,\mathcal{O}_X)$를 $X$ 위의
표준적인 영차원 순환의 차수로 쓰고 싶다. 완전히 일반적인 판본은 \cite{F}를
참조하라. 우리는 $X$가 Gorenstein인 경우에 쌍대성을 써서 공식을 얻겠지만,
어떤 의미에서는 이는 편법이다(예를 들어 이 방법은 더 높은 차원에서는
작동할 수 없다).

\medskip\noindent
먼저 보조정리 \ref{curves-lemma-duality-dim-1}과 \ref{curves-lemma-duality-dim-1-CM}을 써서
$\mathcal{O}_X$의 Euler 표수와 쌍대화 복합체 또는 쌍대화 가군의 Euler 표수
사이의 관계를 얻는다.

\begin{lemma}
\label{curves-lemma-euler}
$X$를 체 $k$ 위의 차원 $\leq1$인 고유 스킴이라 하자. $\omega_X^\bullet$과
$\omega_X$를 보조정리 \ref{curves-lemma-duality-dim-1}에서와 같이 두면
$$
\chi(X,\mathcal{O}_X)=\chi(X,\omega_X^\bullet)
$$
이다. $X$가 Cohen--Macaulay이고 순수차원 $1$이면
$$
\chi(X,\mathcal{O}_X)=-\chi(X,\omega_X)
$$
이다.
\end{lemma}

\begin{proof}
첫째 공식의 우변을 다음과 같이 정의한다.
$$
\chi(X,\omega_X^\bullet)=
\sum\nolimits_{i\in\mathbf{Z}}(-1)^i\dim_kH^i(X,\omega_X^\bullet).
$$
$\omega_X^\bullet$은 $D^b_{\textit{Coh}}(\mathcal{O}_X)$의 대상이므로 이는 잘 정의된다.
또한
$$
H^i(X,\omega_X^\bullet)=\Ext^i(\mathcal{O}_X,\omega_X^\bullet)
=H^{-i}(X,\mathcal{O}_X)
$$
이고 이는 항상 유한차원이며 $i=0,-1$일 때만 영이 아닐 수 있으므로도 잘
정의된다. 물론 이는 첫째 공식도 증명한다. Cohen--Macaulay인 경우 보조정리
\ref{curves-lemma-duality-dim-1-CM}에 의해 $\omega_X^\bullet=\omega_X[1]$이므로
둘째 공식은 첫째 공식에서 따른다.
\end{proof}

\noindent
우리는 $\omega_X$가 가역인 경우, 즉 $X$가 Gorenstein인 경우에 보조정리
\ref{curves-lemma-euler}를 써서 $\chi(X,\mathcal{O}_X)$에 대해 원하는 공식을 얻을 것이다.
그 명제는 $\omega_X$의 첫째 Chern 류를 $X$에 대응하는 순환 $[X]_1$과
교차한 것의 $-1/2$배가 반정수 계수를 갖고 차수가 $\chi(X,\mathcal{O}_X)$인
$X$ 위의 자연스러운 영차원 순환이라는 것이다. Riemann--Roch의 명제에서
분수가 나타나는 것은 피할 수 없다.

\begin{lemma}[Riemann--Roch]
\label{curves-lemma-rr}
$X$를 체 $k$ 위의 Gorenstein이고 순수차원 $1$인 고유 스킴이라 하자.
$\omega_X$를 보조정리 \ref{curves-lemma-duality-dim-1}에서와 같이 두면 다음이 성립한다.
\begin{enumerate}
\item $\omega_X$는 가역 $\mathcal{O}_X$-가군이다.
\item $\deg(\omega_X)=-2\chi(X,\mathcal{O}_X)$이다.
\item 상수 계수의 국소 자유 $\mathcal{O}_X$-가군 $\mathcal{E}$에 대하여
$$
\chi(X,\mathcal{E})=\deg(\mathcal{E})-
\textstyle{\frac12}\text{rank}(\mathcal{E})\deg(\omega_X)
$$
이고, 모든 $i\in\mathbf{Z}$에 대하여
$\dim_k(H^i(X,\mathcal{E}))=
\dim_k(H^{1-i}(X,\mathcal{E}^\vee\otimes_{\mathcal{O}_X}\omega_X))$이다.
\end{enumerate}
\end{lemma}

\noindent
비특이(정규) 곡선은 Gorenstein이다. 스킴의 쌍대성, 보조정리
\ref{duality-lemma-regular-gorenstein}을 보라.

\begin{proof}
Gorenstein 스킴은 Cohen--Macaulay이고(스킴의 쌍대성, 보조정리
\ref{duality-lemma-gorenstein-CM}), 따라서 보조정리
\ref{curves-lemma-duality-dim-1-CM}에 의해 $\omega_X$는 $X$ 위의 쌍대화 가군임을
상기하자. 쌍대화층이 가역이라는 것은 Gorenstein 성질의 정의에서 거의 바로
따른다. 스킴의 쌍대성, 절 \ref{duality-section-gorenstein}을 보라.
$\omega_X$에 (\ref{curves-equation-rr})를 적용하면
$$
\chi(X,\omega_X)=\deg(c_1(\omega_X)\cap[X]_1)+\chi(X,\mathcal{O}_X)
$$
이다. 이를 보조정리 \ref{curves-lemma-euler}와 결합하면
$$
2\chi(X,\mathcal{O}_X)=-\deg(c_1(\omega_X)\cap[X]_1)=-\deg(\omega_X)
$$
를 얻는데, 둘째 등호는 (\ref{curves-equation-degree-c1})에 의한다. 이를 $\mathcal{E}$에
대한 (\ref{curves-equation-rr})에 다시 대입하면 보조정리의 표시된 공식을 얻는다.
차원의 대칭성은 $X$에 대한 쌍대성에서 따른다. 주의
\ref{curves-remark-rework-duality-locally-free}를 보라.
\end{proof}





\section{몇 가지 소멸 결과}
\label{curves-section-vanishing}

\noindent
이 절에는 몇 가지 매우 약한 소멸 결과가 들어 있다. 더 많고 흥미로운 결과는
절 \ref{curves-section-more-vanishing}을 보라.

\begin{lemma}
\label{curves-lemma-automatic}
$k$를 체라 하고 $X$를 $H^0(X,\mathcal{O}_X)=k$인 $k$ 위의 차원 $1$인 고유
스킴이라 하자. 그러면 $X$는 연결이고 Cohen--Macaulay이며 순수차원 $1$이다.
\end{lemma}

\begin{proof}
$\Gamma(X,\mathcal{O}_X)=k$에는 자명하지 않은 멱등원이 없으므로 $X$는 연결이다.
이 사실만으로도 $X$가 순수차원 $1$임을 알 수 있다(차원 $0$인 기약성분이
있다면 그것은 연결성분일 것이다). 닫힌점들에 지지되는 극대 연접 부분가군을
$\mathcal{I}\subset\mathcal{O}_X$라 하자. $\mathcal{I}$은 존재하고(인자, 보조정리
\ref{divisors-lemma-remove-embedded-points}) 전역 생성된다(다양체, 보조정리
\ref{varieties-lemma-chi-tensor-finite}). $1\in\Gamma(X,\mathcal{O}_X)$은
$\mathcal{I}$의 단면이 아니므로 $\mathcal{I}=0$이다. 따라서 $X$에는 매장점이
없고(인자, 보조정리 \ref{divisors-lemma-remove-embedded-points}), 인자,
보조정리 \ref{divisors-lemma-S1-no-embedded}에 의해 $X$는 $(S_1)$을 만족한다. 그러므로
$X$는 Cohen--Macaulay이다.
\end{proof}

\noindent
이 절에서는 다음 상황에서 작업한다.

\begin{situation}
\label{curves-situation-Cohen-Macaulay-curve}
여기서 $k$는 체이고, $X$는 $H^0(X,\mathcal{O}_X)=k$인 $k$ 위의 차원 $1$인
고유 스킴이다.
\end{situation}

\noindent
보조정리 \ref{curves-lemma-automatic}에 의해 $X$는 Cohen--Macaulay이고 순수차원 $1$이다.
보조정리 \ref{curves-lemma-duality-dim-1}과 \ref{curves-lemma-duality-dim-1-CM}에서 논한
쌍대화 가군 $\omega_X$의 $H^1$은 소멸하지 않는다. 실제로
$\dim_kH^1(X,\omega_X)=\dim_kH^0(X,\mathcal{O}_X)=1$이다. $\omega_X$보다
조금이라도 더 ``양''인 것은 $H^1$이 소멸함을 알 수 있다.

\begin{lemma}
\label{curves-lemma-vanishing}
상황 \ref{curves-situation-Cohen-Macaulay-curve}에서 생각하자.
$H^1(X,\mathcal{Q})=0$인(예를 들어 $\dim(\text{Supp}(\mathcal{Q}))=0$인)
연접 $\mathcal{O}_X$-가군의 완전열
$$
\omega_X\to\mathcal{F}\to\mathcal{Q}\to0
$$
이 주어지면, $H^1(X,\mathcal{F})=0$이거나
$\mathcal{F}=\omega_X\oplus\mathcal{Q}$이다.
\end{lemma}

\begin{proof}
(괄호 안의 명제는 스킴의 코호몰로지, 보조정리
\ref{coherent-lemma-coherent-support-dimension-0}에서 따른다.)
$H^0(X,\mathcal{O}_X)=k$는 $H^1(X,\omega_X)$의 쌍대이므로(절
\ref{curves-section-Riemann-Roch} 참조) $\dim H^1(X,\omega_X)=1$이다. 층 $\omega_X$는
연접 $\mathcal{O}_X$-가군들의 범주 위에서 함자
$\mathcal{F}\mapsto\Hom_k(H^1(X,\mathcal{F}),k)$를 표현한다(스킴의 쌍대성,
보조정리 \ref{duality-lemma-dualizing-module-proper-over-A}). 명제와 같은 완전열을
생각하고 $H^1(X,\mathcal{F})\not =0$이라고 하자. $H^1(X,\mathcal{Q})=0$이므로
$H^1(X,\omega_X)\to H^1(X,\mathcal{F})$는 동형이다. 위에 말한 $\omega_X$의
보편성에 의해, $H^1$에 대한 작용이 이 동형의 역인 사상
$\mathcal{F}\to\omega_X$가 존재한다. 합성
$\omega_X\to\mathcal{F}\to\omega_X$는 (보편성에 의해) 항등이므로 증명된다.
\end{proof}

\begin{lemma}
\label{curves-lemma-vanishing-twist}
상황 \ref{curves-situation-Cohen-Macaulay-curve}에서 생각하자. $\mathcal{L}$을
전역 생성되고 $\mathcal{O}_X$와 동형이 아닌 가역 $\mathcal{O}_X$-가군이라 하자.
그러면 $H^1(X,\omega_X\otimes\mathcal{L})=0$이다.
\end{lemma}

\begin{proof}
절 \ref{curves-section-Riemann-Roch}에서 논한 쌍대성에 의해
$H^0(X,\mathcal{L}^{\otimes-1})=0$임을 보이면 된다. 그렇지 않으면
$\mathcal{L}^{\otimes-1}$의 전역단면 $t$와 $st\not =0$인 $\mathcal{L}$의
전역단면 $s$를 고를 수 있다. 그러나 $H^0(X,\mathcal{O}_X)=k$라는 가정에 의해
$st$는 $1$의 상수배이다. 그러면 $\mathcal{L}\cong\mathcal{O}_X$이어서 모순이다.
\end{proof}

\begin{lemma}
\label{curves-lemma-globally-generated}
상황 \ref{curves-situation-Cohen-Macaulay-curve}에서 생각하자.
$\dim(\text{Supp}(\mathcal{Q}))=0$, $\dim_kH^0(X,\mathcal{Q})\geq2$인 연접
$\mathcal{O}_X$-가군들의 완전열
$$
\omega_X\to\mathcal{F}\to\mathcal{Q}\to0
$$
이 주어졌다고 하자. 또한 $\mathcal{Q}'\to\mathcal{Q}$가 단사가 되는 영이 아닌
부분가군 $\mathcal{Q}'\subset\mathcal{F}$가 없다고 하자. 그러면 전역단면들로
생성되는 $\mathcal{F}$의 부분가군은 $\mathcal{Q}$ 위로 전사한다.
\end{lemma}

\begin{proof}
$\mathcal{F}'\subset\mathcal{F}$를 전역단면들과 $\omega_X\to\mathcal{F}$의
상으로 생성되는 부분가군이라 하자. $\dim_kH^0(X,\mathcal{Q})\geq2$이고
$\dim_kH^1(X,\omega_X)=\dim_kH^0(X,\mathcal{O}_X)=1$이므로
$\mathcal{F}'\to\mathcal{Q}$는 영이 아니며
$\omega_X\to\mathcal{F}'$은 동형이 아니다. 따라서 보조정리
\ref{curves-lemma-vanishing}과 $\mathcal{F}$에 대한 가정에 의해
$H^1(X,\mathcal{F}')=0$이다. 짧은 완전열
$$
0\to\mathcal{F}'\to\mathcal{F}\to
\mathcal{Q}/\Im(\mathcal{F}'\to\mathcal{Q})\to0
$$
을 생각하자. 오른쪽 몫이 영이 아니면 $H^0(X,\mathcal{F})$가
$H^0(X,\mathcal{F}')$보다 커져 모순이다.
\end{proof}

\noindent
다음은 전역 생성에 관한 한 예이다.

\begin{lemma}
\label{curves-lemma-globally-generated-curve}
상황 \ref{curves-situation-Cohen-Macaulay-curve}에서 $X$가 정역이라고 가정하자.
$0\to\omega_X\to\mathcal{F}\to\mathcal{Q}\to0$을 연접 $\mathcal{O}_X$-가군들의
짧은 완전열이라 하자. $\mathcal{F}$는 꼬임이 없고,
$\dim(\text{Supp}(\mathcal{Q}))=0$, $\dim_kH^0(X,\mathcal{Q})\geq2$라 하자.
그러면 $\mathcal{F}$는 전역 생성된다.
\end{lemma}

\begin{proof}
전역단면들로 생성되는 부분가군을 $\mathcal{F}'$이라 하자. 보조정리
\ref{curves-lemma-globally-generated}에 의해 $\mathcal{F}'\to\mathcal{Q}$는 전사이고,
특히 $\mathcal{F}'\not =0$이다. $X$가 곡선이므로
$\mathcal{F}'\subset\mathcal{F}$는 계수 $1$인 층들의 포함이고, 따라서
$\mathcal{Q}'=\mathcal{F}/\mathcal{F}'$은 유한개의 점에 지지된다.
모순을 얻기 위해 $\mathcal{Q}'$이 영이 아니라고 하자. 그러면
$H^1(X,\mathcal{F}')\not =0$이다. 보편성에 의해 영이 아닌 사상
$\mathcal{F}'\to\omega_X$를 얻는다(스킴의 쌍대성, 보조정리
\ref{duality-lemma-dualizing-module-proper-over-A}). 합성
$\mathcal{F}'\to\omega_X\to\mathcal{F}$의 상은 전역단면들로 생성되므로
$\mathcal{F}'$ 안에 있다. 따라서 영이 아닌 자기 사상
$\mathcal{F}'\to\mathcal{F}'$을 얻는다. $\mathcal{F}'$은 고유 곡선 위의 계수
$1$인 꼬임 없는 층이므로 이것은 자기동형이어야 한다(세부사항은 생략한다).
그러면 $\mathcal{F}'$은 $\omega_X\subset\mathcal{F}$에 포함되는데, 이는
$\mathcal{F}'\to\mathcal{Q}$의 전사성과 모순이다.
\end{proof}

\begin{lemma}
\label{curves-lemma-tensor-omega-with-globally-generated-invertible}
상황 \ref{curves-situation-Cohen-Macaulay-curve}에서 생각하자. $\mathcal{L}$을
$\deg(\mathcal{L})\geq2$인 매우 풍부한 가역 $\mathcal{O}_X$-가군이라 하자.
그러면 $\omega_X\otimes_{\mathcal{O}_X}\mathcal{L}$은 전역 생성된다.
\end{lemma}

\begin{proof}
$k$가 대수적으로 닫혔다고 가정하자. 닫힌점 $x\in X$를 잡자.
$C_i\subset X$를 기약성분들이라 하고, 각 $i$에 대하여 $x_i\in C_i$를
일반점이라 하자. 다양체, 보조정리
\ref{varieties-lemma-very-ample-vanish-at-point}에 의해, $x$에서는 영이 되지만
모든 $i$에 대해 $x_i$에서는 영이 되지 않는 단면 $s\in H^0(X,\mathcal{L})$을
고를 수 있다. 이에 대응하는 $s$의 가군 사상
$s:\mathcal{O}_X\to\mathcal{L}$은 단사이고, 그 여핵
$\mathcal{Q}$는 유한개의 점에 지지되며 $H^0(X,\mathcal{Q})\geq2$이다.
이에 대응하는 완전열
$$
0\to\omega_X\to\omega_X\otimes\mathcal{L}
\to\omega_X\otimes\mathcal{Q}\to0
$$
을 생각하자. 보조정리 \ref{curves-lemma-globally-generated}에 의해 전역단면들로 생성되는
가군은 $\omega_X\otimes\mathcal{Q}$ 위로 전사한다. $x$가 임의였으므로 보조정리가
증명된다. 일부 세부사항은 생략한다.

\medskip\noindent
$k$가 대수적으로 닫히지 않은 경우를 대수적으로 닫힌 경우로 환원하겠다.
독자는 증명의 나머지를 건너뛰어도 좋다. $k$의 대수적 폐포 $\overline{k}$를
택하고 밑변환 $X_{\overline{k}}$를 생각하자. $X_{\overline{k}}\to\Spec(\overline{k})$가
상황 \ref{curves-situation-Cohen-Macaulay-curve}의 예임을 확인하자. 평탄 밑변환에
의해(스킴의 코호몰로지, 보조정리
\ref{coherent-lemma-flat-base-change-cohomology})
$H^0(X_{\overline{k}},\mathcal{O})=\overline{k}$이다. 스킴 $X_{\overline{k}}$는
$\overline{k}$ 위에서 고유이고(스킴의 사상, 보조정리
\ref{morphisms-lemma-base-change-proper}) 순수차원 $1$이다(스킴의 사상, 보조정리
\ref{morphisms-lemma-dimension-fibre-after-base-change}). $\omega_X$를
$X_{\overline{k}}$로 당긴 것은 보조정리 \ref{curves-lemma-sanity-check-duality}에 의해
$X_{\overline{k}}$의 쌍대화 가군이다. $\mathcal{L}$을 $X_{\overline{k}}$로 당긴
것은 매우 풍부하다(스킴의 사상, 보조정리
\ref{morphisms-lemma-very-ample-base-change}). $\mathcal{L}$을
$X_{\overline{k}}$로 당긴 것의 차수는 $X$ 위에서의
$\mathcal{L}$의 차수와 같다(다양체, 보조정리
\ref{varieties-lemma-degree-base-change}). 마지막으로
$\omega_X\otimes\mathcal{L}$이 전역 생성일 필요충분조건은 그 밑변환이
전역 생성인 것이다(다양체, 보조정리
\ref{varieties-lemma-globally-generated-base-change}). 이로써 결과는 밑체가
대수적으로 닫힌 경우의 결과에서 따른다.
\end{proof}




\section{매우 풍부한 가역층}
\label{curves-section-very-ample}

\noindent
대수적으로 닫힌 체 위의 유한형 스킴 $X$에서 가역 가군 $\mathcal{L}$이 매우
풍부일 자주 쓰이는 판정법은 $\mathcal{L}$의 단면들이 점과 접벡터를 분리한다는
것이다(다양체, 절
\ref{varieties-section-separating-points-tangent-vectors}). 다음은 곡선에 대한
또 다른 판정법이다. 다양체, 소절
\ref{varieties-subsection-regularity}와 비교하라.

\begin{lemma}
\label{curves-lemma-criterion-very-ample}
$k$를 체, $X$를 $H^0(X,\mathcal{O}_X)=k$인 $k$ 위의 차원 $1$인 고유 스킴,
$\mathcal{L}$을 가역 $\mathcal{O}_X$-가군이라 하자. 다음을 가정하자.
\begin{enumerate}
\item $\mathcal{L}$은 정칙 전역단면을 가진다.
\item $H^1(X,\mathcal{L})=0$이다.
\item $\mathcal{L}$은 풍부하다.
\end{enumerate}
그러면 $\mathcal{L}^{\otimes6}$은 $k$ 위의 $X$에서 매우 풍부하다.
\end{lemma}

\begin{proof}
$s$를 $\mathcal{L}$의 정칙 전역단면이라 하고, $i:Z=Z(s)\to X$를 $s$의
영점 스킴이라 하자. 인자, 절
\ref{divisors-section-effective-Cartier-invertible}을 보라. 조건 (3)에 의해
$Z\not =\emptyset$이다(작은 세부사항은 생략한다). 짧은 완전열
$$
0\to\mathcal{O}_X\xrightarrow{s}\mathcal{L}\to i_*(\mathcal{L}|_Z)\to0
$$
을 생각하자. $\mathcal{L}$을 텐서하면
$$
0\to\mathcal{L}\to\mathcal{L}^{\otimes2}
\to i_*(\mathcal{L}^{\otimes2}|_Z)\to0
$$
을 얻는다. $Z$의 차원은 $0$이고(인자, 보조정리
\ref{divisors-lemma-effective-Cartier-makes-dimension-drop}), 따라서 Artin 환의
스펙트럼이다(다양체, 보조정리
\ref{varieties-lemma-algebraic-scheme-dim-0}). 그러므로
$\mathcal{L}|_Z\cong\mathcal{O}_Z$이다(대수학, 보조정리
\ref{algebra-lemma-locally-free-semi-local-free}). 위 짧은 완전열은 또한
$H^1(X,\mathcal{L}^{\otimes2})=0$임을 보여 준다(예를 들어 오른쪽 자리에서의
소멸을 보려면 다양체, 보조정리
\ref{varieties-lemma-chi-tensor-finite}를 쓰라). $n\geq1$에 대한 귀납법과 완전열
$$
0\to\mathcal{L}^{\otimes n}\xrightarrow{s}\mathcal{L}^{\otimes n+1}
\to i_*(\mathcal{L}^{\otimes n+1}|_Z)\to0
$$
을 쓰면, $n>0$에 대하여 $H^1(X,\mathcal{L}^{\otimes n})=0$이고
$\mathcal{L}^{\otimes n+1}|_Z\cong\mathcal{O}_Z$의 자명화를 주는
$\mathcal{L}^{\otimes n+1}$의 전역단면 $t_{n+1}$이 존재함을 알 수 있다.

\medskip\noindent
곱셈사상
$$
\mu_n:
H^0(X,\mathcal{L})\otimes_kH^0(X,\mathcal{L}^{\otimes n})
\oplus
H^0(X,\mathcal{L}^{\otimes2})\otimes_kH^0(X,\mathcal{L}^{\otimes n-1})
\longrightarrow H^0(X,\mathcal{L}^{\otimes n+1})
$$
을 생각하자. $n\geq3$이면 이것이 전사라고 주장한다. 이를 보이기 위해 짧은
완전열
$$
0\to\mathcal{L}^{\otimes n}\xrightarrow{s}\mathcal{L}^{\otimes n+1}
\to i_*(\mathcal{L}^{\otimes n+1}|_Z)\to0
$$
을 생각한다. 이 열의 왼쪽에서 오는 $\mathcal{L}^{\otimes n+1}$의 단면들은
$\mu_n$의 상 안에 있다. 한편
$H^0(\mathcal{L}^{\otimes2})\to H^0(\mathcal{L}^{\otimes2}|_Z)$가 전사이고(위 참조)
$t_{n-1}$은 $\mathcal{L}^{\otimes n-1}|_Z$의 한 자명화로 가므로,
$\mu_n(H^0(X,\mathcal{L}^{\otimes2})\otimes t_{n-1})$은
$H^0(X,\mathcal{L}^{\otimes n+1})$의 부분공간으로서
$\mathcal{L}^{\otimes n+1}|_Z$의 전역단면 위로 전사한다. 이로써 주장이 증명된다.

\medskip\noindent
앞 문단의 주장으로부터 등급 $k$-대수
$$
S=\bigoplus\nolimits_{n\geq0}H^0(X,\mathcal{L}^{\otimes n})
$$
가 $k$ 위에서 차수 $0,1,2,3$의 원소들로 생성됨을 얻는다.
$X=\text{Proj}(S)$임을 상기하자. 스킴의 사상, 보조정리
\ref{morphisms-lemma-proper-ample-is-proj}를 보라. 따라서
$S^{(6)}=\bigoplus_nS_{6n}$은 차수 $1$에서 생성된다. 이는 원하는 대로
$\mathcal{L}^{\otimes6}$이 매우 풍부하다는 뜻이다.
\end{proof}

\begin{lemma}
\label{curves-lemma-criterion-very-ample-bis}
$k$를 체, $X$를 $H^0(X,\mathcal{O}_X)=k$인 $k$ 위의 차원 $1$인 고유 스킴,
$\mathcal{L}$을 가역 $\mathcal{O}_X$-가군이라 하자. 다음을 가정하자.
\begin{enumerate}
\item $\mathcal{L}$은 전역 생성된다.
\item $H^1(X,\mathcal{L})=0$이다.
\item $\mathcal{L}$은 풍부하다.
\end{enumerate}
그러면 $\mathcal{L}^{\otimes2}$는 $k$ 위의 $X$에서 매우 풍부하다.
\end{lemma}

\begin{proof}
$s_0,\ldots,s_n$을 $k$ 위의 $H^0(X,\mathcal{L}^{\otimes2})$의 기저로 택하자.
성질 (1)에 의해 $\mathcal{L}^{\otimes2}$는 전역 생성되고 사상
$$
\varphi_{\mathcal{L}^{\otimes2},(s_0,\ldots,s_n)}:
X\longrightarrow\mathbf{P}^n_k
$$
를 얻는다. 구성, 절
\ref{constructions-section-projective-space}을 보라. 보조정리의 주장은 이 사상이
닫힌몰입이라는 것이다. 이를 확인할 때 $k$를 그 대수적 폐포로 바꾸어도 된다.
내림, 보조정리 \ref{descent-lemma-descending-property-closed-immersion}를 보라.
따라서 $k$가 대수적으로 닫혔다고 가정해도 된다.

\medskip\noindent
$k$가 대수적으로 닫혔다고 하자. 각 일반점 $\eta_i\in X$에 대하여
$V_i\subset H^0(X,\mathcal{L})$를 $\eta_i$에서 영이 되는 단면들의
$k$-부분벡터공간이라 하자. $\mathcal{L}$이 전역 생성되므로
$V_i\not = H^0(X,\mathcal{L})$이다. $X$에는 유한개의 기약성분만 있고 $k$는
무한체이므로, 모든 $i$에 대해 $\eta_i$에서 영이 되지 않는
$s\in H^0(X,\mathcal{L})$를 찾을 수 있다. 그러면 $s$는 $\mathcal{L}$의
정칙 단면이다($X$는 보조정리 \ref{curves-lemma-automatic}에 의해 Cohen--Macaulay이므로
$\mathcal{L}$에는 매장된 연관점이 없기 때문이다).

\medskip\noindent
특히 보조정리 \ref{curves-lemma-criterion-very-ample}의 증명에 나온 명제들이 모두 이
$s$에 대하여 성립한다. 또한 $\mathcal{L}$이 전역 생성되므로
$t|_Z$가 영이 아닌 전역단면 $t\in H^0(X,\mathcal{L})$를 찾을 수 있다
($Z$의 점들이 유한함을 써서 위와 같이 논증하라). 그러면 보조정리
\ref{curves-lemma-criterion-very-ample}의 증명에서 $t$를 써서 곱셈사상
$$
\mu_n:H^0(X,\mathcal{L})\otimes_kH^0(X,\mathcal{L}^{\otimes2})
\longrightarrow H^0(X,\mathcal{L}^{\otimes3})
$$
도 전사임을 알 수 있다. 따라서
$$
S=\bigoplus\nolimits_{n\geq0}H^0(X,\mathcal{L}^{\otimes n})
$$
는 $k$ 위에서 차수 $0,1,2$의 원소들로 생성된다. 보조정리
\ref{curves-lemma-criterion-very-ample}의 증명에서와 같이 논증하면
$S^{(2)}=\bigoplus_nS_{2n}$은 차수 $1$에서 생성된다. 이는 원하는 대로
$\mathcal{L}^{\otimes2}$가 매우 풍부하다는 뜻이다. 일부 세부사항은 생략한다.
\end{proof}







\section{곡선의 종수}
\label{curves-section-genus}

\noindent
$X$가 체 $k$ 위의 매끄럽고 사영이며 기하적으로 기약인 곡선이면, 앞서
$X$의 종수를 $H^1(X,\mathcal{O}_X)$의 차원으로 정의하였다. 곡선의 Picard 스킴, 정의 \ref{pic-definition-genus}를 보라. 이 경우
$H^0(X,\mathcal{O}_X)=k$임을 주목하자. 다양체, 보조정리
\ref{varieties-lemma-regular-functions-proper-variety}를 보라. 이를 다음과 같이
일반화하자.

\begin{definition}
\label{curves-definition-genus}
$k$를 체라 하고 $X$를 $H^0(X,\mathcal{O}_X)=k$인 $k$ 위의 차원 $1$인 고유
스킴이라 하자. $X$의 {\it 종수}는 $g=\dim_kH^1(X,\mathcal{O}_X)$이다.
\end{definition}

\noindent
이를 때때로 $X$의 {\it 산술종수}라고 한다. 문헌에서는 $k$ 위의 고유 곡선
$X$의 산술종수를 때때로
$$
p_a(X)=1-\chi(X,\mathcal{O}_X)=
1-\dim_kH^0(X,\mathcal{O}_X)+\dim_kH^1(X,\mathcal{O}_X)
$$
으로 정의한다. 우리는 $H^0(X,\mathcal{O}_X)=k$를 가정하므로, 우리의 정의가
적용될 때 두 정의는 일치한다. 그러나 다음을 주의하라.
\begin{enumerate}
\item $p_a(X)$는 음수가 될 수 있다.
\item $p_a(X)$는 밑체 $k$에 의존하므로 $p_a(X/k)$라고 써야 한다.
\end{enumerate}
예를 들어 $k=\mathbf Q$이고 $X=\mathbf{P}^1_{\mathbf Q(i)}$이면
$p_a(X/\mathbf Q)=-1$이고 $p_a(X/\mathbf Q(i))=0$이다.

\medskip\noindent
우리 정의의 가정 $H^0(X,\mathcal{O}_X)=k$에는 두 가지 귀결이 있다. 한편으로
밑체에 관한 혼동이 없다는 뜻이다. 다른 한편으로 스킴 $X$가 Cohen--Macaulay이고
순수차원 $1$임을 함의한다(보조정리 \ref{curves-lemma-automatic}). $\omega_X$가 보조정리
\ref{curves-lemma-duality-dim-1}과 \ref{curves-lemma-duality-dim-1-CM}에서와 같은 쌍대화
가군이면, 쌍대성에 의해
\begin{equation}
\label{curves-equation-genus}
g=\dim_kH^1(X,\mathcal{O}_X)=\dim_kH^0(X,\omega_X)
\end{equation}
이다. 주의 \ref{curves-remark-rework-duality-locally-free}를 보라.

\medskip\noindent
$X$가 $k$ 위의 차원 $\leq1$인 고유 스킴이고 $H^0(X,\mathcal{O}_X)$가 밑체
$k$와 같지 않으면, 위 표시식으로 주어진 산술종수 $p_a(X)$ 대신 불변량
$\chi(X,\mathcal{O}_X)$를 사용하겠다. 실제로 \cite[276쪽]{FAC} 및
\cite[Introduction]{Hirzebruch}에서는 $\chi(X,\mathcal{O}_X)$를 산술종수라고
부를 것을 주장한다.

\begin{lemma}
\label{curves-lemma-genus-base-change}
$k'/k$를 체확장이라 하자. $X$를 $H^0(X,\mathcal{O}_X)=k$인 $k$ 위의 차원
$1$인 고유 스킴이라 하자. 그러면 $X_{k'}$는
$H^0(X_{k'},\mathcal{O}_{X_{k'}})=k'$인 $k'$ 위의 차원 $1$인 고유 스킴이다.
또한 $X_{k'}$의 종수는 $X$의 종수와 같다.
\end{lemma}

\begin{proof}
$X_{k'}$의 차원이 $1$이라는 것은 예를 들어 스킴의 사상, 보조정리
\ref{morphisms-lemma-dimension-fibre-after-base-change}에서 따른다.
$X_{k'}\to\Spec(k')$는 스킴의 사상, 보조정리
\ref{morphisms-lemma-base-change-proper}에 의해 고유하다. 등식
$H^0(X_{k'},\mathcal{O}_{X_{k'}})=k'$은 스킴의 코호몰로지, 보조정리
\ref{coherent-lemma-flat-base-change-cohomology}에서 따른다. 종수가 같다는 것도
같은 보조정리에서 따른다.
\end{proof}

\begin{lemma}
\label{curves-lemma-genus-gorenstein}
$k$를 체라 하고 $X$를 $H^0(X,\mathcal{O}_X)=k$인 $k$ 위의 차원 $1$인 고유
스킴이라 하자. $X$가 Gorenstein이면
$$
\deg(\omega_X)=2g-2
$$
이다. 여기서 $g$는 $X$의 종수이고 $\omega_X$는 보조정리
\ref{curves-lemma-duality-dim-1}에서와 같다.
\end{lemma}

\begin{proof}
보조정리 \ref{curves-lemma-rr}에서 바로 따른다.
\end{proof}

\begin{lemma}
\label{curves-lemma-genus-smooth}
$X$를 $H^0(X,\mathcal{O}_X)=k$인 체 $k$ 위의 매끄러운 고유 곡선이라 하자. 그러면
$$
\dim_kH^0(X,\Omega_{X/k})=g
\quad\text{and}\quad
\deg(\Omega_{X/k})=2g-2
$$
이다. 여기서 $g$는 $X$의 종수이다.
\end{lemma}

\begin{proof}
보조정리 \ref{curves-lemma-duality-dim-1}에 의해 $\Omega_{X/k}=\omega_X$이다. 따라서
공식들은 (\ref{curves-equation-genus})과 보조정리 \ref{curves-lemma-genus-gorenstein}에서 따른다.
\end{proof}





\section{평면곡선}
\label{curves-section-plane-curves}

\noindent
$k$를 체라 하자. {\it 평면곡선}이란 $\mathbf{P}^2_k$의 닫힌 부분스킴과 동형인
곡선 $X$를 말한다. 흔히 매장 $X\to\mathbf{P}^2_k$도 주어진 것으로 생각한다.
인자, 예 \ref{divisors-example-closed-subscheme-of-proj}에 의해 곡선은
이에 대응하는 동차 아이디얼
% P07-ERR-0297
$$
I(X)=\Ker\left(k[T_0,T_2,T_2]\longrightarrow
\bigoplus\Gamma(X,\mathcal{O}_X(n))\right)
$$
에 의해 결정된다. 이 상황에서 $\mathbf{P}^2_k$의 닫힌 부분스킴으로서
% P07-ERR-0297
$$
X=\text{Proj}(k[T_0,T_2,T_2]/I)
$$
임을 상기하자. 이 구성들에 대한 더 일반적인 내용은 인자, 예
\ref{divisors-example-closed-subscheme-of-proj} 및 거기에 인용된 참고문헌을
보라. 어떤 동차다항식 $F\in k[T_0,T_1,T_2]$에 대하여 $I(X)=(F)$임을 알 수
있다. 보조정리 \ref{curves-lemma-equation-plane-curve}를 보라. $X$가 기약이므로 $F$도
기약이다. 보조정리 \ref{curves-lemma-plane-curve}를 보라. 또한 $d=\deg(F)$로 두고
짧은 완전열
$$
0\to\mathcal{O}_{\mathbf{P}^2_k}(-d)\xrightarrow{F}
\mathcal{O}_{\mathbf{P}^2_k}\to\mathcal{O}_X\to0
$$
을 살펴보면 $H^0(X,\mathcal{O}_X)=k$이고 $X$의 종수가
$(d-1)(d-2)/2$임을 얻는다. 보조정리 \ref{curves-lemma-genus-plane-curve}의 증명을 보라.

\medskip\noindent
매끄러운 평면곡선을 찾는 가장 쉬운 방법은 방정식을 명시적으로 쓰는 것이다.
$p$로 $k$의 표수를 나타내자. $p$가 $d$를 나누지 않으면
$$
F=T_0^d+T_1^d+T_2^d
$$
로 둘 수 있다. 이에 대응하는 곡선 $X=V_+(F)$를 차수 $d$의 {\it Fermat 곡선}
이라 한다. 각 표준 아핀 조각 $D_+(T_i)$에서 아핀곡선
$$
\Spec(k[x,y]/(x^d+y^d+1))
$$
과 동형인 곡선을 얻으므로 이는 매끄럽다. $d x^{d-1}$과 $d y^{d-1}$이
$k[x,y]/(x^d+y^d+1)$에서 단위 아이디얼을 생성하므로, 대수학, 보조정리
\ref{algebra-lemma-relative-global-complete-intersection-smooth}에 의해 환사상
$k\to k[x,y]/(x^d+y^d+1)$은 매끄럽다. $p|d$이지만 $p\not =3$이면 방정식
$$
F=T_0^{d-1}T_1+T_1^{d-1}T_2+T_2^{d-1}T_0
$$
을 쓸 수 있다. 실제로 아핀 조각에서는 $x+x^{d-1}y+y^{d-1}$과 그 편도함수
$1-x^{d-2}y$, $x^{d-1}-y^{d-2}$를 얻으며, 이 셋의 공통 영점집합은
공집합이다.\footnote{실제로 $x^{d-1}=y^{d-2}$이면
$0=x+x^{d-1}y+y^{d-1}=x+2x^{d-1}y$이다. $1=x^{d-2}y$이므로 $x\not =0$이고,
따라서 $0=1+2x^{d-2}y=3$을 얻는데 $3=0$이 아니면 모순이다.}
표수 $3$에서의 예는 독자에게 맡긴다.

\medskip\noindent
더 일반적으로 임의의 체 $k$와 임의의 $n,d$에 대하여 $\mathbf{P}^n_k$ 안에는
차수 $d$인 매끄러운 초곡면이 존재한다. 예를 들어 \cite{Poonen}을 보라.

\medskip\noindent
물론 이런 방법으로는 종수가 삼각수인 매끄러운 곡선만 얻는다. 임의의 종수인
매끄러운 곡선을 얻으려면 $\mathbf{P}^1\times\mathbf{P}^1$ 위에 놓인 매끄러운
곡선을 찾을 수 있다. % P07-ERR-0298
(향후 참고문헌을 여기에 넣을 것).

\begin{lemma}
\label{curves-lemma-equation-plane-curve}
$Z\subset\mathbf{P}^2_k$를 순수차원 $1$이고 매장점이 없는 닫힌 부분스킴이라
하자(동치로 $Z$는 Cohen--Macaulay이다). 그러면 $Z$에 대응하는 아이디얼
$I(Z)\subset k[T_0,T_1,T_2]$는 주 아이디얼이다.
\end{lemma}

\begin{proof}
이는 인자, 보조정리
\ref{divisors-lemma-equation-codim-1-in-projective-space}의 특수한 경우이다
(다양체, 보조정리
\ref{varieties-lemma-equation-codim-1-in-projective-space}도 보라).
괄호 안의 명제는 $1$차원 Noether 스킴이 Cohen--Macaulay일 필요충분조건이
매장점이 없는 것이라는 사실에서 따른다. 인자, 보조정리
\ref{divisors-lemma-noetherian-dim-1-CM-no-embedded-points}를 보라.
\end{proof}

\begin{lemma}
\label{curves-lemma-plane-curve}
$Z\subset\mathbf{P}^2_k$를 보조정리 \ref{curves-lemma-equation-plane-curve}에서와 같이 두고,
어떤 $F\in k[T_0,T_1,T_2]$에 대하여 $I(Z)=(F)$라 하자. 그러면 $Z$가 곡선일
필요충분조건은 $F$가 기약인 것이다.
\end{lemma}

\begin{proof}
$F$가 가약이고 $F=F'F''$이라 하자. $F'$으로 정의되는 $\mathbf{P}^2_k$의
닫힌 부분스킴을 $Z'$이라 하자. 분명히 $Z'\subset Z$이고 $Z'\not = Z$이다.
$Z'$도 차원 $1$이므로 $Z$가 환원이 아니거나 $Z$가 기약이 아니다. 반대로
$Z=\sum a_iD_i$라 쓰자. 여기서 $D_i$들은 $Z$의 기약성분들이다. 인자,
보조정리 \ref{divisors-lemma-codim-1-part}와
\ref{divisors-lemma-codimension-1-is-effective-Cartier}를 보라.
$F_i\in k[T_0,T_1,T_2]$를 $D_i$의 아이디얼을 생성하는 동차다항식이라 하자.
그러면 $F$와 $\prod F_i^{a_i}$가 $\mathbf{P}^2_k$의 같은 닫힌 부분스킴을
잘라냄이 분명하다. 따라서 둘 다 $Z$의 아이디얼을 생성하므로 어떤
$\lambda\in k^*$에 대하여 $F=\lambda\prod F_i^{a_i}$이다. 그러므로 $F$가
기약이면 $Z$는 소인자, 즉 곡선이다.
\end{proof}

\begin{lemma}
\label{curves-lemma-genus-plane-curve}
$Z\subset\mathbf{P}^2_k$를 보조정리 \ref{curves-lemma-equation-plane-curve}에서와 같이 두고,
어떤 $F\in k[T_0,T_1,T_2]$에 대하여 $I(Z)=(F)$라 하자. 그러면
$H^0(Z,\mathcal{O}_Z)=k$이고, $d=\deg(F)$라 할 때 $Z$의 종수는
$(d-1)(d-2)/2$이다.
\end{lemma}

\begin{proof}
$S=k[T_0,T_1,T_2]$라 하자. 등급가군의 완전열
$$
0\to S(-d)\xrightarrow{F}S\to S/(F)\to0
$$
이 있다. 주어진 닫힌몰입을 $i:Z\to\mathbf{P}^2_k$라 하자. % P07-ERR-0299
완전함자 $\widetilde{\ }$를 적용하면(구성, 보조정리
\ref{constructions-lemma-proj-sheaves})
$$
0\to\mathcal{O}_{\mathbf{P}^2_k}(-d)\to\mathcal{O}_{\mathbf{P}^2_k}
\to i_*\mathcal{O}_Z\to0
$$
을 얻는다. $F$가 $Z$의 아이디얼을 생성하기 때문이다.
$\mathcal{O}_{\mathbf{P}^2_k}(-d)$와 $\mathcal{O}_{\mathbf{P}^2_k}$의 코호몰로지군은
스킴의 코호몰로지, 보조정리
\ref{coherent-lemma-cohomology-projective-space-over-ring}에 주어져 있다.
한편 스킴의 코호몰로지, 보조정리
\ref{coherent-lemma-relative-affine-cohomology}에 의해
$H^q(Z,\mathcal{O}_Z)=H^q(\mathbf{P}^2_k,i_*\mathcal{O}_Z)$이다. 긴 완전
코호몰로지 열을 적용하면 먼저
$$
k=H^0(\mathbf{P}^2_k,\mathcal{O}_{\mathbf{P}^2_k})
\longrightarrow H^0(Z,\mathcal{O}_Z)
$$
가 동형임을 얻고, 이어서 경계사상
$$
H^1(Z,\mathcal{O}_Z)\longrightarrow
H^2(\mathbf{P}^2_k,\mathcal{O}_{\mathbf{P}^2_k}(-d))
\cong k[T_0,T_1,T_2]_{d-3}
$$
이 동형임을 얻는다. 이 공간의 차원이 $(d-1)(d-2)/2$임은 쉽게 알 수
있으므로 증명이 끝난다.
\end{proof}

\begin{lemma}
\label{curves-lemma-smooth-plane-curve-point-over-separable}
$Z\subset\mathbf{P}^2_k$를 보조정리 \ref{curves-lemma-equation-plane-curve}에서와 같이 두고,
어떤 $F\in k[T_0,T_1,T_2]$에 대하여 $I(Z)=(F)$라 하자.
$Z\to\Spec(k)$가 적어도 한 점에서 매끄럽고 $k$가 무한체이면, 매끄러운
자리에 속하고 $\kappa(z)/k$가 차수 $d$ 이하인 유한 분리가 되는 닫힌점
$z\in Z$가 존재한다.
\end{lemma}

\begin{proof}
$z'\in Z$가 $Z\to\Spec(k)$가 매끄러운 점이라고 하자. 필요하면 좌표를
다시 번호 매겨 $z'$이 $D_+(T_0)$에 속한다고 가정하자.
$f=F(1,x,y)\in k[x,y]$라 두자. 그러면 % P07-ERR-0300
$Z\cap D_+(X_0)$는 $k[x,y]/(f)$의 스펙트럼과 동형이다. $f_x,f_y$를 각각
$x,y$에 대한 $f$의 편도함수라 하자. $z'$은 $Z/k$의 매끄러운 점이므로
$f_x$ 또는 $f_y$가 $z'$에서 영이 아니다(대수학, 절
\ref{algebra-section-smooth}의 논의를 보라). % P07-ERR-0301
좌표를 다시 번호 매겨 $f_y$가 $z'$에서 영이 아니라고 가정해도 된다.
따라서 열린 부분스킴 % P07-ERR-0300
$V\subset Z\cap D_+(X_0)$가 존재하여 사영
$$
p:V\longrightarrow\Spec(k[x])
$$
은 \`etale이다. $y$의 다항식으로 본 $f$의 차수는 $d$ 이하이므로 사영 $p$의
올들의 차수는 $d$ 이하이다(스킴의 사상, 절
\ref{morphisms-section-universally-bounded}의 논의를 보라). 또한 $p$가
\'etale이므로 $p$의 상은 열린집합 $U\subset\Spec(k[x])$이다. 마지막으로
$k$가 무한체이므로 $U$의 $k$-유리점들의 집합 $U(k)$는 무한이고, 특히
공집합이 아니다. 아무 $t\in U(k)$나 택하고 $t$로 가는 점 $z\in V$를 택하면
$z$가 원하는 점이다.
\end{proof}




\section{종수 영인 곡선}
\label{curves-section-genus-zero}

\noindent
뒤에서 고유한 종수 영인 곡선의 모양을 알아야 한다. Gorenstein 고유 종수 영
곡선은 차수 $2$인 평면곡선, 즉 원뿔곡선임을 알 수 있다. 보조정리
\ref{curves-lemma-genus-zero}를 보라. 일반적인 고유 종수 영 곡선은 (더 큰 체 위의)
비특이 곡선에서 밀어붙임 절차를 통해 얻어진다. 보조정리
\ref{curves-lemma-genus-zero-singular}를 보라. 비특이 곡선은 Gorenstein이므로 이 두
결과가 가능한 모든 경우를 포괄한다.

\begin{lemma}
\label{curves-lemma-genus-zero-pic}
$X$를 $H^0(X,\mathcal{O}_X)=k$인 체 $k$ 위의 고유 곡선이라 하자. $X$의
종수가 $0$이면 차수 $0$인 모든 가역 $\mathcal{O}_X$-가군 $\mathcal{L}$은 자명하다.
\end{lemma}

\begin{proof}
Riemann--Roch(보조정리 \ref{curves-lemma-rr})에 의해
$\dim_kH^0(X,\mathcal{L})\geq0+1-0=1$이다. 따라서 $\mathcal{L}$은 영이 아닌
단면을 가지며, 다양체, 보조정리
\ref{varieties-lemma-check-invertible-sheaf-trivial}에 의해
$\mathcal{L}\cong\mathcal{O}_X$이다.
\end{proof}

\begin{lemma}
\label{curves-lemma-genus-zero-positive-degree}
$X$를 $H^0(X,\mathcal{O}_X)=k$인 체 $k$ 위의 고유 곡선이라 하자. $X$의
종수가 $0$이라고 가정하자. $\mathcal{L}$을 차수 $d>0$인 가역
$\mathcal{O}_X$-가군이라 하자. 그러면 다음이 성립한다.
\begin{enumerate}
\item $\dim_kH^0(X,\mathcal{L})=d+1$이고 $\dim_kH^1(X,\mathcal{L})=0$이다.
\item $\mathcal{L}$은 매우 풍부하고 $\mathbf{P}^d_k$로 가는 닫힌몰입을 정의한다.
\end{enumerate}
\end{lemma}

\begin{proof}
차수와 종수의 정의에 의해
$$
\dim_kH^0(X,\mathcal{L})-\dim_kH^1(X,\mathcal{L})=d+1
$$
이다. $s$를 $\mathcal{L}$의 영이 아닌 단면이라 하자. 그러면 $s$의 영점
스킴은 유효 Cartier 인자 $D\subset X$이고, $\mathcal{L}=\mathcal{O}_X(D)$이며
짧은 완전열
$$
0\to\mathcal{O}_X\to\mathcal{L}\to\mathcal{L}|_D\to0
$$
이 있다. 인자, 보조정리 \ref{divisors-lemma-characterize-OD}와 주의
\ref{divisors-remark-ses-regular-section}을 보라. 가정에 의해
$H^1(X,\mathcal{O}_X)=0$이므로
$H^0(X,\mathcal{L})\to H^0(X,\mathcal{L}|_D)$는 전사이다.
$\mathcal{L}|_D$는 전역 생성되므로($\dim(D)=0$이기 때문이다. 다양체,
보조정리 \ref{varieties-lemma-chi-tensor-finite} 참조) 가역 가군 $\mathcal{L}$도
전역 생성된다. 사실 $D$는 Artin 스킴이므로
$\mathcal{L}|_D\cong\mathcal{O}_D$이고,\footnote{우리의 경우에는
$D\to\Spec(k)$가 유한이므로 인자, 보조정리
\ref{divisors-lemma-finite-trivialize-invertible-upstairs}에서 따른다.}
따라서 $D$에 제한했을 때 $\mathcal{L}|_D$를 생성하는 $\mathcal{L}$의 단면
$t$를 찾을 수 있다. % P07-ERR-0302
위 짧은 완전열은 $H^1(X,\mathcal{L})=0$임도 보여 준다.

\medskip\noindent
$n\geq1$에 대하여 곱셈사상
$$
\mu_n:H^0(X,\mathcal{L})\otimes_kH^0(X,\mathcal{L}^{\otimes n})
\longrightarrow H^0(X,\mathcal{L}^{\otimes n+1})
$$
을 생각하자. 이것이 전사라고 주장한다. 이를 보이기 위해 짧은 완전열
$$
0\to\mathcal{L}^{\otimes n}\xrightarrow{s}\mathcal{L}^{\otimes n+1}
\to\mathcal{L}^{\otimes n+1}|_D\to0
$$
을 생각하자. 이 열의 왼쪽에서 오는 $\mathcal{L}^{\otimes n+1}$의 단면들은
$\mu_n$의 상 안에 있다. 한편
$H^0(\mathcal{L})\to H^0(\mathcal{L}|_D)$가 전사이고 $t^n$은
$\mathcal{L}^{\otimes n}|_D$의 한 자명화로 가므로,
$\mu_n(H^0(X,\mathcal{L})\otimes t^n)$은
$H^0(X,\mathcal{L}^{\otimes n+1})$의 부분공간으로서
$\mathcal{L}^{\otimes n+1}|_D$의 전역단면 위로 전사한다. 이로써 주장이 증명된다.

\medskip\noindent
다양체, 보조정리 \ref{varieties-lemma-ample-curve}에 의해 $\mathcal{L}$은 풍부하다.
따라서 스킴의 사상, 보조정리 \ref{morphisms-lemma-proper-ample-is-proj}는 동형
$$
X\longrightarrow\text{Proj}\left(
\bigoplus\nolimits_{n\geq0}H^0(X,\mathcal{L}^{\otimes n})\right)
$$
을 준다. 모든 $n\geq1$에 대하여 $\mu_n$이 전사이므로 오른쪽의 등급대수는
$H^0(X,\mathcal{L})$ 위의 대칭대수의 몫이다. $H^0(X,\mathcal{L})$의 $k$-기저
$s_0,\ldots,s_d$를 택하면 이는 $d+1$개 변수의 다항식대수의 몫임을 알 수
있다. 등급환의 몫은 $\text{Proj}$의 닫힌몰입에 대응하므로(구성,
보조정리 \ref{constructions-lemma-surjective-graded-rings-generated-degree-1-map-proj})
닫힌몰입 $X\to\mathbf{P}^d_k$를 얻는다. 이 사상이 $\mathcal{L}$의 단면
$s_0,\ldots,s_d$에 대응하는 구성, 보조정리
\ref{constructions-lemma-projective-space}의 사상이라는 확인은 생략한다.
\end{proof}

\begin{lemma}
\label{curves-lemma-genus-zero}
$X$를 $H^0(X,\mathcal{O}_X)=k$인 체 $k$ 위의 고유 곡선이라 하자. $X$가
Gorenstein이고 종수가 $0$이면 $X$는 차수 $2$인 평면곡선과 동형이다.
\end{lemma}

\begin{proof}
$\omega_X$를 보조정리 \ref{curves-lemma-duality-dim-1}에서와 같이 두고, 가역층
$\mathcal{L}=\omega_X^{\otimes-1}$을 생각하자. 보조정리
\ref{curves-lemma-genus-gorenstein}에 의해 $\deg(\omega_X)=-2$이고 따라서
$\deg(\mathcal{L})=2$이다. 보조정리 \ref{curves-lemma-genus-zero-positive-degree}에 의해
$\mathcal{L}$의 전역단면들로 이루어진 $k$-벡터공간의 기저 $s_0,s_1,s_2$를 택하면 닫힌몰입
$$
\varphi_{(\mathcal{L},(s_0,s_1,s_2))}:X\longrightarrow\mathbf{P}^2_k
$$
를 얻는다. 따라서 $X$는 어떤 차수 $d$의 평면곡선이다. 그 방정식을
$F\in k[T_0,T_1,T_2]_d$라 하자(보조정리
\ref{curves-lemma-equation-plane-curve}). $X$의 종수가 $0$이므로 $d$는 $1$ 또는
$2$이다(보조정리 \ref{curves-lemma-genus-plane-curve}). $F$는 $\varphi(X)$에 제한하면
영단면이므로 $F(s_0,s_1,s_2)$는 % P07-ERR-0303
$\mathcal{L}^{\otimes2}$의 영단면이다. $s_0,s_1,s_2$는 선형독립이므로 $F$는
일차식일 수 없다. 즉 $d=\deg(F)\geq2$이다. 따라서 $d=2$이고 증명이 끝난다.
\end{proof}

\begin{proposition}[사영직선의 특징짓기]
\label{curves-proposition-projective-line}
$k$를 체, $X$를 $k$ 위의 고유 곡선이라 하자. 다음은 동치이다.
\begin{enumerate}
\item $X\cong\mathbf{P}^1_k$이다.
\item $X$는 $k$ 위에서 매끄럽고 기하적으로 기약이며, $X$의 종수가 $0$이고,
$X$는 홀수 차수인 가역 가군을 가진다.
\item $X$는 $k$ 위에서 기하적으로 정역이고 $X$의 종수가 $0$이며 $X$는
Gorenstein이고, $X$는 홀수 차수인 가역층을 가진다.
\item $H^0(X,\mathcal{O}_X)=k$이고 $X$의 종수가 $0$이며 $X$가 Gorenstein이고,
 $X$는 홀수 차수인 가역층을 가진다.
\item $X$는 $k$ 위에서 기하적으로 정역이고 $X$의 종수가 $0$이며, $X$는 차수 $1$인
가역 $\mathcal{O}_X$-가군을 가진다.
\item $H^0(X,\mathcal{O}_X)=k$이고 $X$의 종수가 $0$이며, $X$는 차수 $1$인 가역
$\mathcal{O}_X$-가군을 가진다.
\item $H^1(X,\mathcal{O}_X)=0$이고 $X$는 차수 $1$인 가역
$\mathcal{O}_X$-가군을 가진다.
\item $H^1(X,\mathcal{O}_X)=0$이고, $\mathcal{O}_{X,x_i}$가 정규이며
$\gcd([\kappa(x_i):k])=1$인 닫힌점 $x_1,\ldots,x_n$이 $X$에 존재한다.
\item % P07-ERR-0304
여기에 더 추가할 것.
\end{enumerate}
\end{proposition}

\begin{proof}
조건 (2)--(8)이 각각 (1)을 함의함을 증명하고, (1)이 (2)--(8)을 함의한다는
확인은 생략하겠다.

\medskip\noindent
(2)를 가정하자. $k$ 위의 매끄러운 스킴은 기하적으로 환원이고(다양체,
보조정리 \ref{varieties-lemma-smooth-geometrically-normal}) 정칙이다(다양체,
보조정리 \ref{varieties-lemma-smooth-regular}). 따라서 $X$는 Gorenstein이다
(스킴의 쌍대성, 보조정리 \ref{duality-lemma-regular-gorenstein}).
그러므로 (3)으로 환원된다.

\medskip\noindent
(3)을 가정하자. $X$는 $k$ 위에서 기하적으로 정역이므로 다양체, 보조정리
\ref{varieties-lemma-regular-functions-proper-variety}에 의해
$H^0(X,\mathcal{O}_X)=k$이다. 따라서 (4)로 환원된다. % P07-ERR-0305

\medskip\noindent
(4)를 가정하자. $X$는 Gorenstein이므로 보조정리 \ref{curves-lemma-duality-dim-1}의
쌍대화 가군 $\omega_X$는 보조정리 \ref{curves-lemma-genus-gorenstein}에 의해 차수가
$\deg(\omega_X)=-2$이다. 홀수 차수인 가역 가군이 존재한다는 가정과 결합하면
차수 $1$인 가역 가군이 존재한다. 이로써 (6)으로 환원된다.

\medskip\noindent
(5)를 가정하자. $X$는 $k$ 위에서 기하적으로 정역이므로 다양체, 보조정리
\ref{varieties-lemma-regular-functions-proper-variety}에 의해
$H^0(X,\mathcal{O}_X)=k$이다. 따라서 (6)으로 환원된다. % P07-ERR-0305

\medskip\noindent
(6)을 가정하자. 보조정리 \ref{curves-lemma-genus-zero-positive-degree}에 의해
$X\cong\mathbf{P}^1_k$이다.

\medskip\noindent
(7)을 가정하자. $\kappa=H^0(X,\mathcal{O}_X)$는 $k$ 위의 유한 체확장이다.
다양체, 보조정리 \ref{varieties-lemma-regular-functions-proper-variety}를 보라.
$d=[\kappa:k]>1$이면 모든 가역층의 차수는 $d$로 나누어지므로 차수 $1$인
가역층이 존재할 수 없다. 따라서 $d=1$이고 (6)의 경우로 환원된다.

\medskip\noindent
(8)을 가정하자. $\kappa=H^0(X,\mathcal{O}_X)$는 $k$ 위의 유한 체확장이다.
다양체, 보조정리 \ref{varieties-lemma-regular-functions-proper-variety}를 보라.
$\kappa\subset\kappa(x_i)$이므로 차수들의 최대공약수에 대한 가정에 의해
$k=\kappa$이다. 같은 조건에 의해
$1=\sum a_i[\kappa(x_i):k]$가 되도록 정수 $a_i$를 찾을 수 있다.
각 $x_i$는 다양체, 보조정리
\ref{varieties-lemma-regular-point-on-curve}에 의해 $X$ 위의 유효 Cartier
인자를 정의하므로, 가역 가군 $\mathcal{O}_X(\sum a_ix_i)$를 생각할 수 있다.
우리가 고른 $a_i$들에 의해 % P07-ERR-0306
$\mathcal{L}$의 차수는 $1$이다. 따라서 보조정리
\ref{curves-lemma-genus-zero-positive-degree}에 의해 $X\cong\mathbf{P}^1_k$이다.
\end{proof}

\begin{lemma}
\label{curves-lemma-genus-zero-singular}
$X$를 $H^0(X,\mathcal{O}_X)=k$인 체 $k$ 위의 고유 곡선이라 하자. $X$가
특이이고 종수가 $0$이라고 가정하자. 그러면 다음 그림이 존재한다.
$$
\xymatrix{
x' \ar[d] \ar[r] & X' \ar[d]^\nu \ar[r] & \Spec(k') \ar[d] \\
x \ar[r] & X \ar[r] & \Spec(k)
}
$$
여기서
\begin{enumerate}
\item $k'/k$는 자명하지 않은 유한확장이다.
\item $X'\cong\mathbf{P}^1_{k'}$이다.
\item $x'$은 $X'$의 $k'$-유리점이다.
\item $x$는 $X$의 $k$-유리점이다.
\item $X'\setminus\{x'\}\to X\setminus\{x\}$는 동형이다.
\item $k'/k=\kappa(x')/\kappa(x)$가 점 $x$ 위의 마천루층을 나타낼 때
$0\to\mathcal{O}_X\to\nu_*\mathcal{O}_{X'}\to k'/k\to0$은 짧은 완전열이다.
\end{enumerate}
\end{lemma}

\begin{proof}
$\nu:X'\to X$를 $X$의 정규화라 하자. 다양체, 절
\ref{varieties-section-normalization}과
\ref{varieties-section-normalization-one-dimensional}을 보라. $X$가 특이이므로
$\nu$는 동형이 아니다. 그러면
$k'=H^0(X',\mathcal{O}_{X'})$는 $k$의 유한확장이다(다양체, 보조정리
\ref{varieties-lemma-regular-functions-proper-variety}). 짧은 완전열
$$
0\to\mathcal{O}_X\to\nu_*\mathcal{O}_{X'}\to\mathcal{Q}\to0
$$
과 $\mathcal{Q}$가 유한개의 닫힌점에 지지된다는 사실로부터 다음을 얻는다.
\begin{enumerate}
\item $H^1(X',\mathcal{O}_{X'})=0$, 즉 $k'$ 위의 곡선으로서 $X'$의 종수는 $0$이다.
\item 짧은 완전열 $0\to k\to k'\to H^0(X,\mathcal{Q})\to0$이 있다.
\end{enumerate}
특히 $k'/k$는 자명하지 않은 확장이다.

\medskip\noindent
이제 흔히 {\it 도체 아이디얼}이라 불리는
$$
\mathcal{I}=\SheafHom_{\mathcal{O}_X}(\nu_*\mathcal{O}_{X'},\mathcal{O}_X)
$$
를 생각하자. 이는 준연접 $\mathcal{O}_X$-가군이다. 사상
$\varphi\mapsto\varphi(1)$을 통해 $\mathcal{I}$을 $\mathcal{O}_X$의 아이디얼로
본다. 따라서 $\mathcal{I}(U)$는
$f\left(\nu_*\mathcal{O}_{X'}(U)\right)\subset\mathcal{O}_X(U)$인
$f\in\mathcal{O}_X(U)$들의 집합이다. 달리 말하면 $f$가 $\mathcal{Q}$를
소멸시킨다는 조건이다. 동치로 다음 정의 완전열이 있다. % P07-ERR-0309
$$
0\to\mathcal{I}\to\mathcal{O}_X\to
\SheafHom_{\mathcal{O}_X}(\mathcal{Q},\mathcal{Q}).
$$
$U\subset X$를 $\mathcal{Q}$의 지지를 포함하는 아핀 열린집합이라 하자.
그러면 $V=\mathcal{Q}(U)=H^0(X,\mathcal{Q})$는 차원 $n-1$인 $k$-벡터공간이다.
% P07-ERR-0307
$\mathcal{O}_X(U)\to\Hom_k(V,V)$의 상은 가환 부분대수이므로 $k$ 위의 차원이
$\leq n-1$이다(이는 행렬대수의 가환 부분대수의 성질이다. 세부사항은
생략한다). % P07-ERR-0308
따라서 짧은 완전열
$$
0\to\mathcal{I}\to\mathcal{O}_X\to\mathcal{A}\to0
$$
이 있으며 $\text{Supp}(\mathcal{A})=\text{Supp}(\mathcal{Q})$이고
$\dim_kH^0(X,\mathcal{A})\leq n-1$이다. 한편 표현
$\mathcal{I}=\SheafHom_{\mathcal{O}_X}(\nu_*\mathcal{O}_{X'},\mathcal{O}_X)$은
$\mathcal{I}$에 $\nu_*\mathcal{O}_{X'}$-가군 구조를 주며, 이에 대해 포함사상
$\mathcal{I}\to\nu_*\mathcal{O}_{X'}$는 $\nu_*\mathcal{O}_{X'}$-가군 사상이다.
그러므로 어떤 준연접 아이디얼층
$\mathcal{I}'\subset\mathcal{O}_{X'}$에 대하여
$\mathcal{I}=\nu_*\mathcal{I}'$이다. 스킴의 사상, 보조정리
\ref{morphisms-lemma-affine-equivalence-modules}를 보라. $\mathcal{A}'$을 다음
여핵으로 정의하자.
$$
0\to\mathcal{I}'\to\mathcal{O}_{X'}\to\mathcal{A}'\to0.
$$
지금까지의 완전열들을 결합하면 짧은 완전열
$0\to\mathcal{A}\to\nu_*\mathcal{A}'\to\mathcal{Q}\to0$을 얻는다. 위의 추정과
$\dim_kH^0(X,\mathcal{Q})=n-1$을 결합하면
$$
\dim_kH^0(X',\mathcal{A}')=
\dim_kH^0(X,\mathcal{A})+\dim_kH^0(X,\mathcal{Q})\leq2n-2
$$
이다. 그러나 $X'$은 $k'$ 위의 곡선이므로 좌변은 $n$으로 나누어떨어진다
(다양체, 보조정리 \ref{varieties-lemma-divisible}). $\mathcal{A}$와
$\mathcal{A}'$은 영일 수 없으므로 $\dim_kH^0(X',\mathcal{A}')=n$이고, 이는
$\mathcal{I}'$이 어떤 $k'$-유리점 $x'$의 아이디얼층이라는 뜻이다.
명제 \ref{curves-proposition-projective-line}에 의해
$X'\cong\mathbf{P}^1_{k'}$이다. 위 등식들로 돌아가면
$\dim_kH^0(X,\mathcal{A})=1$을 얻는다. 이는 $\mathcal{I}$이 어떤 $k$-유리점
$x$의 아이디얼층이라는 뜻이다. 그러면 마천루층으로서
$\mathcal{A}=\kappa(x)=k$이고 $\mathcal{A}'=\kappa(x')=k'$이다. 위의 완전열들을
비교하면 보조정리에 명시된 구조층에 관한 결과를 곧바로 얻는다.
\end{proof}

\begin{example}
\label{curves-example-squish-on-P1}
사실 보조정리 \ref{curves-lemma-genus-zero-singular}에 기술된 상황은 자명하지 않은 모든
유한확장 $k'/k$에 대하여 일어난다. 실제로
$$
A=\{f\in k'[x]\mid f(0)\in k\}
$$
를 생각하자. $A$의 스펙트럼은 아핀곡선이며, $x^{-1}$을 $y$로 보내는 동형
$A_x\cong B_y$를 써서 이를 $B=k'[y]$의 스펙트럼과 붙일 수 있다. 그 결과는
$H^0(X,\mathcal{O}_X)=k$이고 극대 아이디얼 $A\cap(x)$에 대응하는 특이점 $x$를
갖는 고유 곡선 $X$이다. 보조정리에서와 꼭 마찬가지로 $X$의 정규화는
$\mathbf{P}^1_{k'}$이다.
\end{example}







\section{기하종수}
\label{curves-section-geometric-genus}

\noindent
$X$가 $H^0(X,\mathcal{O}_X)=k$인 $k$ 위의 고유하고 {\bf 매끄러운} 곡선이면
$$
p_g(X)=\dim_kH^0(X,\Omega_{X/k})
$$
를 $X$의 {\it 기하종수}라고 한다. 보조정리 \ref{curves-lemma-genus-smooth}에 의해
$X$의 기하종수는 (산술)종수와 일치한다. 그러나 더 높은 차원에서는 기하종수와
산술종수가 서로 다르다. 주의 \ref{curves-remark-genus-higher-dimension}을 보라.

\medskip\noindent
특이 곡선에 대해서는 기하종수를 다음과 같이 정의한다.

\begin{definition}
\label{curves-definition-geometric-genus}
$k$를 체, $X$를 $k$ 위의 기하적으로 기약인 곡선이라 하자. $X$의
{\it 기하종수}는 보조정리 \ref{curves-lemma-smooth-models}에서와 같이 필요하면 $k$의
확장체 위에서 정의된 $X$의 매끄러운 사영모형의 종수이다.
\end{definition}

\noindent
$k$가 완전체이면 $X$의 비특이 사영모형 $Y$는 매끄럽고(보조정리
\ref{curves-lemma-nonsingular-model-smooth}), $X$의 기하종수는 단순히 $Y$의 종수이다.
그러나 $k$가 완전체가 아니면 이는 참이 아닐 수 있다. 이 경우
$(X_K)_{red}$의 비특이 사영모형 $Y_K$가 매끄러운 사영곡선이 되도록 하는
확장 $K/k$를 택하고, $X$의 기하종수를 $Y_K$의 종수로 정의한다. 보조정리
\ref{curves-lemma-smooth-models}와 \ref{curves-lemma-genus-base-change}에 의해 이는 잘 정의된다.

\begin{remark}
\label{curves-remark-genus-higher-dimension}
$X$가 대수적으로 닫힌 체 $k$ 위의 $d$차원 고유 매끄러운 다양체라고 하자.
그러면 {\it 산술종수}는 흔히
$p_a(X)=(-1)^d(\chi(X,\mathcal{O}_X)-1)$로, {\it 기하종수}는
$p_g(X)=\dim_kH^0(X,\Omega^d_{X/k})$로 정의된다. 이 상황에서는 여전히
$\omega_X\cong\Omega^d_{X/k}$이지만 산술종수와 기하종수가 더는 일치하지 않는다.
예를 들어 $d=2$이면
\begin{align*}
p_a(X)-p_g(X)&=
h^0(X,\mathcal{O}_X)-h^1(X,\mathcal{O}_X)+h^2(X,\mathcal{O}_X)-1
-h^0(X,\Omega^2_{X/k})\\
&=-h^1(X,\mathcal{O}_X)+h^2(X,\mathcal{O}_X)-h^0(X,\omega_X)\\
&=-h^1(X,\mathcal{O}_X)
\end{align*}
이다. 여기서 $h^i(X,\mathcal{F})=\dim_kH^i(X,\mathcal{F})$이고 마지막 등호는
쌍대성에서 따른다. 따라서 곡면에서는 차 $p_g(X)-p_a(X)$가 항상 음이 아니며,
이를 때때로 곡면의 비정칙도라고 한다. $X=C_1\times C_2$가 종수가 각각
$g_1,g_2$인 매끄러운 사영곡선들의 곱이면 비정칙도는 $g_1+g_2$이다.
\end{remark}





\section{Riemann--Hurwitz}
\label{curves-section-riemann-hurewitz}

\noindent
$k$를 체라 하고 $f:X\to Y$를 $k$ 위의 매끄러운 곡선들의 사상이라 하자.
그러면 스킴의 사상, 보조정리 \ref{morphisms-lemma-triangle-differentials}에 의해
표준 완전열
$$
f^*\Omega_{Y/k}\xrightarrow{\text{d}f}\Omega_{X/k}
\longrightarrow\Omega_{X/Y}\longrightarrow0
$$
을 얻는다. $X$와 $Y$가 매끄러우므로 층 $\Omega_{X/k}$와 $\Omega_{Y/k}$는
가역 가군이다. 스킴의 사상, 보조정리
\ref{morphisms-lemma-smooth-omega-finite-locally-free}를 보라. 첫 사상이 영이
아니라고, 즉 $f$가 일반적으로 \'etale이라고 가정하자. 보조정리
\ref{curves-lemma-generically-etale}를 보라. $R\subset X$를 $f$의 different
$\mathfrak{D}_f$가 잘라내는 닫힌 부분스킴이라 하자. 판별식, 보조정리
\ref{discriminant-lemma-discriminant-quasi-finite-morphism-smooth}에 의해 이는
$\text{d}f$의 소멸자리와 같고 유효 Cartier 인자이며,
$$
f^*\Omega_{Y/k}\otimes_{\mathcal{O}_X}\mathcal{O}_X(R)=\Omega_{X/k}
$$
를 얻는다. 특히 $X,Y$가 사영이고
$k=H^0(Y,\mathcal{O}_Y)=H^0(X,\mathcal{O}_X)$이며 $X,Y$의 종수가 각각
$g_X,g_Y$이면 Riemann--Hurwitz 공식
\begin{align*}
2g_X-2&=\deg(\Omega_{X/k})\\
&=\deg(f^*\Omega_{Y/k}\otimes_{\mathcal{O}_X}\mathcal{O}_X(R))\\
&=\deg(f)\deg(\Omega_{Y/k})+\deg(R)\\
&=\deg(f)(2g_Y-2)+\deg(R)
\end{align*}
을 얻는다. 첫째와 마지막 등호는 보조정리 \ref{curves-lemma-genus-smooth}에서 따른다.
둘째 등호는 위에 주어진 가역층들의 동형에서 따른다. 셋째 등호는 차수의
가법성(다양체, 보조정리 \ref{varieties-lemma-degree-tensor-product}), 당김의
차수 공식(다양체, 보조정리
\ref{varieties-lemma-degree-pullback-map-proper-curves}), 그리고 마지막으로
$\mathcal{O}_X(R)$의 차수 공식(다양체, 보조정리
\ref{varieties-lemma-degree-effective-Cartier-divisor})에서 따른다.
% P07-ERR-0310

\medskip\noindent
Riemann--Hurwitz 공식을 사용하려면
$\deg(R)=\dim_k\Gamma(R,\mathcal{O}_R)$를 계산해야 한다. $k$ 위 영차원 스킴의
구조(예를 들어 다양체, 보조정리
\ref{varieties-lemma-algebraic-scheme-dim-0} 참조)에 의해 $R$은 $k$ 위에서
각 $\mathcal{O}_{R,x}$가 유한차원인 Artin 국소환들의 스펙트럼의 유한 분리합
$R=\coprod_{x\in R}\Spec(\mathcal{O}_{R,x})$이다. 따라서
$$
\deg(R)=\sum\nolimits_{x\in R}\dim_k\mathcal{O}_{R,x}
=\sum\nolimits_{x\in R}d_x[\kappa(x):k]
$$
이고, 여기서
$$
d_x=\text{length}_{\mathcal{O}_{R,x}}\mathcal{O}_{R,x}
=\text{length}_{\mathcal{O}_{X,x}}\mathcal{O}_{R,x}
$$
는 $R$에서 $x$의 중복도이다(대수학, 보조정리
\ref{algebra-lemma-pushdown-module} 참조). 상이 $y\in Y$인 닫힌점 $x\in X$를
잡자. 줄기를 살펴보면 완전열
$$
\Omega_{Y/k,y}\to\Omega_{X/k,x}\to\Omega_{X/Y,x}\to0
$$
을 얻는다. (계수 $1$인 자유) 가군 $\Omega_{X/k,x}$와 $\Omega_{Y/k,y}$의
국소 생성원을 각각 $\eta_x$, $\eta_y$로 택하면, 어떤 영이 아닌
$h\in\mathcal{O}_{X,x}$에 대하여 $\eta_y\mapsto h\eta_x$이다. 위 완전열에 의해
$\mathcal{O}_{X,x}$-가군으로서
$\Omega_{X/Y,x}\cong\mathcal{O}_{X,x}/h\mathcal{O}_{X,x}$이다. 인자 $R$은
$h$로 잘라내지므로 $\mathcal{O}_{R,x}=\mathcal{O}_{X,x}/h\mathcal{O}_{X,x}$이다.
따라서 다음 등식들을 얻는다.
\begin{align*}
d_x&=\text{length}_{\mathcal{O}_{X,x}}(\mathcal{O}_{R,x})\\
&=\text{length}_{\mathcal{O}_{X,x}}(\mathcal{O}_{X,x}/h\mathcal{O}_{X,x})\\
&=\text{length}_{\mathcal{O}_{X,x}}(\Omega_{X/Y,x})\\
&=\text{ord}_{\mathcal{O}_{X,x}}(h)\\
&=\text{ord}_{\mathcal{O}_{X,x}}(``\eta_y/\eta_x")
\end{align*}
첫째 등호는 $d_x$의 정의에 의한다. 둘째와 셋째 등호는 위에서 보았다.
넷째 등호는 $\text{ord}$의 정의이다. 대수학, 정의
\ref{algebra-definition-ord}를 보라. $\mathcal{O}_{X,x}$는 이산값매김환이므로
정수 $\text{ord}_{\mathcal{O}_{X,x}}(h)$는 바로 $h$의 값매김이다. 다섯째
등호는 기억을 돕기 위한 표기이다. % P07-ERR-0311

\medskip\noindent
다른 불변량들로 중복도 $d_x$를 ``계산''할 수 있는 경우를 보자.
$\kappa(x)$가 $k$ 위에서 분리이면, $\mathcal{O}_{X,x}$와
$\mathcal{O}_{Y,y}$의 균등화원 $s,t$에 대하여 각각
$\eta_x=\text{d}s$, $\eta_y=\text{d}t$로 택할 수 있다(보조정리
\ref{curves-lemma-uniformizer-works}). 확장
$\mathcal{O}_{Y,y}\subset\mathcal{O}_{X,x}$의 분기지수를 $e_x$라 하면 어떤 단위
$u\in\mathcal{O}_{X,x}$에 대하여 $t\mapsto us^{e_x}$이다. 따라서
% P07-ERR-0312
$$
\eta_y=\text{d}t=\text{d}(us^{e_x})
=e s^{e_x-1}u\text{d}s+s^{e_x}\text{d}u.
$$
어떤 $w\in\mathcal{O}_{X,x}$에 대하여 $\text{d}u=w\text{d}s$라 쓰면
% P07-ERR-0312
$$
``\eta_y/\eta_x"=e s^{e_x-1}u+s^{e_x}w=(e_xu+sw)s^{e_x-1}.
$$
따라서 $\kappa(x)$의 표수가 $p>0$이고 $p$가 $e_x$를 나누는 경우를 제외하면
이것의 소멸차수는 $e_x-1$이고, 그 예외적인 경우에는 $>e_x-1$이다.

\medskip\noindent
위의 모든 것을 결합하면 $k$의 표수가 영일 때
$$
2g_X-2=(2g_Y-2)\deg(f)+
\sum\nolimits_{x\in X}(e_x-1)[\kappa(x):k]
$$
임을 얻는다. 여기서 $e_x$는 $\mathcal{O}_{Y,f(x)}$ 위의
$\mathcal{O}_{X,x}$의 분기지수이다. 이 정확한 공식은 모든 분기가 순한 경우,
즉 잉여체 확장 $\kappa(x)/\kappa(y)$가 분리이고 $e_x$가 $k$의 표수와 서로
소인 경우에만 성립하지만, 위의 논증만으로는 이를 증명하기에 충분하지 않다.
보조정리 \ref{curves-lemma-rhe}와 그 증명을 참조하라.

\begin{lemma}
\label{curves-lemma-generically-etale}
\begin{slogan}
매끄러운 곡선들의 사상이 분리일 필요충분조건은 거의 모든 곳에서
\'etale인 것이다. % P07-ERR-0313
\end{slogan}
$k$를 체라 하고 $f:X\to Y$를 $k$ 위의 매끄러운 곡선들의 사상이라 하자.
다음은 동치이다.
\begin{enumerate}
\item $\text{d}f:f^*\Omega_{Y/k}\to\Omega_{X/k}$는 영이 아니다.
\item $\Omega_{X/Y}$는 $X$의 진 닫힌 부분집합에 지지된다.
\item 공집합이 아닌 열린집합 $U\subset X$가 존재하여
$f|_U:U\to Y$는 비분기이다.
\item 공집합이 아닌 열린집합 $U\subset X$가 존재하여
$f|_U:U\to Y$는 \'etale이다.
\item 함수체의 확장 $k(X)/k(Y)$는 유한 분리이다.
\end{enumerate}
\end{lemma}

\begin{proof}
$X$와 $Y$가 매끄러우므로 층 $\Omega_{X/k}$와 $\Omega_{Y/k}$는 가역
가군이다. 스킴의 사상, 보조정리
\ref{morphisms-lemma-smooth-omega-finite-locally-free}를 보라. 스킴의 사상,
보조정리 \ref{morphisms-lemma-triangle-differentials}의 완전열
$$
f^*\Omega_{Y/k}\longrightarrow\Omega_{X/k}
\longrightarrow\Omega_{X/Y}\longrightarrow0
$$
을 쓰면 (1)과 (2)가 동치이고, 이들은
$f^*\Omega_{Y/k}\to\Omega_{X/k}$가 일반점에서 영이 아니라는 조건과
동치임을 알 수 있다. (2)와 (3)의 동치는 스킴의 사상, 보조정리
\ref{morphisms-lemma-unramified-omega-zero}에서 따른다. (3)과 (4)의 동치는
스킴의 사상, 보조정리 \ref{morphisms-lemma-flat-unramified-etale} 및 평탄성이
자동이라는 사실(보조정리 \ref{curves-lemma-flat})에서 따른다. (5)와 (4)의 동치는
대수학, 보조정리 \ref{algebra-lemma-smooth-at-generic-point}를 쓰면 된다. 일부
세부사항은 생략한다.
\end{proof}

\begin{lemma}
\label{curves-lemma-rh}
$f:X\to Y$를 체 $k$ 위의 매끄러운 고유 곡선들의 사상으로서 보조정리
\ref{curves-lemma-generically-etale}의 동치인 조건들을 만족한다고 하자.
$k=H^0(Y,\mathcal{O}_Y)=H^0(X,\mathcal{O}_X)$이고 $X,Y$의 종수가 각각
$g_X,g_Y$이면
$$
2g_X-2=(2g_Y-2)\deg(f)+\deg(R)
$$
이다. 여기서 $R\subset X$는 $f$의 different가 잘라내는 유효 Cartier 인자이다.
\end{lemma}

\begin{proof}
위의 논의를 보라. 우리는 판별식, 보조정리
\ref{discriminant-lemma-discriminant-quasi-finite-morphism-smooth}, 보조정리
\ref{curves-lemma-genus-smooth}, 그리고 다양체, 보조정리
\ref{varieties-lemma-degree-tensor-product}와
\ref{varieties-lemma-degree-pullback-map-proper-curves}를 사용하였다.
\end{proof}

\begin{lemma}
\label{curves-lemma-uniformizer-works}
$X\to\Spec(k)$가 닫힌점 $x\in X$에서 상대차원 $1$인 매끄러운 사상이라 하자.
$\kappa(x)$가 $k$ 위에서 분리이면, 이산값매김환 $\mathcal{O}_{X,x}$의 임의의
균등화원 $s$에 대하여 원소 $\text{d}s$는 $\mathcal{O}_{X,x}$ 위에서
$\Omega_{X/k,x}$를 자유 생성한다.
\end{lemma}

\begin{proof}
대수학, 보조정리 \ref{algebra-lemma-characterize-smooth-over-field}에 의해
$\mathcal{O}_{X,x}$는 이산값매김환이다. $x$가 닫혔으므로 $\kappa(x)$는 $k$의
유한확장이다. 따라서 $\kappa(x)/k$가 분리이면, 임의의 균등화원 $s$의 상은
$\Omega_{X/k,x}\otimes_{\mathcal{O}_{X,x}}\kappa(x)$에서 영이 아니다(대수학,
보조정리 \ref{algebra-lemma-computation-differential}). $\Omega_{X/k,x}$는
$\mathcal{O}_{X,x}$ 위의 계수 $1$인 자유 가군이므로 결과가 따른다.
\end{proof}

\begin{lemma}
\label{curves-lemma-rhe}
표기와 가정은 보조정리 \ref{curves-lemma-rh}에서와 같이 두자. 닫힌점 $x\in X$에
대하여 $d_x$를 $R$에서 $x$의 중복도라 하자. 그러면
$$
2g_X-2=(2g_Y-2)\deg(f)+\sum\nolimits d_x[\kappa(x):k]
$$
이다. 또한 다음이 성립한다.
\begin{enumerate}
\item $d_x=\text{length}_{\mathcal{O}_{X,x}}(\Omega_{X/Y,x})$이다.
\item $e_x$를 $\mathcal{O}_{Y,y}$ 위의 $\mathcal{O}_{X,x}$의 분기지수라 하면
$d_x\geq e_x-1$이다.
\item $d_x=e_x-1$일 필요충분조건은 $\mathcal{O}_{X,x}$가
$\mathcal{O}_{Y,y}$ 위에서 순하게 분기하는 것이다.
\end{enumerate}
\end{lemma}

\begin{proof}
보조정리 \ref{curves-lemma-rh}와 위의 논의(다양체, 보조정리
\ref{varieties-lemma-algebraic-scheme-dim-0} 및 대수학, 보조정리
\ref{algebra-lemma-pushdown-module}를 사용하였다)에 의해 $R$에서 $x$의
중복도 $d_x$에 관한 결과들만 증명하면 된다. (1)은 위의 논의에서 증명하였다.
그 논의에서는 $\kappa(x)$가 $k$ 위에서 분리인 경우에만 (2), (3)을
증명하였다. 증명의 나머지에서는 준유한 Gorenstein 사상의 different에 관한
내용을 써서 (2), (3)을 통일적으로 다룬다.

\medskip\noindent
먼저 $f$가 준유한 Gorenstein 사상임을 주목하자. 예를 들어 $f$가 평탄
준유한 사상이고 $X$가 Gorenstein이기 때문이다(스킴의 쌍대성, 보조정리
\ref{duality-lemma-flat-morphism-from-gorenstein-scheme} 참조). 또는 위에서
사용한 판별식, 보조정리
\ref{discriminant-lemma-discriminant-quasi-finite-morphism-smooth}의 증명에서
이를 보였다. 따라서 판별식, 보조정리
\ref{discriminant-lemma-gorenstein-quasi-finite}에 의해 $\omega_{X/Y}$는
가역이고, $X$를 열린집합으로 바꾸거나 어떤 $Y'\to Y$로 밑변환한 뒤에도
마찬가지이다. 아래에서는 더 말하지 않고 이를 사용하겠다.

\medskip\noindent
$x\in U$, $y\in V$, $f(U)\subset V$이고 $x$가 $y$ 위에 놓인 $U$의 유일한
점이 되도록 아핀 열린집합 $U\subset X$, $V\subset Y$를 택하자.
$U=\Spec(A)$, $V=\Spec(B)$라 쓰자. 그러면 $R\cap U$는
$f|_U:U\to V$의 different이다. 판별식, 보조정리
\ref{discriminant-lemma-base-change-different}에 의해 우리의 경우 different의
형성은 임의의 밑변환과 가환한다. $U,V$의 선택에 의해
$$
A\otimes_B\kappa(y)=
\mathcal{O}_{X,x}\otimes_{\mathcal{O}_{Y,y}}\kappa(y)
=\mathcal{O}_{X,x}/(s^{e_x})
$$
이다. 여기서 $e_x$는 명제의 분기지수이다. $C=\mathcal{O}_{X,x}/(s^{e_x})$를
$\kappa(y)$ 위의 유한 대수로 보고, $\mathfrak{D}_{C/\kappa(y)}$를
판별식, 정의 \ref{discriminant-definition-different}의 의미에서
$\kappa(y)$ 위의 $C$의 different라 하자. 다음을 보이면 충분하다.
$\mathfrak{D}_{C/\kappa(y)}$가 영이 아닐 필요충분조건은 확장
$\mathcal{O}_{Y,y}\subset\mathcal{O}_{X,x}$가 순하게 분기하는 것이고, 순하게
분기하는 경우 $\mathfrak{D}_{C/\kappa(y)}$는 $C$에서 $s^{e_x-1}$로 생성되는
아이디얼이다. 순한 분기란 정확히 $\kappa(x)/\kappa(y)$가 분리이고
$\kappa(y)$의 표수가 $e_x$를 나누지 않는다는 뜻임을 상기하자. 한편
$C/\kappa(y)$의 different가 영이 아닐 필요충분조건은
$\tau_{C/\kappa(y)}\in\omega_{C/\kappa(y)}$가 영이 아닌 것이다. 실제로
$\omega_{C/\kappa(y)}$는 ($\omega_{A/B}$의 밑변환이므로) 가역 $C$-가군이고,
따라서 계수 $1$인 자유 가군이다. 생성원을 $\lambda$라 하고 어떤 $h\in C$에
대하여 $\tau_{C/\kappa(y)}=h\lambda$라 쓰자. 그러면
$\mathfrak{D}_{C/\kappa(y)}=(h)\subset C$이므로 주장이 따른다. 판별식,
보조정리 \ref{discriminant-lemma-tau-nonzero}에 의해
$\tau_{C/\kappa(y)}\not =0$일 필요충분조건은 $\kappa(x)/\kappa(y)$가 분리이고
$e_x$가 표수와 서로 소인 것이다. 마지막으로 $\tau_{C/\kappa(y)}$가 영이
아니더라도 $s\tau_{C/\kappa(y)}=0$이다. 왜냐하면
$s\tau_{C/\kappa(y)}:C\to\kappa(y)$는 $c$를 멱영 연산자 $sc$의 대각합으로
보내고 그 값은 영이기 때문이다. 따라서 $sh=0$이고
$h\in(s^{e_x-1})$이다. 그러므로 항상
$\mathfrak{D}_{C/\kappa(y)}\subset(s^{e_x-1})$이다.
$(s^{e_x-1})\subset C$가 영이 아닌 가장 작은 아이디얼이므로 마지막 명제가
증명된다.
\end{proof}





\section{비분리 사상}
\label{curves-section-inseparable}

\noindent
비분리 사상 아래에서 종수가 어떻게 변하는지에 관한 몇 가지 비고를 한다.

\begin{lemma}
\label{curves-lemma-dominated-by-smooth}
$k$를 체라 하고 $f:X\to Y$를 $k$ 위의 곡선들의 전사사상이라 하자.
$X$가 $k$ 위에서 매끄럽고 $Y$가 정규이면 $Y$는 $k$ 위에서 매끄럽다.
\end{lemma}

\begin{proof}
$y\in Y$를 잡고 $y$로 가는 $x\in X$를 택하자. 다양체, 보조정리
\ref{varieties-lemma-flat-under-smooth}에 의해 $f$가 $x$에서 평탄임을 보이면
충분하다. 이는 보조정리 \ref{curves-lemma-flat}에서 따른다.
\end{proof}

\begin{lemma}
\label{curves-lemma-purely-inseparable}
$k$를 표수 $p>0$인 체라 하고 $f:X\to Y$를 $k$ 위의 고유 비특이 곡선들의
상수가 아닌 사상이라 하자. 함수체의 확장 $k(X)/k(Y)$가 순수 비분리이면
분해
$$
X=X_0\to X_1\to\ldots\to X_n=Y
$$
가 존재하여 각 $X_i$는 고유 비특이 곡선이고 $X_i\to X_{i+1}$은 차수 $p$인
사상으로서 $k(X_{i+1})\subset k(X_i)$가 비분리이다.
\end{lemma}

\begin{proof}
정리 \ref{curves-theorem-curves-rational-maps}과, 체의 유한 순수 비분리 확장은 항상
차수 $p$인 (비분리) 확장들의 열로 얻을 수 있다는 사실에서 따른다. 체,
보조정리 \ref{fields-lemma-finite-purely-inseparable}를 보라.
\end{proof}

\begin{lemma}
\label{curves-lemma-inseparable-deg-p-smooth}
$k$를 표수 $p>0$인 체라 하고 $f:X\to Y$를 $k$ 위의 고유 비특이 곡선들의
상수가 아닌 사상이라 하자. $X$가 매끄럽고 $k(Y)\subset k(X)$가 차수 $p$인
비분리 확장이면, $f$가 $F_{X/k}$가 되게 하는 유일한 동형
$Y=X^{(p)}$가 존재한다.
\end{lemma}

\begin{proof}
상대 Frobenius 사상 $F_{X/k}:X\to X^{(p)}$은 다양체, 절
\ref{varieties-section-frobenius}에서 구성하였다. $X^{(p)}$은 $X$의
밑변환이므로 $k$ 위의 매끄러운 고유 곡선이다. 다양체, 보조정리
\ref{varieties-lemma-inseparable-deg-p-smooth}에 의해 $F_{X/k}$의 차수는
$p$이다. 따라서 $k(X^{(p)})$와 $k(Y)$는 모두 $k(X)$의 부분체이고
$[k(X):k(Y)]=[k(X):k(X^{(p)})]=p$이다. 보조정리를 증명하려면 $k(X)$ 안에서
$k(Y)=k(X^{(p)})$임을 보이면 충분하다. 정리
\ref{curves-theorem-curves-rational-maps}를 보라.

\medskip\noindent
$K=k(X)$라 쓰고 사상 $\text{d}:K\to\Omega_{K/k}$를 생각하자. 보조정리
\ref{curves-lemma-generically-etale}에 의해 $k(Y)$는 $\text{d}$의 핵에 포함된다.
다양체, 보조정리 \ref{varieties-lemma-relative-frobenius-omega}에 의해
$k(X^{(p)})$도 $\text{d}$의 핵에 포함된다. % P07-ERR-0317
$X$가 매끄러운 곡선이므로 $\Omega_{K/k}$는 $K$ 위의 차원 $1$인
벡터공간이다. 그러면 대수학 심화, 보조정리
\ref{more-algebra-lemma-p-basis}는 $\Ker(\text{d})=kK^p$이고
$[K:kK^p]=p$임을 함의한다. 따라서 차수를 비교하면
$k(Y)=kK^p=k(X^{(p)})$이다.
\end{proof}

\begin{lemma}
\label{curves-lemma-purely-inseparable-smooth}
$k$를 표수 $p>0$인 체라 하고 $f:X\to Y$를 $k$ 위의 고유 비특이 곡선들의
상수가 아닌 사상이라 하자. $X$가 매끄럽고 $k(Y)\subset k(X)$가 순수
비분리이면, $f$가 $X/k$의 $n$중 상대 Frobenius가 되게 하는 유일한
$n\geq0$과 유일한 동형 $Y=X^{(p^n)}$이 존재한다.
\end{lemma}

\begin{proof}
$X/k$의 $n$중 상대 Frobenius는 다양체, 주의
\ref{varieties-remark-n-fold-relative-frobenius}에 정의되어 있다. 보조정리는
보조정리 \ref{curves-lemma-inseparable-deg-p-smooth}와
\ref{curves-lemma-purely-inseparable}을 결합하면 따른다.
\end{proof}

\begin{lemma}
\label{curves-lemma-purely-inseparable-smooth-genus}
$k$를 표수 $p>0$인 체라 하고 $f:X\to Y$를 $k$ 위의 고유 비특이 곡선들의
상수가 아닌 사상이라 하자. 다음을 가정하자.
\begin{enumerate}
\item $X$는 매끄럽다.
\item $H^0(X,\mathcal{O}_X)=k$이다.
\item $k(X)/k(Y)$는 순수 비분리이다.
\end{enumerate}
그러면 $Y$는 매끄럽고 $H^0(Y,\mathcal{O}_Y)=k$이며 $Y$의 종수는 $X$의
종수와 같다.
\end{lemma}

\begin{proof}
보조정리 \ref{curves-lemma-purely-inseparable-smooth}에 의해
$Y=X^{(p^n)}$은 $F_{\Spec(k)}^n$에 의한 $X$의 밑변환이다. 따라서 $Y$는
매끄럽고 코호몰로지와 종수에 관한 결과는 보조정리
\ref{curves-lemma-genus-base-change}에서 따른다.
\end{proof}

\begin{example}
\label{curves-example-inseparable}
이 예는 고유 비특이 곡선들의 순수 비분리 사상 아래에서 종수가 변할 수 있음을
보여 준다. $k$를 표수 $3$인 체라 하고, $3$제곱이 아닌 원소 $a\in k$가
존재한다고 가정하자. 예를 들어 $k=\mathbf F_3(a)$이면 된다. $X$를 절
\ref{curves-section-plane-curves}에서와 같이 동차방정식
$$
F=T_1^2T_0-T_2^3+aT_0^3
$$
으로 정의된 평면곡선이라 하자. 아핀 조각 $D_+(T_0)$에서 좌표
$x=T_1/T_0$, $y=T_2/T_0$를 쓰면 $x^2-y^3+a=0$을 얻고, 이는 비특이
아핀곡선을 정의한다. 또한 무한원점 $(0:1:0)$은 매끄러운 점이다. 따라서
$X$는 종수 $1$인 비특이 사영곡선이다(보조정리
\ref{curves-lemma-genus-plane-curve}). 한편 $D_+(T_0)$에서 $(x,y)$를
$x\in\mathbf A^1_k\subset\mathbf{P}^1_k$로 보내는 사상
$f:X\to\mathbf{P}^1_k$를 생각하자. 그러면 $f$는 $k$ 위의 고유 비특이 곡선들의
사상이고 차수 $p=3$인 비분리 함수체 확장을 유도하지만, $X$의 종수는 $1$이고
$\mathbf{P}^1_k$의 종수는 $0$이다.
\end{example}

\begin{proposition}
\label{curves-proposition-unwind-morphism-smooth}
$k$를 표수 $p>0$인 체라 하고 $f:X\to Y$를 $k$ 위의 고유 매끄러운 곡선들의
상수가 아닌 사상이라 하자. 그러면 $f$를
$$
X\longrightarrow X^{(p^n)}\longrightarrow Y
$$
로 분해할 수 있다. 여기서 $X^{(p^n)}\to Y$는 분리 체확장
$k(X^{(p^n)})/k(Y)$를 유도하는 고유 매끄러운 곡선들의 상수가 아닌 사상이고,
$$
X^{(p^n)}=X\times_{\Spec(k),F_{\Spec(k)}^n}\Spec(k),
$$
이며 $X\to X^{(p^n)}$은 $X$의 $n$중 상대 Frobenius이다.
\end{proposition}

\begin{proof}
체, 보조정리 \ref{fields-lemma-separable-first}에 의해 중간확장
$k(X)/E/k(Y)$가 존재하여 $k(X)/E$는 순수 비분리이고 $E/k(Y)$는 분리이다.
정리 \ref{curves-theorem-curves-rational-maps}에 의해 이는 $f$의 분해
$X\to Z\to Y$에 대응하며 $Z$는 비특이 고유 곡선이다. 사상 $X\to Z$에
보조정리 \ref{curves-lemma-purely-inseparable-smooth}을 적용하면 된다.
\end{proof}

\begin{lemma}
\label{curves-lemma-inseparable-linear-system}
$k$를 표수 $p>0$인 체, $X$를 $k$ 위의 매끄러운 고유 곡선이라 하자.
$(\mathcal{L},V)$를 $r\geq1$인 $\mathfrak g^r_d$라 하자. 그러면 다음 둘 중
하나가 성립한다.
\begin{enumerate}
\item 이에 대응하는 사상 $X\to\mathbf{P}^1_k$(보조정리
\ref{curves-lemma-linear-series})가 일반적으로 \'etale인(즉 보조정리
\ref{curves-lemma-generically-etale}에서와 같은) $\mathfrak g^1_d$가 존재한다.
\item $d'\leq d/p$인 $X^{(p)}$ 위의 $\mathfrak g^r_{d'}$가 존재한다.
\end{enumerate}
\end{lemma}

\begin{proof}
$k$-선형독립인 두 원소 $s,t\in V$를 택하자. 그러면 $f=s/t$는 선형계
$(\mathcal{L},ks+kt)$에 대응하는 사상 $X\to\mathbf{P}^1_k$를 정의하는
유리함수이다. 이 사상이 일반적으로 \'etale이 아니면 명제
\ref{curves-proposition-unwind-morphism-smooth}에 의해 $f\in k(X^{(p)})$이다.
이제 $V$의 기저 $s_0,\ldots,s_r$를 택하고 $\mathcal{L}'\subset\mathcal{L}$을
$s_0,\ldots,s_r$로 생성되는 가역층이라 하자. $k(X)$에서
$f_i=s_i/s_0$로 두자. 각 쌍 $(s_0,s_i)$에 대하여
$f_i\in k(X^{(p)})$이면 사상
% P07-ERR-0318
$$
\varphi=\varphi_{(\mathcal{L}',(s_0,\ldots,s_r)}:
X\longrightarrow\mathbf{P}^r_k=\text{Proj}(k[T_0,\ldots,T_r])
$$
은 $X^{(p)}$을 통하여 분해된다. 실제로 이는 아핀 열린집합 $D_+(T_0)$ 위에서
참이고, 아핀 부분 위의 사상을 보조정리 \ref{curves-lemma-extend-over-normal-curve}에 의해
매끄러운 곡선 $X^{(p)}$ 전체로 연장할 수 있다. 표기를 도입하여 $\varphi$의
분해를
$$
X\xrightarrow{F_{X/k}}X^{(p)}\xrightarrow{\psi}\mathbf{P}^r_k
$$
라 하자. 그러면 % P07-ERR-0319
$\mathcal N=\psi^*\mathcal{O}_{\mathbf{P}^1_k}(1)$은 가역
$\mathcal{O}_{X^{(p)}}$-가군이고 $\mathcal{L}'=F_{X/k}^*\mathcal N$이며,
$\psi^*T_0,\ldots,\psi^*T_r$은 $k$-선형독립이다($X$로 당기면
$s_0,\ldots,s_r$이 되기 때문이다). 마지막으로
$$
d=\deg(\mathcal{L})\geq\deg(\mathcal{L}')
=\deg(F_{X/k})\deg(\mathcal N)=p\deg(\mathcal N)
$$
이어서 원하는 결과를 얻는다. 여기서 다양체, 보조정리
\ref{varieties-lemma-check-invertible-sheaf-trivial},
\ref{varieties-lemma-degree-pullback-map-proper-curves},
\ref{varieties-lemma-inseparable-deg-p-smooth}를 사용하였다.
\end{proof}

\begin{lemma}
\label{curves-lemma-point-over-separable-extension}
$k$를 체라 하고 $X$를 $H^0(X,\mathcal{O}_X)=k$이고 종수가 $g\geq2$인
$k$ 위의 매끄러운 고유 곡선이라 하자. 그러면 $\kappa(x)/k$가 차수
$\leq2g-2$인 분리확장이 되는 닫힌점 $x\in X$가 존재한다.
\end{lemma}

\begin{proof}
$\omega=\Omega_{X/k}$라 두자. 보조정리 \ref{curves-lemma-genus-smooth}에 의해 이는
차수 $2g-2$이고 $g$개의 전역단면을 가진다. 따라서
$\mathfrak g^{g-1}_{2g-2}$가 있다. 자명한 보조정리
\ref{curves-lemma-linear-series-trivial-existence}에 의해 $\mathfrak g^1_{2g-2}$가
존재하고, 보조정리 \ref{curves-lemma-g1d}에 의해 차수 $d\leq2g-2$인 사상
$$
\varphi:X\longrightarrow\mathbf{P}^1_k
$$
를 얻는다. $\varphi$는 평탄하고(보조정리 \ref{curves-lemma-flat}) 유한하므로(보조정리
\ref{curves-lemma-finite}) 차수 $d$인 유한 국소 자유 사상이다(스킴의 사상, 보조정리
\ref{morphisms-lemma-finite-flat}). 아무 유리점 $t\in\mathbf{P}^1_k$나 택하고
$\varphi(x)=t$인 점 $x\in X$를 택하자. 그러면 예를 들어 스킴의 사상, 보조정리
\ref{morphisms-lemma-finite-locally-free-universally-bounded}와
\ref{morphisms-lemma-characterize-universally-bounded}에 의해
$$
d\geq[\kappa(x):\kappa(t)]=[\kappa(x):k]
$$
이다. 따라서 $k$가 완전체이면(예를 들어 표수가 영이거나 유한체이면)
보조정리가 증명된다. 그러므로 다음 문단에서 논할 경우만 남는다.

\medskip\noindent
$k$가 표수 $p>0$인 무한체라고 가정하자. 위와 같이 $X$가
$\mathfrak g^{g-1}_{2g-2}$를 가짐을 사용하겠다. 매끄러운 고유 곡선
$X^{(p)}$은 $X$와 같은 종수를 가진다. 따라서 그 종수는 $>0$이다. 보조정리
\ref{curves-lemma-grd-inequalities}에 의해 $X^{(p)}$은 어떤 $d\leq g-1$에 대해서도
$\mathfrak g^{g-1}_d$를 갖지 않는다. 우리의 $\mathfrak g^{g-1}_{2g-2}$에
보조정리 \ref{curves-lemma-inseparable-linear-system}을 적용하고
% P07-ERR-0320
$2g-2/p\leq g-1$임을 주목하면 가능성 (2)는 일어나지 않는다. 따라서
일반적으로 \'etale이고(보조정리의 의미에서) 차수가 $\leq2g-2$인 사상
$$
\varphi:X\longrightarrow\mathbf{P}^1_k
$$
를 얻는다. $U\subset X$를 $\varphi$가 \'etale인 공집합이 아닌 열린
부분스킴이라 하자. 그러면 $\varphi(U)\subset\mathbf{P}^1_k$는 공집합이 아닌
Zariski 열린집합이고, $k$가 무한체이므로 $k$-유리점
$t\in\varphi(U)$를 택할 수 있다. $\varphi(u)=t$인 점 $u\in U$를 택하자.
그러면 $\kappa(u)/\kappa(t)$는 분리이고(스킴의 사상, 보조정리
\ref{morphisms-lemma-etale-over-field}), $\kappa(t)=k$이며 이전과 같이
$[\kappa(u):k]\leq2g-2$이다.
\end{proof}

\noindent
다음 보조정리는 사실 이 절에 속하지 않지만 달리 둘 만한 좋은 곳을 모르겠다.

\begin{lemma}
\label{curves-lemma-ramification-to-algebraic-closure}
$X$를 체 $k$ 위의 매끄러운 곡선이라 하자. $\overline{x}\in X_{\overline{k}}$를
상이 $x\in X$인 닫힌점이라 하자. 그러면
$\mathcal{O}_{X,x}\subset\mathcal{O}_{X_{\overline{k}},\overline{x}}$의 분기지수는
$\kappa(x)/k$의 비분리 차수이다.
\end{lemma}

\begin{proof}
$X$를 줄여서 \'etale 사상 $\pi:X\to\mathbf A^1_k$가 존재한다고 가정해도
된다. 스킴의 사상, 보조정리
\ref{morphisms-lemma-smooth-etale-over-affine-space}를 보라. 그러면 국소환들의
그림
$$
\xymatrix{
\mathcal{O}_{X_{\overline{k}},\overline{x}}&
\mathcal{O}_{\mathbf A^1_{\overline{k}},\pi(\overline{x})}\ar[l]\\
\mathcal{O}_{X,x}\ar[u]&
\mathcal{O}_{\mathbf A^1_k,\pi(x)}\ar[l]\ar[u]
}
$$
을 생각할 수 있다. 수평 화살표들은 \'etale 사상에 대응하므로 분기지수가
$1$이다. 또한 확장 $\kappa(x)/\kappa(\pi(x))$은 분리이므로 $\kappa(x)$와
$\kappa(\pi(x))$는 $k$ 위에서 같은 비분리 차수를 가진다. 분기지수의
곱셈성에 의해 $x$가 아핀직선의 점인 경우만 증명하면 충분하다.

\medskip\noindent
$X=\mathbf A^1_k$라고 가정하자. 이 경우 $x$에서 $X$의 국소환은
$$
\mathcal{O}_{X,x}=k[t]_{(P)}
$$
의 꼴이며 $P$는 $k$ 위의 기약 일계수 다항식이다. 체, 보조정리
\ref{fields-lemma-irreducible-polynomials}에 의해 어떤 분리다항식
$Q\in k[t]$에 대하여 $P(t)=Q(t^q)$이다. $\kappa(x)=k[t]/(P)$가
$k$ 위에서 비분리 차수 $q$임을 주목하자. 한편 $\overline{k}$ 위에서
$Q(t)=\prod(t-\alpha_i)$로 인수분해할 수 있고 $\alpha_i$들은 서로 다르다.
유일한 $\beta_i\in\overline{k}$에 대하여 $\alpha_i=\beta_i^q$라 쓰자.
그러면 점 $\overline{x}$는 $\beta_i$들 중 하나에 대응하며, 환사상
$$
k[t]_{(P)}\longrightarrow\overline{k}[t]_{(t-\beta_i)}
$$
의 분기지수는 실제로 $q$이다. 균등화원 $P$의 상이 단위를 곱한
$(t-\beta_i)^q$이기 때문이다. 이로써 결론을 얻는다.
\end{proof}




\section{밀어붙임}
\label{curves-section-pushouts}

\noindent
$k$를 체라 하자. 다음 $k$ 위 스킴들의 실선 그림을 생각하자.
$$
\xymatrix{
Z' \ar[d] \ar[r]_{i'} & X' \ar@{..>}[d]^a \\
Z \ar@{..>}[r]^i & X
}
$$
다음을 가정한다.
\begin{enumerate}
\item[(a)] $X'$은 $k$ 위의 차원 $\leq1$인 분리 유한형 스킴이다.
\item[(b)] % P07-ERR-0321
$i:Z'\to X'$은 닫힌몰입이다.
\item[(c)] $Z'$과 $Z$는 $\Spec(k)$ 위에서 유한이다.
\item[(d)] $Z'\to Z$는 전사이다.
\end{enumerate}
이 상황에서 $X'$의 임의의 유한 점집합은 어떤 아핀 열린집합에 포함된다.
다양체, 명제
\ref{varieties-proposition-finite-set-of-points-of-codim-1-in-affine}을 보라.
따라서 스킴의 사상 심화, 명제
\ref{more-morphisms-proposition-pushout-along-closed-immersion-and-integral}의
가정들이 성립하고 다음을 얻는다. % P07-ERR-0322
\begin{enumerate}
\item 밀어붙임 $X=Z\amalg_{Z'}X'$이 스킴의 범주에서 존재한다.
\item $i:Z\to X$는 닫힌몰입이다.
\item $a:X'\to X$는 정역이고 전사이다. % P07-ERR-0323
\item 스킴의 사상 심화, 보조정리
\ref{more-morphisms-lemma-pushout-separated}에 의해 $X\to\Spec(k)$는 분리이다.
\item 스킴의 사상 심화, 보조정리
\ref{more-morphisms-lemma-pushout-finite-type}에 의해 $X\to\Spec(k)$는 유한형이다.
\item 따라서 스킴의 사상, 보조정리 \ref{morphisms-lemma-finite-integral}과
\ref{morphisms-lemma-permanence-finite-type}에 의해 $a:X'\to X$는 유한이다.
\item $X'\to\Spec(k)$가 고유이면 스킴의 사상, 보조정리
\ref{morphisms-lemma-image-proper-is-proper}에 의해 $X\to\Spec(k)$는 고유이다.
\end{enumerate}

\noindent
다음 보조정리는 상당히 일반화할 수 있다.

\begin{lemma}
\label{curves-lemma-complete-local-ring-pushout}
위의 상황에서 $Z=\Spec(k')$이고 $k'$가 체이며,
$Z'=\Spec(k'_1\times\ldots\times k'_n)$이고 $k'_i/k'$가 유한 체확장이라고
하자. $x\in X$를 $Z\to X$의 상이라 하고 $x'_i\in X'$를
$\Spec(k'_i)\to X'$의 상이라 하자. 그러면 올곱 그림
$$
\xymatrix{
\prod\nolimits_{i=1,\ldots,n}k'_i&
\prod\nolimits_{i=1,\ldots,n}\mathcal{O}_{X',x'_i}^\wedge\ar[l]\\
k'\ar[u]&\mathcal{O}_{X,x}^\wedge\ar[u]\ar[l]
}
$$
을 얻는다. 여기서 수평 화살표들은 잉여체로 가는 사상들로 주어진다.
\end{lemma}

\begin{proof}
$X$에서 $x$의 아핀 열린 근방 $\Spec(A)$를 택하자. 그 역상을
$\Spec(A')\subset X'$이라 하자. 구성상 올곱 그림
$$
\xymatrix{
\prod\nolimits_{i=1,\ldots,n}k'_i&A'\ar[l]\\
k'\ar[u]&A\ar[u]\ar[l]
}
$$
을 얻는다. 모든 것이 $A$ 위에서 유한이므로 $x$에 대응하는 극대 아이디얼
$\mathfrak m\subset A$에 관하여 완비한 뒤에도 이 그림은 올곱 그림이다
(대수학, 보조정리 \ref{algebra-lemma-completion-flat}). 마지막으로 대수학,
보조정리 \ref{algebra-lemma-completion-finite-extension}를 적용하여 $A'$의 완비를
식별하면 된다.
\end{proof}




\section{붙이기와 찌그러뜨리기}
\label{curves-section-glueing-squishing}

\noindent
아래에서 $k[\epsilon]$은 다양체, 정의
\ref{varieties-definition-dual-numbers}에 정의된 $k$ 위의 쌍대수 대수를 나타낸다.
% P07-ERR-0324

\begin{lemma}
\label{curves-lemma-no-in-between-over-k}
$k$를 대수적으로 닫힌 체라 하자. 환확장 $k\subset A$에서 $A$가 정확히 두 개의
$k$-부분대수를 가지면 $A=k\times k$이거나 $A=k[\epsilon]$이다.
\end{lemma}

\begin{proof}
가정은 $k\not = A$이고 모든 부분환 $k\subset C\subset A$가 $k$ 또는 $A$와
같다는 뜻이다. $t\in A$, $t\not \in k$를 택하자. 그러면 $A$는 $k$ 위에서
$t$로 생성된다. 따라서 어떤 아이디얼 $I$에 대하여 $A=k[x]/I$이다.
$I=(0)$이면 허용되지 않는 부분대수 $k[x^2]$가 있다. 그렇지 않으면 $I$는
어떤 일계수 다항식 $P$로 생성된다. % P07-ERR-0325
$P=\prod_{i=1}^d(t-a_i)$라 쓰자. $d>2$이면
$(t-a_1)(t-a_2)$로 생성되는 부분대수가 모순을 준다. 따라서 $d=2$이다.
$a_1\not = a_2$이면 $A=k\times k$이고, $a_1=a_2$이면 $A=k[\epsilon]$이다.
\end{proof}

\begin{example}[점 붙이기]
\label{curves-example-glue-points}
$k$를 대수적으로 닫힌 체라 하고 $f:X'\to X$를 대수적 $k$-스킴들의
사상이라 하자. 다음이 성립하면 $X$는 $X'$에서 $a$와 $b$를 붙여서 얻어진다고
한다.
\begin{enumerate}
\item $a,b\in X'(k)$는 서로 다른 점이고 같은 점 $x\in X(k)$로 간다.
\item $f$는 유한이고
$f^{-1}(X\setminus\{x\})\to X\setminus\{x\}$는 동형이다.
\item 짧은 완전열
$$
0\to\mathcal{O}_X\to f_*\mathcal{O}_{X'}\xrightarrow{a-b}x_*k\to0
$$
이 있으며, 여기서 오른쪽 화살표는 $f_*\mathcal{O}_{X'}$의 국소단면 $h$를
차 $h(a)-h(b)\in k$로 보낸다. % P07-ERR-0326
\end{enumerate}
이 경우 짧은 완전열
$$
0\to\mathcal{O}_X^*\to f_*\mathcal{O}_{X'}^*
\xrightarrow{ab^{-1}}x_*k^*\to0
$$
도 있으며, 여기서 오른쪽 화살표는 $f_*\mathcal{O}_{X'}^*$의 국소단면 $h$를
곱셈적 차 $h(a)h(b)^{-1}\in k^*$로 보낸다. % P07-ERR-0326
\end{example}

\begin{example}[접벡터 찌그러뜨리기]
\label{curves-example-squish-tangent-vector}
$k$를 대수적으로 닫힌 체라 하고 $f:X'\to X$를 대수적 $k$-스킴들의
사상이라 하자. 다음이 성립하면 $X$는 $X'$에서 접벡터 $\vartheta$를
찌그러뜨려서 얻어진다고 한다.
\begin{enumerate}
\item $\vartheta:\Spec(k[\epsilon])\to X'$은 $k$ 위의 닫힌몰입이고,
$f\circ\vartheta$는 한 점 $x\in X(k)$를 통하여 분해된다.
\item $f$는 유한이고
$f^{-1}(X\setminus\{x\})\to X\setminus\{x\}$는 동형이다.
\item 짧은 완전열
$$
0\to\mathcal{O}_X\to f_*\mathcal{O}_{X'}
\xrightarrow{\vartheta}x_*k\to0
$$
이 있으며, 여기서 오른쪽 화살표는 $f_*\mathcal{O}_{X'}$의 국소단면 $h$를
$\vartheta^\sharp(h)\in k[\epsilon]$에서 $\epsilon$의 계수로 보낸다.
% P07-ERR-0326
\end{enumerate}
이 경우 짧은 완전열
$$
0\to\mathcal{O}_X^*\to f_*\mathcal{O}_{X'}^*
\xrightarrow{\vartheta}x_*k\to0
$$
도 있으며, 여기서 오른쪽 화살표는 $f_*\mathcal{O}_{X'}^*$의 국소단면 $h$를
$\text{d}\log(\vartheta^\sharp(h))$로 보낸다. 여기서
$\text{d}\log:k[\epsilon]^*\to k$는 $a+b\epsilon$을 $b/a\in k$로 보내는
가환군의 준동형이다. % P07-ERR-0326
\end{example}

\begin{lemma}
\label{curves-lemma-factor-almost-isomorphism}
$k$를 대수적으로 닫힌 체라 하자. $f:X'\to X$를
$\mathcal{O}_X\subset f_*\mathcal{O}_{X'}$인 대수적 $k$-스킴들의
유한사상이라 하자. $f$는 유한개의 점을 제외하면 동형이라고 가정하자. 그러면 분해
$$
X'=X_n\to X_{n-1}\to\ldots\to X_1\to X_0=X
$$
가 존재하여 각 $X_i\to X_{i-1}$은 두 점을 붙이거나 접벡터를 찌그러뜨리는
사상이다(예 \ref{curves-example-glue-points}와
\ref{curves-example-squish-tangent-vector} 참조).
\end{lemma}

\begin{proof}
$U\subset X$를 $f$가 동형인 극대 열린집합이라 하자. 그러면
$X\setminus U=\{x_1,\ldots,x_n\}$이고 $x_i\in X(k)$이다. $f$의 분해
$X'\to Y\to X$ 중 두 사상이 모두 유한이고
$$
\mathcal{O}_X\subset g_*\mathcal{O}_Y\subset f_*\mathcal{O}_{X'}
$$
인 것들을 생각하자. 여기서 $g:Y\to X$는 주어진 사상이다. 가정에 의해
어떤 $i$에 대하여 $x=x_i$인 경우를 제외하면
$\mathcal{O}_{X,x}\to(f_*\mathcal{O}_{X'})_x$는 동형이다. % P07-ERR-0327
따라서 여핵
$$
f_*\mathcal{O}_{X'}/\mathcal{O}_X=\bigoplus\mathcal{Q}_i
$$
는 $x_1,\ldots,x_n$에 지지되는 마천루층 $\mathcal{Q}_i$들의 직합이다. 표시된 몫은 연접
$\mathcal{O}_X$-가군이므로 $\mathcal{Q}_i$는 $\mathcal{O}_{X,x_i}$ 위에서 유한
길이를 가진다. 따라서 이 길이들의 합, 즉 전체 여핵의 길이에 대한 귀납법을
쓸 수 있다.

\medskip\noindent
$n>1$이면 $\mathcal{Q}_1$의 역상을 취하여
$\mathcal{O}_X$-부분대수 $\mathcal{A}\subset f_*\mathcal{O}_{X'}$를 정의할 수
있다. 이는 자명하지 않은 분해를 주고 귀납가정으로 결론을 얻는다.

\medskip\noindent
$n=1$이라고 가정하고 $x=x_1$이라 줄여 쓰자. 유한 $k$-대수확장
$$
A=\mathcal{O}_{X,x}\subset(f_*\mathcal{O}_{X'})_x=B
$$
를 생각하자. $\mathcal{Q}=\mathcal{Q}_1$은 값이 $B/A$인 마천루층임을 주목하자.
$k$-부분대수 $A\subset A+\mathfrak m_AB\subset B$가 있다. 두 포함이 모두
진포함이면 자명하지 않은 분해를 얻고 위와 같이 귀납가정으로 결론을 얻는다.
$A+\mathfrak m_AB=B$이면 Nakayama 보조정리에 의해 $A=B$이고, 그러면 $f$는
동형이어서 증명할 것이 없다. % P07-ERR-0328
따라서 $B=A+\mathfrak m_AB$라고 가정해도 된다. $C=B/\mathfrak m_AB$라 하자.
$C$가 $2$개보다 많은 $k$-부분대수를 가지면 그 역상을 $B$에서 취하여 $A$와
$B$ 사이의 부분대수를 얻는다. 따라서 $C$가 정확히 $2$개의 $k$-부분대수를
가진다고 가정해도 된다. 보조정리 \ref{curves-lemma-no-in-between-over-k}에 의해
$C=k\times k$이거나 $C=k[\epsilon]$이다. 이에 따라 $f$는 두 점을 붙이는
사상이거나 접벡터를 찌그러뜨리는 사상이다. % P07-ERR-0329
\end{proof}

\begin{lemma}
\label{curves-lemma-glue-points}
$k$를 대수적으로 닫힌 체라 하자. $f:X'\to X$가 예
\ref{curves-example-glue-points}에서와 같이 두 점 $a,b$를 붙이는 사상이면 완전열
$$
k^*\to\Pic(X)\to\Pic(X')\to0
$$
이 있다. $a$와 $b$가 $X'$의 서로 다른 연결성분에 있으면 첫 사상은 영이고,
$X'$이 고유이며 $a,b$가 $X'$의 같은 연결성분에 있으면 첫 사상은 단사이다.
\end{lemma}

\begin{proof}
다양체, 보조정리 \ref{varieties-lemma-surjective-pic-birational-finite}에 의해
$\Pic(X)\to\Pic(X')$은 전사이다. 짧은 완전열
$$
0\to\mathcal{O}_X^*\to f_*\mathcal{O}_{X'}^*
\xrightarrow{ab^{-1}}x_*k^*\to0
$$
을 쓰면
$$
H^0(X',\mathcal{O}_{X'}^*)\xrightarrow{ab^{-1}}k^*\to
H^1(X,\mathcal{O}_X^*)\to H^1(X,f_*\mathcal{O}_{X'}^*)
$$
을 얻는다. $H^1(X,f_*\mathcal{O}_{X'}^*)\subset H^1(X',\mathcal{O}_{X'}^*)$이다
(예를 들어 Leray 스펙트럼열을 쓰라. 코호몰로지, 보조정리
\ref{cohomology-lemma-Leray} 참조). 따라서
$\Pic(X)\to\Pic(X')$의 핵은
$ab^{-1}:H^0(X',\mathcal{O}_{X'}^*)\to k^*$의 여핵이다. $a,b$가 $X'$의
서로 다른 연결성분에 있으면 $ab^{-1}$은 전사이다. $k$가 대수적으로 닫혔으므로
$k$ 위의 환원 연결 고유 스킴의 모든 정칙함수는 $k$의 원소에서 온다.
다양체, 보조정리
\ref{varieties-lemma-proper-geometrically-reduced-global-sections}를 보라.
따라서 $X'$이 고유이고 $a,b$가 같은 연결성분에 있으면 $ab^{-1}$은 영이다.
\end{proof}

\begin{lemma}
\label{curves-lemma-squish-tangent-vector}
$k$를 대수적으로 닫힌 체라 하자. $f:X'\to X$가 예
\ref{curves-example-squish-tangent-vector}에서와 같이 접벡터 $\vartheta$를
찌그러뜨리는 사상이면 완전열
$$
(k,+)\to\Pic(X)\to\Pic(X')\to0
$$
이 있고, $X'$이 고유이고 환원이면 첫 사상은 단사이다.
\end{lemma}

\begin{proof}
다양체, 보조정리 \ref{varieties-lemma-surjective-pic-birational-finite}에 의해
$\Pic(X)\to\Pic(X')$은 전사이다. 예
\ref{curves-example-squish-tangent-vector}의 짧은 완전열
$$
0\to\mathcal{O}_X^*\to f_*\mathcal{O}_{X'}^*
\xrightarrow{\vartheta}x_*k\to0
$$
을 쓰면
$$
H^0(X',\mathcal{O}_{X'}^*)\xrightarrow{\vartheta}k\to
H^1(X,\mathcal{O}_X^*)\to H^1(X,f_*\mathcal{O}_{X'}^*)
$$
을 얻는다. $H^1(X,f_*\mathcal{O}_{X'}^*)\subset H^1(X',\mathcal{O}_{X'}^*)$이다
(예를 들어 Leray 스펙트럼열을 쓰라. 코호몰로지, 보조정리
\ref{cohomology-lemma-Leray} 참조). 따라서
$\Pic(X)\to\Pic(X')$의 핵은 사상
$\vartheta:H^0(X',\mathcal{O}_{X'}^*)\to k$의 여핵이다. $k$가 대수적으로
닫혔으므로 $k$ 위의 환원 연결 고유 스킴의 모든 정칙함수는 $k$의 원소에서
온다. 다양체, 보조정리
\ref{varieties-lemma-proper-geometrically-reduced-global-sections}를 보라.
이로써 보조정리의 마지막 명제가 증명된다. % P07-ERR-0330
\end{proof}


\section{다중교차 특이점과 마디 특이점}
\label{curves-section-multicross}

\noindent
이 절에서는 가능한 가장 단순한 곡선 특이점들을 논한다.

\medskip\noindent
$k$를 체라 하자. 완비 국소 $k$-대수
\begin{equation}
\label{curves-equation-multicross}
A=\{(f_1,\ldots,f_n)\in k[[t]]\times\ldots\times k[[t]]\mid
f_1(0)=\ldots=f_n(0)\}
\end{equation}
를 생각하자. 다양체, 정의 \ref{varieties-definition-wedge}에 도입된
용어로 $A$는 $k$ 위의 변수가 $1$개인 멱급수환 $n$개의 쐐기이다.
$k[[t]]\times\ldots\times k[[t]]$가 $A$의 전분수환 안에서의 정수적 폐포임을
주목하자. 따라서 $A$의 $\delta$-불변량은 $n-1$이다. $x_i$를 $A$의
$(0,\ldots,0,t,0,\ldots,0)$으로 보내서 얻는 동형
$$
k[[x_1,\ldots,x_n]]/(\{x_ix_j\}_{i\not = j})\longrightarrow A
$$
이 있다. 따라서 $\dim(A)=1$이고 $\dim_k\mathfrak m/\mathfrak m^2=n$이다.
특히 $A$가 정칙일 필요충분조건은 $n=1$인 것이다.

\begin{lemma}
\label{curves-lemma-multicross-algebra}
$k$를 분리적으로 닫힌 체라 하자. $A$를 잉여체가 $k$인 $1$차원 환원 Nagata
국소 $k$-대수라 하자. 그러면
$$
\delta\text{-invariant }A \geq \text{number of branches of }A-1
$$
이다. 등호가 성립하면 $A^\wedge$은 (\ref{curves-equation-multicross})와 같은 꼴이다.
\end{lemma}

\begin{proof}
$A$의 잉여체가 분리적으로 닫혔으므로 $A$의 가지 수는 $A$의 기하적 가지
수와 같다. 대수학 심화, 정의
\ref{more-algebra-definition-number-of-branches}를 보라. 부등식은 다양체,
보조정리 \ref{varieties-lemma-delta-number-branches-inequality}에서 따른다.
등호가 성립한다고 하자. $A$를 $A$의 완비로 바꾸어도 된다. 이로써 가지 수나
$\delta$-불변량은 변하지 않는다. 대수학 심화, 보조정리
\ref{more-algebra-lemma-one-dimensional-number-of-branches}와 다양체, 보조정리
\ref{varieties-lemma-delta-same-after-completion}을 보라. 이제 $A$는 엄밀
Hensel이다(대수학, 보조정리 \ref{algebra-lemma-complete-henselian}). 다양체,
보조정리 \ref{varieties-lemma-delta-number-branches-inequality-sh}에 의해 $A$는
완비 이산값매김환들의 쐐기이다. 이들 각각은 대수학, 보조정리
\ref{algebra-lemma-regular-complete-containing-coefficient-field}에 의해
$k[[t]]$와 동형이다. 따라서 $A$는 (\ref{curves-equation-multicross})와 같은 꼴이다.
\end{proof}

\begin{definition}
\label{curves-definition-multicross}
$k$를 대수적으로 닫힌 체, $X$를 대수적 $1$차원 $k$-스킴,
$x\in X$를 닫힌점이라 하자. 어떤 $n\geq2$에 대하여 완비
$\mathcal{O}_{X,x}^\wedge$이 (\ref{curves-equation-multicross})와 동형이면 $x$가
{\it 다중교차 특이점}을 정의한다고 한다. $n=2$이면 $x$를 {\it 마디점},
또는 {\it 보통 이중점}이라고 하며, 또는 그 점이 {\it 마디 특이점}을 정의한다고 한다.
\end{definition}

\noindent
어떤 의미에서 이들은 대수적으로 닫힌 체 위의 곡선이 가질 수 있는 가장
단순한 종류의 특이점이다.

\begin{lemma}
\label{curves-lemma-multicross}
$k$를 대수적으로 닫힌 체, $X$를 환원 대수적 $1$차원 $k$-스킴,
$x\in X$라 하자. 다음은 동치이다.
\begin{enumerate}
\item $x$는 다중교차 특이점을 정의한다.
\item $x$에서 $X$의 $\delta$-불변량은 $x$에서 $X$의 가지 수보다 $1$ 작다.
\item 사상들의 열
$U_n\to U_{n-1}\to\ldots\to U_0=U\subset X$가 존재한다. 여기서 $U$는
$x$의 열린 근방이고 $U_n$은 비특이이며, 각 $U_i\to U_{i-1}$은 예
\ref{curves-example-glue-points}에서와 같이 두 점을 붙이는 사상이다.
\end{enumerate}
\end{lemma}

\begin{proof}
(1)과 (2)의 동치는 보조정리 \ref{curves-lemma-multicross-algebra}이다.

\medskip\noindent
(3)을 가정하자. $i$에 대한 내림귀납법으로 $U_i$의 모든 특이점이 다중교차임을
보이겠다. $U_n$에는 특이점이 없으므로 이는 $U_n$에 대하여 참이다.
$U_i$가 $U_{i+1}$의 $a,b\in U_{i+1}$를 한 점 $c\in U_i$로 붙여서 얻어진다면
$$
\mathcal{O}_{U_i,c}^\wedge\subset
\mathcal{O}_{U_{i+1},a}^\wedge\times\mathcal{O}_{U_{i+1},b}^\wedge
$$
은 $k$ 안에서 같은 잉여류를 갖는 원소들의 집합이다. 따라서 $c$에서의 가지
수는 $a,b$에서의 가지 수의 합이고, $c$에서의 $\delta$-불변량은 $a,b$에서의
$\delta$-불변량의 합보다 $1$ 크다(표시된 포함의 여차원이 $1$이기 때문이다).
따라서 원하는 대로 (2)가 성립한다.

\medskip\noindent
동치인 조건 (1), (2)를 가정하자. $x$가 $U$의 유일한 특이점이 되도록 열린
$U\subset X$를 택할 수 있다. 정규화 사상에 보조정리
\ref{curves-lemma-factor-almost-isomorphism}을 적용하여
$$
U^\nu=U_n\to U_{n-1}\to\ldots\to U_1\to U_0=U
$$
를 얻는다. 어느 단계에서도 접벡터를 찌그러뜨리지 않음을 보이면 된다.
$U_{i+1}\to U_i$가 $u\in U_i$ 위에 놓인 $u'\in U_{i+1}$에서 접벡터
$\theta$를 찌그러뜨리는, $i$가 최소인 단계라고 하자. 위와 같이 논증하면
$u_i$는 다중교차 특이점이다($U_i\to\ldots\to U_0$의 사상들이 점들의 쌍을
붙이기 때문이다). % P07-ERR-0331 P07-ERR-0332
그러나 이제 $u'$과 $u$에서 가지 수는 같고, $u$에서 $U_i$의
$\delta$-불변량은 $u'$에서 $U_{i+1}$의 $\delta$-불변량보다 $1$ 크다.
보조정리 \ref{curves-lemma-multicross-algebra}에 의해 이는 $u$가 다중교차 특이점일 수
없음을 뜻하여 모순이다.
\end{proof}

\begin{lemma}
\label{curves-lemma-multicross-gorenstein-is-nodal}
$k$를 대수적으로 닫힌 체, $X$를 환원 대수적 $1$차원 $k$-스킴,
$x\in X$를 다중교차 특이점(정의 \ref{curves-definition-multicross})이라 하자.
$X$가 Gorenstein이면 $x$는 마디점이다.
\end{lemma}

\begin{proof}
사상 $\mathcal{O}_{X,x}\to\mathcal{O}_{X,x}^\wedge$은 평탄하고,
$\kappa(x)=\mathcal{O}_{X,x}^\wedge/\mathfrak m_x\mathcal{O}_{X,x}^\wedge$이라는
의미에서 비분기이다(대수학 심화, 절
\ref{more-algebra-section-permanence-completion} 참조). 따라서 $X$가
Gorenstein이면 $\mathcal{O}_{X,x}$가 Gorenstein이고, 쌍대화 복합체,
보조정리 \ref{dualizing-lemma-flat-under-gorenstein}에 의해
$\mathcal{O}_{X,x}^\wedge$도 Gorenstein이다. % P07-ERR-0333
따라서 (\ref{curves-equation-multicross})의 환 $A$가 $n\geq2$일 때 Gorenstein일
필요충분조건이 $n=2$임을 보이면 충분하다.

\medskip\noindent
$n=2$이면 $A=k[[x,y]]/(xy)$는 완전교차이므로 Gorenstein이다. 예를 들어
$k[[x,y]]\to A$에 스킴의 쌍대성, 보조정리
\ref{duality-lemma-gorenstein-lci}를 적용하고, 정칙 국소환 $k[[x,y]]$가
쌍대화 복합체, 보조정리 \ref{dualizing-lemma-regular-gorenstein}에 의해
Gorenstein이라는 사실을 쓰면 된다.

\medskip\noindent
$n>2$라고 하자. $A$가 Gorenstein이라면 $A$는 $A$ 위의 쌍대화 복합체일
것이다(스킴의 쌍대성, 정의
\ref{duality-definition-gorenstein}). % P07-ERR-0334
그러면 어떤 $n\in\mathbf{Z}$에 대하여 $R\Hom(k,A)$가 $k[n]$과 같을 것이다.
쌍대화 복합체, 보조정리 \ref{dualizing-lemma-find-function}을 보라.
% P07-ERR-0335
따라서 $n$의 값에 따라 $\Ext^1_A(k,A)\cong k$이거나
$\Ext^1_A(k,A)=0$이다(사실 $n$은 $-1$이어야 하지만 여기서는 중요하지 않다).
완전열
$$
0\to\mathfrak m_A\to A\to k\to0
$$
을 쓰면
$$
\Ext^1_A(k,A)=\Hom_A(\mathfrak m_A,A)/A
$$
를 얻는다. 여기서 $A\to\Hom_A(\mathfrak m_A,A)$는
$a\mapsto(a'\mapsto aa')$로 주어진다. $e_i\in\Hom_A(\mathfrak m_A,A)$를
$(f_1,\ldots,f_n)\in\mathfrak m_A$를
$(0,\ldots,0,f_i,0,\ldots,0)$으로 보내는 원소라 하자. 독자는
$e_1,\ldots,e_{n-1}$이 $\Hom_A(\mathfrak m_A,A)/A$에서 $k$-선형독립임을
쉽게 확인할 수 있다. 따라서
$\dim_k\Ext^1_A(k,A)\geq n-1\geq2$이고, 증명이 끝난다.
($e_1+\ldots+e_n$은 $A\to\Hom_A(\mathfrak m_A,A)$에서 $1$의 상임을 주목하라.)
\end{proof}




\section{Picard 군의 꼬임}
\label{curves-section-torsion-in-pic}

\noindent
이 절에서는 체 위의 $1$차원 고유 스킴의 Picard 군에서 꼬임을 한정한다.
이를 곡선의 반안정 축소를 연구할 때 사용할 것이다.

\medskip\noindent
보조정리 \ref{curves-lemma-torsion-picard-smooth-projective}의 결과를 얻는 초등적인
방법은 없는 듯하다. 증명을 분석하면 두 핵심 재료가 있다. (1) 매끄러운
사영곡선 위의 차수 영인 가역층들을 분류하는 Abel 다양체가 있고, (2) Abel
다양체 위의 꼬임점의 구조를 결정할 수 있다.

\begin{lemma}
\label{curves-lemma-torsion-picard-smooth-projective}
$k$를 대수적으로 닫힌 체, $X$를 $k$ 위의 종수 $g$인 매끄러운 사영곡선이라
하자.
\begin{enumerate}
\item $n\geq1$이 $k$에서 가역이면
$\Pic(X)[n]\cong(\mathbf{Z}/n\mathbf{Z})^{\oplus2g}$이다.
\item $k$의 표수가 $p>0$이면, 어떤 정수 $0\leq f\leq g$가 존재하여 모든
$m\geq1$에 대하여
$\Pic(X)[p^m]\cong(\mathbf{Z}/p^m\mathbf{Z})^{\oplus f}$이다.
\end{enumerate}
\end{lemma}

\begin{proof}
$\Pic^0(X)\subset\Pic(X)$를 차수 $0$인 가역층들의 부분군이라 하자.
달리 말하면 짧은 완전열
$$
0\to\Pic^0(X)\to\Pic(X)\xrightarrow{\deg}\mathbf{Z}\to0
$$
이 있다. 군 $\Pic^0(X)$는 군스킴 $\underline{\Picardfunctor}^0_{X/k}$의
$k$-점들의 군이다. 곡선의 Picard 스킴, 보조정리
\ref{pic-lemma-picard-pieces}를 보라. 같은 보조정리에 의해
$\underline{\Picardfunctor}^0_{X/k}$는 군체, 정의
\ref{groupoids-definition-abelian-variety}의 의미에서 $k$ 위의 $g$차원 Abel
다양체이다. 따라서 군체, 명제
\ref{groupoids-proposition-review-abelian-varieties}의 결과로 결론을 얻는다.
\end{proof}

\begin{lemma}
\label{curves-lemma-torsion-picard-becomes-visible}
$k$를 체라 하고 $n$이 $k$의 표수와 서로 소라고 하자. $X$를
$H^0(X,\mathcal{O}_X)=k$이고 종수가 $g$인 $k$ 위의 매끄러운 고유 곡선이라 하자.
\begin{enumerate}
\item $g=1$이면, $X_{k'}$가 $k'$-유리점을 가지고
$\Pic(X_{k'})[n]\cong(\mathbf{Z}/n\mathbf{Z})^{\oplus2}$가 되는 유한 분리확장
$k'/k$가 존재한다.
\item $g\geq2$이면, $X_{k'}$가 $k'$-유리점을 가지고
$\Pic(X_{k'})[n]\cong(\mathbf{Z}/n\mathbf{Z})^{\oplus2g}$가 되는 유한 분리확장
$k'/k$가 존재하며 $[k':k]\leq(2g-2)(n^{2g})!$이다.
\end{enumerate}
\end{lemma}

\begin{proof}
$g\geq2$라고 하자. 먼저 보조정리 \ref{curves-lemma-point-over-separable-extension}에
의해 차수 $2g-2$ 이하의 유한 분리확장을 택하여 $X$가 유리점을 갖게 할 수
있다. 따라서 $X$가 $k$-유리점 $x\in X(k)$를 가진다고 가정해도 되지만,
이제 $[k':k]\leq(n^{2g})!$로 보조정리를 증명해야 한다. 한 대수적 폐포 안에서
분리 대수적 폐포를 $k\subset k^{sep}\subset\overline{k}$로 택하자. 보조정리
\ref{curves-lemma-torsion-picard-smooth-projective}에 의해
$$
\Pic(X_{\overline{k}})[n]\cong(\mathbf{Z}/n\mathbf{Z})^{\oplus2g}
$$
이다. 곡선의 Picard 스킴, 보조정리 \ref{pic-lemma-torsion-descends}에 의해
$$
\Pic(X_{k^{sep}})[n]\cong(\mathbf{Z}/n\mathbf{Z})^{\oplus2g}
$$
이다. 곡선의 Picard 스킴, 보조정리
\ref{pic-lemma-torsion-descends}에 의해 연속 작용
% P07-ERR-0336
$$
\text{Gal}(k^{sep}/k)\longrightarrow
\text{Aut}(\Pic(X_{k^{sep}})[n]
$$
이 있고, 이 사상의 핵에 대한 고정체 $k'$에 대하여 보조정리가 성립한다.
작용이 연속이므로 핵은 열린 부분군이고, 따라서 $k'/k$는 유한이다. Galois
이론에 의해 $\text{Gal}(k'/k)$는 표시된 화살표의 상이다. 원소 수가
$n^{2g}$인 집합의 치환군의 원소 수가 $(n^{2g})!$이므로 Galois 이론에 의해
$[k':k]\leq(n^{2g})!$이다. (물론 이는 상계
$|\text{GL}_{2g}(\mathbf{Z}/n\mathbf{Z})|$로도 보조정리를 증명하지만, 여기서는
어떤 상계가 존재한다는 것만 필요하다.)

\medskip\noindent
종수가 $1$이면 $X$가 유리점을 갖게 되는 유한 분리 체확장의 차수에는 상계가
없다(세부사항은 생략한다). 그래도 그러한 확장은 존재한다. 예를 들어
다양체, 보조정리 \ref{varieties-lemma-smooth-separable-closed-points-dense}를
보라. 증명의 나머지는 $g\geq2$인 경우와 같다.
\end{proof}

\begin{proposition}
\label{curves-proposition-torsion-picard-reduced-proper}
$k$를 대수적으로 닫힌 체라 하고 $X$를 환원이고 연결이며 차원 $1$인 $k$ 위의
고유 스킴이라 하자. $g$를 $X$의 종수라 하고 $g_{geom}$을 $X$의 기약성분들의
기하종수의 합이라 하자. $k$의 표수와 다른 임의의 소수 $\ell$에 대하여
$$
\dim_{\mathbf F_\ell}\Pic(X)[\ell]\leq g+g_{geom}
$$
이고, 등호가 성립할 필요충분조건은 $X$의 모든 특이점이 다중교차인 것이다.
\end{proposition}

\begin{proof}
$\nu:X^\nu\to X$를 정규화라 하자(다양체, 보조정리
\ref{varieties-lemma-prepare-delta-invariant}). 보조정리
\ref{curves-lemma-factor-almost-isomorphism}에서와 같은 분해
$$
X^\nu=X_n\to X_{n-1}\to\ldots\to X_1\to X_0=X
$$
를 택하자. $h^0_i=\dim_kH^0(X_i,\mathcal{O}_{X_i})$,
$h^1_i=\dim_kH^1(X_i,\mathcal{O}_{X_i})$라 쓰자. 보조정리
\ref{curves-lemma-glue-points}와 \ref{curves-lemma-squish-tangent-vector}에 의해 각
$n>i\geq0$에 대하여 다음 세 경우 중 하나가 성립한다. % P07-ERR-0337
\begin{enumerate}
\item $X_i$는 서로 다른 연결성분에 있는 $a,b\in X_{i+1}$를 붙여서 얻어진다.
이 경우 $\Pic(X_i)=\Pic(X_{i+1})$,
$h^0_{i+1}=h^0_i+1$, $h^1_{i+1}=h^1_i$이다.
\item $X_i$는 같은 연결성분에 있는 $a,b\in X_{i+1}$를 붙여서 얻어진다.
이 경우 짧은 완전열
$$
0\to k^*\to\Pic(X_i)\to\Pic(X_{i+1})\to0
$$
이 있고 $h^0_{i+1}=h^0_i$, $h^1_{i+1}=h^1_i-1$이다.
\item $X_i$는 $X_{i+1}$의 접벡터를 찌그러뜨려서 얻어진다. 이 경우 짧은 완전열
$$
0\to(k,+)\to\Pic(X_i)\to\Pic(X_{i+1})\to0
$$
이 있고 $h^0_{i+1}=h^0_i$, $h^1_{i+1}=h^1_i-1$이다.
\end{enumerate}
구조층의 코호몰로지군들의 차원에 관한 명제는 예
\ref{curves-example-glue-points}와 \ref{curves-example-squish-tangent-vector}의 완전열들을
쓰면 된다. $k$는 대수적으로 닫혔고 그 표수는 $\ell$과 서로 소이므로
$(k,+)$와 $k^*$는 $\ell$-가분이고, 그 $\ell$-꼬임은
$(k,+)[\ell]=0$, $k^*[\ell]\cong\mathbf F_\ell$이다. 따라서
$$
\dim_{\mathbf F_\ell}\Pic(X_{i+1})[\ell]-
\dim_{\mathbf F_\ell}\Pic(X_i)[\ell]
$$
은 경우 (2)에서 $-1$인 것을 제외하면 영이다. 이 과정의 끝에는 매끄러운
사영 $k$-곡선들의 분리합인 정규화 $X^\nu=X_n$을 얻는다. 따라서
\begin{enumerate}
\item $h^1_n=g_{geom}$이고,
\item $\dim_{\mathbf F_\ell}\Pic(X_n)[\ell]=2g_{geom}$이다.
\end{enumerate}
마지막 등식은 보조정리 \ref{curves-lemma-torsion-picard-smooth-projective}에서 따른다.
% P07-ERR-0338
$g=h^1_0$이므로 (2), (3)형 단계의 수는
$h^1_0-h^1_n=g-g_{geom}$ 이하이다. % P07-ERR-0340
위에서 계산한 계수 차이로부터 % P07-ERR-0339
$$
\dim_{\mathbf F_\ell}\Pic(X)[\ell]\leq
g-g_{geom}+2g_{geom}=g+g_{geom}
$$
이고, 등호가 성립할 필요충분조건은 (3)형 단계가 하나도 없는 것임을 얻는다.
보조정리 \ref{curves-lemma-multicross}에 의해 이는 정확히 $X$의 모든 특이점이
다중교차라는 뜻이다. 반대로 모든 특이점이 다중교차이면 보조정리
\ref{curves-lemma-multicross}는 위와 같은 열
$X^\nu=X_n\to\ldots\to X_0=X$을 (3)형 단계 없이 찾을 수 있음을 보장하고,
따라서 보조정리에서 등호가 성립한다. % P07-ERR-0341
($x$의 모든 특이점에 작동하는 하나의 열을 얻으려면 각 점의 국소 열들을
붙이면 된다. 일부 세부사항은 생략한다.) % P07-ERR-0342
\end{proof}





\section{종수와 기하종수}
\label{curves-section-genus-geometric-genus}

\noindent
$k$를 대수적 폐포 $\overline{k}$를 갖는 체라 하고 $X$를 $k$ 위의 차원
$\leq1$인 고유 스킴이라 하자. $g_{geom}(X/k)$를 $X_{\overline{k}}$의 차원
$1$인 기약성분들의 기하종수의 합으로 정의한다.

\begin{lemma}
\label{curves-lemma-bound-geometric-genus}
$k$를 체라 하고 $X$를 $k$ 위의 차원 $\leq1$인 고유 스킴이라 하자. 그러면
$$
g_{geom}(X/k)=\sum\nolimits_{C\subset X}g_{geom}(C/k)
$$
이다. 여기서 합은 차원 $1$인 기약성분 $C\subset X$들을 돈다.
\end{lemma}

\begin{proof}
정의와 다음 사실에서 바로 따른다. $X_{\overline{k}}$의 기약성분
$\overline{Z}$는 $X$의 어떤 기약성분 $Z$ 위로 전사하고(다양체, 보조정리
\ref{varieties-lemma-image-irreducible}), $\overline{Z}$의 일반점에 스킴의 사상,
보조정리 \ref{morphisms-lemma-dimension-fibre-after-base-change}를 적용하면 두
성분의 차원이 같음을 알 수 있다.
\end{proof}

\begin{lemma}
\label{curves-lemma-geometric-genus-normalization}
$k$를 체라 하고 $X$를 $k$ 위의 차원 $\leq1$인 고유 스킴이라 하자. 그러면
\begin{enumerate}
\item $g_{geom}(X/k)=g_{geom}(X_{red}/k)$이다.
\item $X'\to X$가 쌍유리 고유 사상이면
$g_{geom}(X'/k)=g_{geom}(X/k)$이다.
\item $X^\nu\to X$가 정규화 사상이면
$g_{geom}(X^\nu/k)=g_{geom}(X/k)$이다.
\end{enumerate}
\end{lemma}

\begin{proof}
(1)은 보조정리 \ref{curves-lemma-bound-geometric-genus}에서 바로 따른다.
$X'\to X$가 고유 쌍유리이면 이는 유한이고 조밀한 열린집합 위에서 동형이다
(다양체, 보조정리 \ref{varieties-lemma-finite-in-codim-1}과
\ref{varieties-lemma-modification-normal-iso-over-codimension-1} 참조).
따라서 $X'_{\overline{k}}\to X_{\overline{k}}$는 조밀한 열린집합 위에서 동형이다.
그러므로 $X'_{\overline{k}}$와 $X_{\overline{k}}$의 기약성분들은 서로 일대일
대응하고 대응하는 성분들은 동형인 함수체를 가진다. 특히 이 성분들은 동형인
비특이 사영모형을 가지며 따라서 기하종수도 같다. 이로써 (2)가 증명된다.
(3)은 (1), (2)와 $X^\nu\to X_{red}$가 쌍유리라는 사실(스킴의 사상, 보조정리
\ref{morphisms-lemma-normalization-birational})에서 따른다.
\end{proof}

\begin{lemma}
\label{curves-lemma-genus-goes-down}
$k$를 체라 하고 $X$를 $k$ 위의 차원 $\leq1$인 고유 스킴이라 하자.
$f:Y\to X$를 유한사상이라 하고, 어떤 조밀한 열린집합 $U\subset X$ 위에서
$f$가 닫힌몰입이라고 하자. 그러면
$$
\dim_kH^1(X,\mathcal{O}_X)\geq\dim_kH^1(Y,\mathcal{O}_Y)
$$
이다.
\end{lemma}

\begin{proof}
$X$ 위의 연접층들의 완전열
$$
0\to\mathcal G\to\mathcal{O}_X\to f_*\mathcal{O}_Y\to\mathcal{F}\to0
$$
을 생각하자. 가정에 의해 $\mathcal{F}$는 유한개의 닫힌점에 지지되고, 따라서
고차 코호몰로지는 소멸한다(다양체, 보조정리
\ref{varieties-lemma-chi-tensor-finite}). 한편 코호몰로지, 명제
\ref{cohomology-proposition-vanishing-Noetherian}에 의해
$H^2(X,\mathcal G)=0$이다. 따라서 유도된 사상
$H^1(X,\mathcal{O}_X)\to H^1(X,f_*\mathcal{O}_Y)$는 전사이다. 스킴의 코호몰로지, 보조정리 \ref{coherent-lemma-relative-affine-cohomology}에 의해
$H^1(X,f_*\mathcal{O}_Y)=H^1(Y,\mathcal{O}_Y)$이므로 결론이 따른다.
\end{proof}

\begin{lemma}
\label{curves-lemma-genus-normalization}
$k$를 체라 하고 $X$를 $k$ 위의 차원 $\leq1$인 고유 스킴이라 하자.
$X'\to X$가 쌍유리 고유 사상이면
$$
\dim_kH^1(X,\mathcal{O}_X)\geq\dim_kH^1(X',\mathcal{O}_{X'})
$$
이다. $X$가 환원이고 $H^0(X,\mathcal{O}_X)\to H^0(X',\mathcal{O}_{X'})$가
전사이며 등호가 성립하면 $X'=X$이다.
\end{lemma}

\begin{proof}
$f:X'\to X$가 고유 쌍유리이면 이는 유한이고 조밀한 열린집합 위에서
동형이다(다양체, 보조정리 \ref{varieties-lemma-finite-in-codim-1}과
\ref{varieties-lemma-modification-normal-iso-over-codimension-1} 참조).
따라서 부등식은 보조정리 \ref{curves-lemma-genus-goes-down}에서 따른다.
% P07-ERR-0343
$X$가 환원이라고 하자. 그러면 $\mathcal{O}_X\to f_*\mathcal{O}_{X'}$은
단사이고 짧은 완전열
$$
0\to\mathcal{O}_X\to f_*\mathcal{O}_{X'}\to\mathcal{F}\to0
$$
을 얻는다. 둘째 명제의 가정 아래에서 긴 완전 코호몰로지 열로부터
$H^0(X,\mathcal{F})=0$을 얻는다. $\mathcal{F}$는 전역 생성되므로(다양체,
보조정리 \ref{varieties-lemma-chi-tensor-finite}) $\mathcal{F}=0$이고
$\mathcal{O}_X=f_*\mathcal{O}_{X'}$이다. % P07-ERR-0344
$f$가 아핀이므로 이는 $X=X'$임을 함의한다.
\end{proof}

\begin{lemma}
\label{curves-lemma-bound-geometric-genus-curve}
$k$를 체, $C$를 $k$ 위의 고유 곡선이라 하자.
$\kappa=H^0(C,\mathcal{O}_C)$라 두자. 그러면
$$
[\kappa:k]_s\dim_\kappa H^1(C,\mathcal{O}_C)\geq g_{geom}(C/k)
$$
이다.
\end{lemma}

\begin{proof}
다양체, 보조정리 \ref{varieties-lemma-regular-functions-proper-variety}에 의해
$\kappa$는 체이고 $k$의 유한확장이다. 체, 보조정리
\ref{fields-lemma-separable-degree}에 의해
$[\kappa:k]_s=|\Mor_k(\kappa,\overline{k})|$이고, 따라서
$\Spec(\kappa\otimes_k\overline{k})$는 각각 잉여체가 $\overline{k}$인
$[\kappa:k]_s$개의 점을 가진다. 따라서 집합론적으로
$$
C_{\overline{k}}=\bigcup\nolimits_{t\in\Spec(\kappa\otimes_k\overline{k})}C_t
$$
이다. 여기서 $t$를 $k$-대수사상 $t:\kappa\to\overline{k}$로 보고
$C_t=C\times_{\Spec(\kappa),t}\Spec(\overline{k})$로 둔다. 따라서
$g_{geom}(C/k)=\sum_tg_{geom}(C_t/\overline{k})$이고 합의 첨자집합의 크기는
$[\kappa:k]_s$이다. 코호몰로지와 밑변환(스킴의 코호몰로지, 보조정리
\ref{coherent-lemma-flat-base-change-cohomology})에 의해
$$
H^0(C_t,\mathcal{O}_{C_t})=\overline{k}
\quad\text{and}\quad
\dim_{\overline{k}}H^1(C_t,\mathcal{O}_{C_t})=
\dim_\kappa H^1(C,\mathcal{O}_C)
$$
이다. 정규화 $C_t^\nu$는 $C_t$의 기약성분들의 비특이 사영모형들의 분리합이다
(스킴의 사상, 보조정리
\ref{morphisms-lemma-normalization-in-terms-of-components}). 따라서
$\dim_{\overline{k}}H^1(C_t^\nu,\mathcal{O}_{C_t^\nu})$는
$g_{geom}(C_t/\overline{k})$와 같다. 보조정리 \ref{curves-lemma-genus-goes-down}에 의해
$$
\dim_{\overline{k}}H^1(C_t,\mathcal{O}_{C_t})\geq
\dim_{\overline{k}}H^1(C_t^\nu,\mathcal{O}_{C_t^\nu})
$$
이므로 증명이 끝난다.
\end{proof}

\begin{lemma}
\label{curves-lemma-bound-torsion-simple}
$k$를 체라 하고 $X$를 $k$ 위의 차원 $\leq1$인 고유 스킴이라 하자.
$\ell$을 $k$에서 가역인 소수라 하자. 그러면
$$
\dim_{\mathbf F_\ell}\Pic(X)[\ell]\leq
\dim_kH^1(X,\mathcal{O}_X)+g_{geom}(X/k)
$$
이다. 여기서 $g_{geom}(X/k)$는 위에서 정의한 것과 같다.
\end{lemma}

\begin{proof}
다양체, 보조정리 \ref{varieties-lemma-change-fields-pic}에 의해
$\Pic(X)\to\Pic(X_{\overline{k}})$는 단사이다. 스킴의 코호몰로지, 보조정리
\ref{coherent-lemma-flat-base-change-cohomology}에 의해
$\dim_kH^1(X,\mathcal{O}_X)$는
$\dim_{\overline{k}}H^1(X_{\overline{k}},\mathcal{O}_{X_{\overline{k}}})$와 같다.
따라서 $k$가 대수적으로 닫혔다고 가정해도 된다.

\medskip\noindent
$X_{red}$를 $X$의 환원이라 하자. 전사
$\mathcal{O}_X\to\mathcal{O}_{X_{red}}$는 전사
$H^1(X,\mathcal{O}_X)\to H^1(X,\mathcal{O}_{X_{red}})$를 유도한다. 준연접층의
코호몰로지가 차수 $\geq2$에서 소멸하기 때문이다(코호몰로지, 명제
\ref{cohomology-proposition-vanishing-Noetherian}).
$X_{red}\to X$는 $\overline{k}$ 위의 기약성분들에 동형을 유도하고 Picard 군의
$\ell$-꼬임에도 동형을 유도하므로(곡선의 Picard 스킴, 보조정리
\ref{pic-lemma-torsion-descends}) $X$를 $X_{red}$로 바꾸어도 된다. 이로써
명제 \ref{curves-proposition-torsion-picard-reduced-proper}로 환원된다.
\end{proof}





\section{마디곡선}
\label{curves-section-nodal}

\noindent
대수적으로 닫힌 체 위의 보통 이중점은 이미 정의하였다. 정의
\ref{curves-definition-multicross}를 보라. 즉 $x\in X$가 대수적으로 닫힌 체 $k$ 위의
$1$차원 대수적 스킴의 닫힌점이면, $x$가 보통 이중점일 필요충분조건은
$$
\mathcal{O}_{X,x}^\wedge\cong k[[x,y]]/(xy)
$$
인 것이다. 절 \ref{curves-section-multicross}에서
(\ref{curves-equation-multicross}) 뒤의 논의를 보라.

\begin{definition}
\label{curves-definition-nodal}
$k$를 체, $X$를 $1$차원 국소 대수적 $k$-스킴이라 하자.
\begin{enumerate}
\item $x$로 가는 보통 이중점 $\overline{x}\in X_{\overline{k}}$가 존재하면
닫힌점 $x\in X$를 {\it 마디점}, 또는 {\it 보통 이중점}이라고 하며, 또는
그 점이 {\it 마디 특이점}을 정의한다고 한다.
\item $X$의 모든 닫힌점이 구조사상 $X\to\Spec(k)$의 매끄러운 자리에 있거나
보통 이중점이면, $X$의 {\it 특이점들이 기껏해야 마디형}이라고 한다.
\end{enumerate}
\end{definition}

\noindent
흔히 $X$의 특이점들이 기껏해야 마디형인 $1$차원 대수적 스킴 $X$를
{\it 마디곡선}이라고 한다. 때로는 마디곡선이 고유일 것도 요구한다. 이렇게
정의한 마디곡선은 기약일 필요가 없으므로, 이는 곡선을 차원 $1$인 다양체로
정의한 앞의 정의와 충돌한다.

\begin{lemma}
\label{curves-lemma-reduced-quotient-regular-ring-dim-2}
$(A,\mathfrak m)$을 차원 $2$인 정칙 국소환이라 하고
$I\subset\mathfrak m$을 아이디얼이라 하자.
\begin{enumerate}
\item $A/I$가 환원이면 $I=(0)$이거나 $I=\mathfrak m$이거나, 어떤 영이 아닌
$f\in\mathfrak m$에 대하여 $I=(f)$이다.
\item $A/I$의 깊이가 $1$이면 어떤 영이 아닌 $f\in\mathfrak m$에 대하여
$I=(f)$이다.
\end{enumerate}
\end{lemma}

\begin{proof}
$I\not =0$이라고 하고 $I=(f_1,\ldots,f_r)$라 쓰자. $A$는 UFD이므로(대수학 심화, 보조정리 \ref{more-algebra-lemma-regular-local-UFD}),
$f$를 $f_1,\ldots,f_r$의 최대공약원으로 하여 $f_i=fg_i$라 쓸 수 있다.
따라서 $g_1,\ldots,g_r$의 최대공약원은 $1$이고, 이는
$g_1,\ldots,g_r$ 위에 높이 $1$인 소 아이디얼이 없다는 뜻이다. 그러면
$(g_1,\ldots,g_r)=A$여서 $I=(f)$이거나, 그렇지 않다면 $\dim(A)=2$에 의해
$V(g_1,\ldots,g_r)=\{\mathfrak m\}$, 즉
$\mathfrak m=\sqrt{(g_1,\ldots,g_r)}$이다.

\medskip\noindent
$A/I$가 환원, 즉 $I$가 근기 아이디얼이라고 하자. $f$가 단위이면 $I$가
근기이므로 $I=\mathfrak m$이다. $f\in\mathfrak m$이면 $f^n$의 $A/I$에서의
상은 영이다. 환원성에 의해 $f\in I$이고 따라서 $I=(f)$이다.

\medskip\noindent
$A/I$의 깊이가 $1$이라고 하자. 그러면 $\mathfrak m$은 $A/I$의 연관
소 아이디얼이 아니다. $I$을 법으로 한 $f$의 류는 $g_1,\ldots,g_r$에 의해
소멸하므로, 이는 $f$의 류가 $A/I$에서 영임을 함의한다. 따라서 원하는 대로
$I=(f)$이다.
\end{proof}

\noindent
$\kappa$를 체, $V$를 $\kappa$ 위의 벡터공간이라 하자. 유도된
$\kappa$-선형사상 $V^\vee\to V$가 동형이면
$q\in\text{Sym}^2_\kappa(V)$를 {\it 비퇴화}라고 한다. $V$의 어떤
$\kappa$-기저 $x_1,\ldots,x_n$에 대하여
$q=\sum_{i\leq j}a_{ij}x_ix_j$이면, 이는 행렬
$$
\left(
\begin{matrix}
2a_{11}&a_{12}&\ldots\\
a_{12}&2a_{22}&\ldots\\
\ldots&\ldots&\ldots
\end{matrix}
\right)
$$
의 행렬식이 영이 아니라는 뜻이다. 이는 $x_i$들에 대한 $q$의 편도함수들이
스킴론적으로 $0$을 잘라낸다는 조건과 동치이다.

\begin{lemma}
\label{curves-lemma-nodal-algebraic}
$k$를 체, $(A,\mathfrak m,\kappa)$를 Noether 국소 $k$-대수라 하자.
다음은 동치이다.
\begin{enumerate}
\item $\kappa/k$는 분리이고 $A$는 환원이며
$\dim_\kappa(\mathfrak m/\mathfrak m^2)=2$이고,
$\mathfrak m^2/\mathfrak m^3$에서 영으로 가는 비퇴화 원소
$q\in\text{Sym}^2_\kappa(\mathfrak m/\mathfrak m^2)$가 존재한다.
\item $\kappa/k$는 분리이고 $\text{depth}(A)=1$이며
$\dim_\kappa(\mathfrak m/\mathfrak m^2)=2$이고,
$\mathfrak m^2/\mathfrak m^3$에서 영으로 가는 비퇴화 원소
$q\in\text{Sym}^2_\kappa(\mathfrak m/\mathfrak m^2)$가 존재한다.
\item $\kappa$ 위의 비퇴화 이차형식 $ax^2+bxy+cy^2$에 대하여 $k$-대수로서
$A^\wedge\cong\kappa[[x,y]]/(ax^2+bxy+cy^2)$이고, $\kappa/k$는 분리이다.
\end{enumerate}
\end{lemma}

\begin{proof}
(3)을 가정하자. $ax^2+bxy+cy^2$는 기약이거나 서로 동반이 아닌 두 소원의
곱이므로 $A^\wedge$은 환원이다. 따라서 $A\subset A^\wedge$도 환원이다.
그러므로 (1)이 성립한다.

\medskip\noindent
(1)을 가정하자. $A$가 Artin이면 $\mathfrak m\not =(0)$이므로 환원일 수 없다.
따라서 $\dim(A)\geq1$이고, 대수학, 보조정리
\ref{algebra-lemma-criterion-reduced}에 의해 $\text{depth}(A)\geq1$이다.
한편 $\dim(A)=2$이면 $A$가 정칙인데, 이는 대수학, 보조정리
\ref{algebra-lemma-regular-graded}에 의해 $q$의 존재와 모순이다. 따라서
$\dim(A)\leq1$이고, 대수학, 보조정리 \ref{algebra-lemma-bound-depth}에 의해
$\text{depth}(A)=1$이다. 그러므로 (2)가 성립한다.

\medskip\noindent
(2)를 가정하자. $A$의 깊이는 $A^\wedge$의 깊이와 같고(대수학 심화,
보조정리 \ref{more-algebra-lemma-completion-depth}) 다른 조건들은 완비화에 영향을
받지 않으므로 $A$가 완비라고 가정해도 된다. 대수학 심화, 보조정리
\ref{more-algebra-lemma-lift-residue-field}에서와 같은 $\kappa\to A$를 택하자.
$\dim_\kappa(\mathfrak m/\mathfrak m^2)=2$이므로 한 기저로 가는
$x_0,y_0\in\mathfrak m$을 택할 수 있다. 규칙 $x\mapsto x_0$, $y\mapsto y_0$로
연속 $\kappa$-대수사상
$$
\kappa[[x,y]]\longrightarrow A
$$
을 얻는다. $q$를
$\text{Sym}^2_\kappa(\mathfrak m/\mathfrak m^2)$에서
$ax_0^2+bx_0y_0+cy_0^2$의 류라 하자. 두 변수 다항식으로
$Q(x,y)=ax^2+bxy+cy^2$라 쓰자. 그러면 어떤 $a_{ij}$가 $A$에 속하고
$$
Q(x_0,y_0)=ax_0^2+bx_0y_0+cy_0^2
=\sum\nolimits_{i+j=3}a_{ij}x_0^iy_0^j
$$
이다. $x_0,y_0$의 선택을 바꾸어 소멸차수를 높일 수 있음을 증명하려 한다.
$x_1,y_1\in\mathfrak m^2$라고 하자. 그러면
$$
Q(x_0+x_1,y_0+y_1)=Q(x_0,y_0)
+(2ax_0+by_0)x_1+(bx_0+2cy_0)y_1\bmod\mathfrak m^4.
$$
$Q$의 비퇴화성은 정확히 $2ax_0+by_0$와 $bx_0+2cy_0$가
$\mathfrak m/\mathfrak m^2$의 $\kappa$-기저라는 뜻이다. 보조정리 앞의 논의를
보라. 따라서
$Q(x_0+x_1,y_0+y_1)\in\mathfrak m^4$가 되도록
$x_1,y_1\in\mathfrak m^2$를 택할 수 있다. 이 과정을 귀납적으로 계속하면
$x_i,y_i\in\mathfrak m^{i+1}$를 찾아
$$
Q(x_0+x_1+\ldots+x_n,y_0+y_1+\ldots+y_n)\in\mathfrak m^{n+3}
$$
이 되게 할 수 있다. $A$가 완비이므로
$x_\infty=\sum x_i$, $y_\infty=\sum y_i$라 두고 $x$를 $x_\infty$로, $y$를
$y_\infty$로 보내는 사상 $\kappa[[x,y]]\longrightarrow A$를 생각할 수 있다.
대수학, 보조정리 \ref{algebra-lemma-completion-generalities}에 의해 이 사상은 전사
$\kappa[[x,y]]/(Q)\longrightarrow A$를 유도한다. 보조정리
\ref{curves-lemma-reduced-quotient-regular-ring-dim-2}에 의해
% P07-ERR-0345
$k[[x,y]]\to A$의 핵은 주 아이디얼이다. 그러나 핵은 $Q$의 진약수를 포함할
수 없다. 그러한 약수는 $x,y$에 관해 차수 $1$이고, 이는
$\dim(\mathfrak m/\mathfrak m^2)=2$와 모순이기 때문이다. 따라서 원하는 대로
$Q$가 핵을 생성한다.
\end{proof}

\begin{lemma}
\label{curves-lemma-2-branches-delta-1}
$k$를 체, $(A,\mathfrak m,\kappa)$를 Nagata 국소 $k$-대수라 하자. 다음은
동치이다.
\begin{enumerate}
\item $k\to A$는 보조정리 \ref{curves-lemma-nodal-algebraic}에서와 같다.
\item $\kappa/k$는 분리이고, $A$는 차원 $1$인 환원환이며,
$A$의 $\delta$-불변량은 $1$이고 $A$의 기하적 가지 수는 $2$이다.
\end{enumerate}
이들이 성립하면, $A$의 전분수환 안에서의 정수적 폐포 $A'$은 $1$개 또는 $2$개의
극대 아이디얼 $\mathfrak m'$을 가지고, 확장 $\kappa(\mathfrak m')/k$는
분리이다.
\end{lemma}

\begin{proof}
두 경우 모두 $A$와 $A^\wedge$은 환원이다. (2)의 경우에는 환원 국소 Nagata
환의 완비가 환원이기 때문이다(대수학 심화, 보조정리
\ref{more-algebra-lemma-completion-reduced}). 두 경우 모두 $A$와 $A^\wedge$의
차원은 $1$이다(대수학 심화, 보조정리
\ref{more-algebra-lemma-completion-dimension}). 다양체, 보조정리
\ref{varieties-lemma-delta-same-after-completion} 및 대수학 심화, 보조정리
\ref{more-algebra-lemma-one-dimensional-number-of-branches}에 의해 $A$와
$A^\wedge$의 $\delta$-불변량 및 기하적 가지 수는 각각 같다. $A'$을
다양체, 보조정리 \ref{varieties-lemma-pre-delta-invariant}에서와 같이 $A$의
전분수환 안에서의 정수적 폐포라 하자. 다양체, 보조정리
\ref{varieties-lemma-normalization-same-after-completion}에 의해
$A'\otimes_AA^\wedge$은 $A^\wedge$에 대하여 같은 역할을 한다. 따라서
$A$를 $A^\wedge$으로 바꾸고 $A$가 완비라고 가정해도 된다.

\medskip\noindent
(1)을 가정하자. $A$의 기하적 가지 수가 둘이고 $\delta$-불변량이 $1$임을
보이면 충분하다. 비퇴화인 $q=ax^2+bxy+cy^2$에 대하여
$A=\kappa[[x,y]]/(ax^2+bxy+cy^2)$라고 가정해도 된다. 두 경우가 있다.

\medskip\noindent
경우 I: $q$가 $\kappa$ 위에서 분해된다. 이 경우 좌표를 바꾸어 $q=xy$라고
가정해도 된다. 그러면
$$
A'=\kappa[[x,y]]/(x)\times\kappa[[x,y]]/(y)
$$
임을 알 수 있다.

\medskip\noindent
경우 II: $q$가 분해되지 않는다. 이 경우 $c\not =0$이고 비퇴화성은
$b^2-4ac\not =0$이라는 뜻이다. 따라서
$\kappa'=\kappa[t]/(a+bt+ct^2)$는 $\kappa$의 차수 $2$인 분리확장이다.
그러면 $t=y/x$는 $A$ 위에서 정수적이고
$$
A'=\kappa'[[x]]
$$
임을 얻는다. 여기서 오른쪽에서 $y$는 $tx$로 간다.

\medskip\noindent
두 경우 모두 $\delta$-불변량이 $1$이고 기하적 가지 수가 $2$임을 직접
확인할 수 있다. 이로써 (1)이 (2)를 함의한다. 또한 보조정리의 마지막
명제도 얻는다.

\medskip\noindent
(2)를 가정하자. 대수학 심화, 보조정리
\ref{more-algebra-lemma-number-of-branches-1}에 의해 $A'$은 두 극대 아이디얼을
가지거나, $A'$이 하나의 극대 아이디얼을 가지고
$[\kappa(\mathfrak m'):\kappa]_s=2$이다.

\medskip\noindent
경우 I: $A'$이 잉여체가 각각 $\kappa_1,\kappa_2$인 두 극대 아이디얼
$\mathfrak m'_1,\mathfrak m'_2$를 가진다. $\delta$-불변량은 $A'/A$의 길이이고
전사 $A'/A\to(\kappa_1\times\kappa_2)/\kappa$가 있으므로
$\kappa=\kappa_1=\kappa_2$이다. $A$는 완비이고(따라서 대수학, 보조정리
\ref{algebra-lemma-complete-henselian}에 의해 Hensel이다) $A'$은 $A$ 위에서
유한이므로 $A'=A_1\times A_2$이다(대수학, 보조정리
\ref{algebra-lemma-finite-over-henselian}). $A'$은 정규환이므로 $A_1,A_2$는
이산값매김환이다. 따라서 대수학 심화, 보조정리
\ref{more-algebra-lemma-power-series-over-residue-field}에 의해 $A_1$과 $A_2$는
$k$-대수로서 $\kappa[[t]]$와 동형이다. $\delta$-불변량이 $1$이므로 $A$는
$A_1,A_2$의 쐐기이다(다양체, 정의
\ref{varieties-definition-wedge}). 이로부터 쉽게
$A\cong\kappa[[x,y]]/(xy)$를 얻는다.

\medskip\noindent
경우 II: $A'$은 잉여체가 $\kappa'$이고
$[\kappa':\kappa]_s=2$인 하나의 극대 아이디얼 $\mathfrak m'$을 가진다.
경우 I과 똑같이 논증하면 $[\kappa':\kappa]=2$이고 $\kappa'$는 $\kappa$
위에서 분리임을 알 수 있다. $A'$은 정규이므로 $A'$은 $\kappa'[[t]]$와 동형이다
(위 참고문헌을 보라). $A'/A$의 길이가 $1$이므로
$$
A=\{f\in\kappa'[[t]]\mid f(0)\in\kappa\}
$$
이다. 간단한 계산으로 $A$가 (1)과 같은 꼴임을 알 수 있다.
\end{proof}

\begin{lemma}
\label{curves-lemma-fitting-ideal-well-defined}
$k$를 체라 하자. $A=k[[x_1,\ldots,x_n]]$이라 하고
$I=(f_1,\ldots,f_m)\subset A$를 아이디얼이라 하자. 임의의 $r\geq0$에 대하여
행렬 $(\partial f_j/\partial x_i)$의 $r\times r$ 소행렬식들로 $A/I$에서
생성되는 아이디얼은 $I$의 생성원 선택이나 $A$의 정칙 매개변수계
$x_1,\ldots,x_n$의 선택에 의존하지 않는다.
\end{lemma}

\begin{proof}
이 보조정리의 ``올바른'' 증명은 이 아이디얼이 $k$ 위의 $A/I$의 연속
미분가군의 $(n-r)$째 Fitting 아이디얼임을 보이는 것이다. 여기서는 직접
증명한다. $g_1,\ldots,g_l$이 $I$의 또 다른 생성원계이면
$g_s=\sum a_{sj}f_j$라 쓸 수 있고 행렬의 등식
$$
(\partial g_s/\partial x_i)=(a_{sj})(\partial f_j/\partial x_i)
+(\partial a_{sj}/\partial x_i f_j)
$$
이 성립한다. 마지막 항은 $A/I$에서 영이다. Cauchy--Binet 공식에 의해
$g_s$들에 대한 소행렬식 아이디얼은 $f_j$들에 대한 아이디얼에 포함된다.
대칭성에 의해 두 아이디얼은 같다. $y_1,\ldots,y_n\in\mathfrak m_A$가 또
다른 정칙 매개변수계이면 행렬 $(\partial y_j/\partial x_i)$는 가역이고,
미분의 연쇄법칙을 쓸 수 있다. 일부 세부사항은 생략한다.
\end{proof}

\begin{lemma}
\label{curves-lemma-fitting-ideal}
$k$를 체라 하자. $A=k[[x_1,\ldots,x_n]]$이라 하고
$I=(f_1,\ldots,f_m)\subset\mathfrak m_A$를 아이디얼이라 하자. 다음은 동치이다.
\begin{enumerate}
\item $k\to A/I$는 보조정리 \ref{curves-lemma-nodal-algebraic}에서와 같다.
\item $A/I$는 환원이고, 행렬 $(\partial f_j/\partial x_i)$의
$(n-1)\times(n-1)$ 소행렬식들은 $I+\mathfrak m_A$를 생성한다.
\item $\text{depth}(A/I)=1$이고, 행렬 $(\partial f_j/\partial x_i)$의
$(n-1)\times(n-1)$ 소행렬식들은 $I+\mathfrak m_A$를 생성한다.
\end{enumerate}
\end{lemma}

\begin{proof}
보조정리 \ref{curves-lemma-fitting-ideal-well-defined}에 의해 증명 도중 좌표계와 생성원계의
선택을 바꾸어도 된다.

\medskip\noindent
(1)이 성립하면 좌표를 바꾸어 $x_1,\ldots,x_{n-2}$가 $A/I$에서 영으로 가고,
어떤 비퇴화 이차식
$a x_{n-1}^2+b x_{n-1}x_n+c x_n^2$에 대하여
$A/I=k[[x_{n-1},x_n]]/(a x_{n-1}^2+b x_{n-1}x_n+c x_n^2)$이라고 할 수
있다. 그러면 직접 계산하여 (2), (3)이 모두 참임을 보일 수 있다.

\medskip\noindent
행렬 $(\partial f_j/\partial x_i)$의 $(n-1)\times(n-1)$ 소행렬식들이
$I+\mathfrak m_A$를 생성한다고 하자. 어떤 $i,j$에 대하여 편도함수
$\partial f_j/\partial x_i$가 $A$의 단위라고 하자. 그러면 매개변수계
$f_j,x_1,\ldots,x_{i-1},\hat x_i,x_{i+1},\ldots,x_n$과 $I$의 생성원계
$f_j,f_1,\ldots,f_{j-1},\hat f_j,f_{j+1},\ldots,f_m$을 쓸 수 있다.
따라서 정칙 매개변수계 $x_1,\ldots,x_n$과 $I$의 생성원
$x_1,f_2,\ldots,f_m$을 얻는다. 다음으로 $i\geq2$, $j\geq2$이고
$\partial f_j/\partial x_i$가 $A$의 단위인 쌍을 찾는다. 그러한 쌍이 있으면
위와 같이 바꾸어 정칙 매개변수계 $x_1,\ldots,x_n$과 $I$의 생성원
$x_1,x_2,f_3,\ldots,f_m$을 가진다고 할 수 있다. 이를 계속하면 유한 번 뒤에
정칙 매개변수계 $x_1,\ldots,x_n$과 $I$의 생성원
$x_1,\ldots,x_t,f_{t+1},\ldots,f_m$을 얻고, 모든 $i,j\geq t+1$에 대하여
$\partial f_j/\partial x_i\in\mathfrak m_A$이다.

\medskip\noindent
이 경우 편도함수 행렬은 블록 형태
$$
\left(
\begin{matrix}
I_{t\times t}&*\\
0&\mathfrak m_A
\end{matrix}
\right)
$$
를 가진다. 따라서 모든 $(n-1)\times(n-1)$ 소행렬식은
$\mathfrak m_A^{n-1-t}$에 속한다. 소행렬식 아이디얼이
$1$을 포함할 것이므로 $I\not =\mathfrak m_A$임을 주목하자. 따라서
$\mathfrak m_A\setminus\mathfrak m_A^2+I$에 원소가 있으므로(그렇지 않으면
Nakayama 보조정리에 의해 $I=\mathfrak m_A$이다) $n-1-t\leq1$이다.
% P07-ERR-0346
그러므로 $t\geq n-2$이다. 위에서 $t\not = n$임을 보았고, 마찬가지로
$t=n-1$이면 허용되지 않는 가역 $(n-1)\times(n-1)$ 소행렬식이 생긴다.
따라서 $t=n-2$이다. 이제 $A/I$는 $k[[x_{n-1},x_n]]$의 몫이고, 보조정리
\ref{curves-lemma-reduced-quotient-regular-ring-dim-2}에 의해 (2), (3)의 두 경우 모두
어떤 $f=f(x_{n-1},x_n)$에 대하여 $I$는
$x_1,\ldots,x_{n-2},f$로 생성된다. 이 경우 소행렬식에 대한 조건은 정확히
$f$의 이차항이 비퇴화라는 뜻, 즉 $A/I$가 보조정리
\ref{curves-lemma-nodal-algebraic}에서와 같다는 뜻이다.
\end{proof}

\begin{lemma}
\label{curves-lemma-nodal}
$k$를 체, $X$를 $1$차원 대수적 $k$-스킴, $x\in X$를 닫힌점이라 하자.
다음은 동치이다.
\begin{enumerate}
\item $x$는 마디점이다.
\item $k\to\mathcal{O}_{X,x}$는 보조정리 \ref{curves-lemma-nodal-algebraic}에서와 같다.
\item $x$로 가는 임의의 $\overline{x}\in X_{\overline{k}}$가 마디 특이점을 정의한다.
\item $\kappa(x)/k$는 분리이고 $\mathcal{O}_{X,x}$는 환원이며,
$\Omega_{X/k}$의 첫째 Fitting 아이디얼은 $\mathcal{O}_{X,x}$에서
$\mathfrak m_x$를 생성한다.
\item $\kappa(x)/k$는 분리이고 $\text{depth}(\mathcal{O}_{X,x})=1$이며,
$\Omega_{X/k}$의 첫째 Fitting 아이디얼은 $\mathcal{O}_{X,x}$에서
$\mathfrak m_x$를 생성한다.
\item $\kappa(x)/k$는 분리이고 $\mathcal{O}_{X,x}$는 환원이며,
$\delta$-불변량이 $1$이고 기하적 가지 수가 $2$이다.
\end{enumerate}
\end{lemma}

\begin{proof}
먼저 $k$가 대수적으로 닫혔다고 하자. 이 경우 (1)과 (3)의 동치는 자명하다.
두 변수의 비퇴화 이차식은 좌표변환을 제외하면 $xy$뿐이므로 (1), (3)과
(2)는 동치이다. (1)과 (6)의 동치는 보조정리
\ref{curves-lemma-multicross-algebra}이다. $X$를 $x$의 아핀 근방으로 바꾸어
$x$를 $0$으로 보내는 닫힌몰입 $X\to\mathbf A^n_k$가 있다고 가정해도 된다.
$f_1,\ldots,f_m\in k[x_1,\ldots,x_n]$을 $\mathbf A^n_k$에서 $X$의 아이디얼
$I$의 생성원이라 하자. 그러면 $\Omega_{X/k}$는
$R=k[x_1,\ldots,x_n]/I$-가군 $\Omega_{R/k}$에 대응하고, 이는 표시
$$
R^{\oplus m}\xrightarrow{(\partial f_j/\partial x_i)}
R^{\oplus n}\to\Omega_{R/k}\to0
$$
을 가진다(대수학, 절 \ref{algebra-section-differentials}와
\ref{algebra-section-netherlander} 참조). 따라서 $\Omega_{R/k}$의 첫째
Fitting 아이디얼은 행렬 $(\partial f_j/\partial x_i)$의
$(n-1)\times(n-1)$ 소행렬식들로 생성되는 아이디얼이다. 그러므로
극대 아이디얼 $(x_1,\ldots,x_n)$에서 $k[x_1,\ldots,x_n]\to R$을 완비한 것에
보조정리 \ref{curves-lemma-fitting-ideal}을 적용하면 (2), (4), (5)가 동치이다.

\medskip\noindent
이제 $k$를 임의의 체라 하자. 경우 (2), (4), (5), (6)에서는 잉여체
$\kappa(x)$가 $k$ 위에서 분리이다. (1), (3)에서도 그러함을 보이자.
$Z\subset X$를 $X$의 특이점자리, 즉 $\Omega_{X/k}$의 첫째 Fitting 아이디얼로 정의되는 닫힌
부분스킴이라 하자. $Z$의 형성은 체확장과 가환한다(인자, 보조정리
\ref{divisors-lemma-base-change-and-fitting-ideal-omega}). (1) 또는 (3)이
참이면 $X_{\overline{k}}$의 한 점 $\overline{x}$가 존재하여
$\overline{x}$는 $Z_{\overline{k}}$의 중복도 $1$인 고립점이다(보조정리의
조건들이 $\overline{k}$ 위에서 동치이기 때문이다). 특히
$Z_{\overline{x}}$는 $\overline{x}$에서 기하적으로 환원이다($\overline{k}$가
대수적으로 닫혔기 때문이다). % P07-ERR-0347
따라서 $Z$는 $x$에서 기하적으로 환원이다(다양체, 보조정리
\ref{varieties-lemma-geometrically-reduced-upstairs}). 특히 $Z$는 $x$에서
환원이고, 따라서 $x$의 어떤 근방에서 $Z=\Spec(\kappa(x))$이며
$\kappa(x)$는 $k$ 위에서 기하적으로 환원이다. 이는 $\kappa(x)/k$가
분리라는 뜻이다(대수학, 보조정리
\ref{algebra-lemma-characterize-separable-field-extensions}).

\medskip\noindent
앞 문단의 논증은 (1) 또는 (3)이 성립하면 $\Omega_{X/k}$의 첫째 Fitting
아이디얼이 $\mathfrak m_x$를 생성함을 보여 준다.
$\mathcal{O}_{X,x}\to\mathcal{O}_{X_{\overline{k}},\overline{x}}$는 평탄하고
$\mathcal{O}_{X_{\overline{k}},\overline{x}}$는 환원이고 깊이 $1$이므로 (4), (5)가
성립한다(대수학, 보조정리 \ref{algebra-lemma-descent-reduced}와
\ref{algebra-lemma-apply-grothendieck}를 쓰라). 반대로 (4)는 대수학, 보조정리
\ref{algebra-lemma-criterion-reduced}에 의해 (5)를 함의한다. (5)가 성립하면
$Z$는 $x$에서 기하적으로 환원이다($\kappa(x)/k$가 분리이고 어떤 근방에서
$Z$가 $x$이기 때문이다). 따라서 $Z_{\overline{k}}$는 $x$ 위에 놓인
$X_{\overline{k}}$의 임의의 점 $\overline{x}$에서 환원이다. 달리 말하면
$\Omega_{X_{\overline{k}}/\overline{k}}$의 첫째 Fitting 아이디얼은
$\mathcal{O}_{X_{\overline{k}},\overline{x}}$에서 $\mathfrak m_{\overline{x}}$를
생성한다. 또한 $\mathcal{O}_{X,x}\to
\mathcal{O}_{X_{\overline{k}},\overline{x}}$가 평탄이므로
$\text{depth}(\mathcal{O}_{X_{\overline{k}},\overline{x}})=1$이다(위 참고문헌 참조).
따라서 $\overline{x}\in X_{\overline{k}}$에 대하여 (5)가 성립하고, 대수적으로
닫힌 체 위에서의 동치에 의해 (3)이 성립한다. 이로써 (1), (3), (4), (5)는
동치이다.

\medskip\noindent
(2)와 (6)의 동치는 보조정리 \ref{curves-lemma-2-branches-delta-1}에서 따른다.

\medskip\noindent
마지막으로 (2) $=$ (6)이 (1) $=$ (3) $=$ (4) $=$ (5)와 동치임을 증명한다.
먼저 다양체, 보조정리
\ref{varieties-lemma-geometric-branches-and-change-of-fields}에 의해 $x$에서
$X$의 기하적 가지 수와 $\overline{x}$에서 $X_{\overline{k}}$의 기하적 가지
수는 같다. 지금까지 얻은 정보로부터 모든 경우에 이 수가 $2$임을 안다.
한편 경우 (1)에는 $X$가 $x$에서 기하적으로 환원임이 분명하고, 따라서
$$
\delta\text{-invariant of }X\text{ at }x\leq
\delta\text{-invariant of }X_{\overline{k}}\text{ at }\overline{x}
$$
이다(다양체, 보조정리
\ref{varieties-lemma-delta-invariant-and-change-of-fields}). 경우 (1)에서 우변은
$1$이므로 $x$에서 $X$의 $\delta$-불변량은 $1$이어야 한다(영이면 다양체,
보조정리 \ref{varieties-lemma-delta-invariant-is-zero}에 의해
$\mathcal{O}_{X,x}$가 이산값매김환이고, 이는 단일가지여서 모순이다). 따라서
(5)가 성립한다. % P07-ERR-0348
반대로 (2) $=$ (5)가 참이면 보조정리 \ref{curves-lemma-2-branches-delta-1}에 의해
$x\in X$에 대하여 다양체, 보조정리
\ref{varieties-lemma-geometrically-normal-in-codim-1}의 가정 (a), (b), (c)가
성립한다. % P07-ERR-0349
따라서 다양체, 보조정리
\ref{varieties-lemma-delta-invariant-and-change-of-fields-better}를 적용하면 위의
표시된 부등식에서 등호를 얻는다. 그러므로
$\overline{x}\in X_{\overline{k}}$에 대하여 (5)가 성립하고, 대수적으로 닫힌
체 위에서 조건들의 동치에 의해 다시 경우 (1)을 얻는다. % P07-ERR-0350
\end{proof}

\begin{remark}[마디점에 대응하는 이차확장]
\label{curves-remark-quadratic-extension}
$k$를 체, $(A,\mathfrak m,\kappa)$를 Noether 국소 $k$-대수라 하자.
$(A,\mathfrak m,\kappa)$가 보조정리 \ref{curves-lemma-nodal-algebraic}에서와 같거나,
$A$가 보조정리 \ref{curves-lemma-2-branches-delta-1}에서와 같은 Nagata 환이거나,
$A$가 완비이고 보조정리 \ref{curves-lemma-fitting-ideal}에서와 같다고 가정하자.
그러면 $A$는 다음과 같이 표준적으로 차수 $2$인 분리 $\kappa$-대수
$\kappa'$을 정의한다.
\begin{enumerate}
\item $q=ax^2+bxy+cy^2$를 보조정리 \ref{curves-lemma-nodal-algebraic}에서와 같은
비퇴화 이차식으로 택하되, $a\not =0$이 되도록 좌표 $x,y$를 택하고
$\kappa'=\kappa[x]/(ax^2+bx+c)$로 둔다.
\item $A'\supset A$를 $A$의 전분수환 안에서의 정수적 폐포라 하고
$\kappa'=A'/\mathfrak mA'$로 둔다.
\item
$\text{Proj}(\bigoplus_{n\geq0}\mathfrak m^n/\mathfrak m^{n+1})
=\Spec(\kappa')$가 되는 $\kappa$-대수를 $\kappa'$이라 한다.
\end{enumerate}
(1)과 (2)의 동치는 보조정리 \ref{curves-lemma-2-branches-delta-1}의 증명에서 보였다.
이들과 (3)의 동치는 생략한다. $X$가 국소 Noether $k$-스킴이고 $x\in X$가
$\mathcal{O}_{X,x}=A$인 점이면, (3)은 $\nu:X^\nu\to X$가 정규화 사상일 때
$\Spec(\kappa')=X^\nu\times_X\Spec(\kappa)$임을 보여 준다.
\end{remark}

\begin{remark}[자명한 이차확장]
\label{curves-remark-trivial-quadratic-extension}
$k$를 체, $(A,\mathfrak m,\kappa)$를 주의
\ref{curves-remark-quadratic-extension}에서와 같이 두고, $\kappa'/\kappa$를 이에
대응하는 차수 $2$인 분리대수라 하자. 다음은 동치이다.
\begin{enumerate}
\item $\kappa$-대수로서 $\kappa'\cong\kappa\times\kappa$이다.
\item 보조정리 \ref{curves-lemma-nodal-algebraic}의 형식 $q$를 $xy$로 택할 수 있다.
\item $A$는 두 개의 가지를 가진다.
\item 보조정리 \ref{curves-lemma-2-branches-delta-1}의 확장 $A'/A$는 두 극대 아이디얼을
가진다.
\item $k$-대수로서 $A^\wedge\cong\kappa[[x,y]]/(xy)$이다.
\end{enumerate}
이 조건들의 동치는 보조정리 \ref{curves-lemma-2-branches-delta-1}의 증명에서 보였다.
$X$가 국소 Noether $k$-스킴이고 $x\in X$가 $\mathcal{O}_{X,x}=A$인 점이면,
이는 정확히 정규화 $X^\nu$에 $x$ 위에 놓이는 두 점 $x_1,x_2$가 있고
$\kappa(x)=\kappa(x_1)=\kappa(x_2)$임을 뜻한다.
\end{remark}

\begin{definition}
\label{curves-definition-split-node}
$k$를 체라 하자. $X$를 $1$차원 대수적 $k$-스킴이라 하고,
$x\in X$를 닫힌점이라 하자. $x$가 마디점이고 $\kappa(x)=k$이며
$A=\mathcal{O}_{X,x}$에 대하여 주의
\ref{curves-remark-trivial-quadratic-extension}의 동치인 명제들이 성립하면,
$x$를 {\it 분할 마디점}이라 한다.
\end{definition}

\noindent
이 개념에 관하여 이미 알고 있는 내용을 진술하는 다음의 필수 보조정리를
적어 둔다.

\begin{lemma}
\label{curves-lemma-split-node}
$k$를 체라 하자. $X$를 $1$차원 대수적 $k$-스킴이라 하고,
$x\in X$를 닫힌점이라 하자. 다음은 동치이다.
\begin{enumerate}
\item $x$는 분할 마디점이다.
\item $x$는 마디점이고, $x$ 위에 놓이면서
$k=\kappa(x_1)=\kappa(x_2)$인 정규화 $X^\nu$의 점이 정확히 두 개
$x_1,x_2$ 있다.
\item $k$-대수로서
$\mathcal{O}_{X,x}^\wedge\cong k[[x,y]]/(xy)$이다.
% P07-ERR-0351: frozen source contains the editorial placeholder below.
\item 여기에 내용을 더 추가하라.
\end{enumerate}
\end{lemma}

\begin{proof}
주의 \ref{curves-remark-trivial-quadratic-extension}의 논의와
보조정리 \ref{curves-lemma-nodal}에서 따른다.
\end{proof}

\begin{lemma}
\label{curves-lemma-node-field-extension}
$K/k$를 체의 확장이라 하자. $X$를 차원 $1$인 국소 대수적
$k$-스킴이라 하자. $y\in X_K$를 상이 $x\in X$인 점이라 하자.
다음은 동치이다.
\begin{enumerate}
\item $x$는 $X$의 닫힌점이며 마디점이다.
% P07-ERR-0352: frozen source uses undefined Y throughout this lemma/proof.
\item $y$는 $Y$의 닫힌점이며 마디점이다.
\end{enumerate}
\end{lemma}

\begin{proof}
$x$가 $X$의 닫힌점이면 $y$도 닫힌점이다(잉여체를 보라).
그러나 역은 일반적으로 성립하지 않는다. 즉, $Y$의 닫힌점이 $X$의
닫히지 않은 점으로 갈 수 있다. 다만 이 경우 $y$는 마디점일 수 없다.
실제로 그러면 $x$는 $X$의 일반점이고 $\mathcal{O}_{X,x}$는 Artin 환이며
모든 Artin 국소환은 기하적으로 단일가지이므로, $X$는 $x$에서 기하적으로
단일가지이다. 따라서 $Y$는 $y$에서 기하적으로 단일가지이다
(다양체, 보조정리
\ref{varieties-lemma-geometrically-unibranch-and-change-of-fields}).
이는 보조정리 \ref{curves-lemma-nodal}에 의해 $y$가 마디점일 수 없음을 뜻한다.
그러므로 이제부터 $x$와 $y$가 모두 닫힌점이라고 가정해도 된다.

\medskip\noindent
대수적 폐포 $\overline{k}$와 $\overline{K}$ 및 주어진 사상
$k\to K$를 확장하는 사상 $\overline{k}\to\overline{K}$를 택하자.
보조정리 \ref{curves-lemma-nodal}의 (1)과 (3)의 동치를 이용하면 $k$와 $K$가
대수적으로 닫힌 경우로 환원된다. 이 경우 보조정리 \ref{curves-lemma-nodal}의
증명에서와 같이 논증하여 $x$에서 $X$의 기하적 가지 수와
$\delta$-불변량이 각각 $y$에서 $Y$의 것과 같음을 알 수 있다.
다른 논증은 다양체, 보조정리
\ref{varieties-lemma-complete-local-ring-structure-as-algebra}에서와 같은
$k$-대수의 동형
$$
k[[x_1,\ldots,x_n]]/(g_1,\ldots,g_m)
\to \mathcal{O}_{X,x}^\wedge
$$
을 택하여 얻을 수 있다. 다양체, 보조정리
\ref{varieties-lemma-base-change-complete-local-ring}에 의해 이것은
$K$-대수의 동형
$$
K[[x_1,\ldots,x_n]]/(g_1,\ldots,g_m)
\to \mathcal{O}_{Y,y}^\wedge
$$
을 준다. 정의에 따라 다음 두 조건이 동치임을 보이면 된다.
$$
k[[x_1,\ldots,x_n]]/(g_1,\ldots,g_m)\cong k[[s,t]]/(st)
$$
$$
K[[x_1,\ldots,x_n]]/(g_1,\ldots,g_m)\cong K[[s,t]]/(st).
$$
독자가 이를 직접 증명해 보기를 권한다. $k$와 $K$가 대수적으로 닫힌
체이므로 이는 이 환들이 보조정리 \ref{curves-lemma-nodal-algebraic}에서와 같은지를
묻는 것과 같다. 보조정리 \ref{curves-lemma-fitting-ideal}을 통하면 이는 행렬
$(\partial g_j/\partial x_i)$의 $(n-1)\times(n-1)$ 소행렬식에 관한 명제로
바뀌며, 이 명제는 사용한 체와 무관함이 분명하다. 세부사항은 생략한다.
\end{proof}

\begin{lemma}
\label{curves-lemma-node-etale-local}
$k$를 체라 하자. $X$를 차원 $1$인 국소 대수적 $k$-스킴이라 하자.
$Y\to X$를 \'etale 사상이라 하고, $y\in Y$를 상이 $x\in X$인 점이라
하자. 다음은 동치이다.
\begin{enumerate}
\item $x$는 $X$의 닫힌점이며 마디점이다.
\item $y$는 $Y$의 닫힌점이며 마디점이다.
\end{enumerate}
\end{lemma}

\begin{proof}
보조정리 \ref{curves-lemma-node-field-extension}에 의해 $k$의 대수적 폐포로
기저변환해도 된다. 그러면 $x$와 $y$의 잉여체는 $k$이다. 따라서 사상
$\mathcal{O}_{X,x}^\wedge\to\mathcal{O}_{Y,y}^\wedge$은 동형이다
(예를 들어 에탈 사상, 보조정리
\ref{etale-lemma-characterize-etale-completions} 또는 대수학 심화,
보조정리 \ref{more-algebra-lemma-flat-unramified}을 보라). 이제 보조정리는
분명하다.
\end{proof}

\begin{lemma}
\label{curves-lemma-node-over-separable-extension}
$k'/k$를 유한 분리 체확장이라 하자. $X$를 차원 $1$인 국소 대수적
$k'$-스킴이라 하고, $x\in X$를 닫힌점이라 하자. 다음은 동치이다.
\begin{enumerate}
\item $x$는 마디점이다.
\item $X$를 국소 대수적 $k$-스킴으로 볼 때 $x$는 마디점이다.
\end{enumerate}
\end{lemma}

\begin{proof}
보조정리 \ref{curves-lemma-nodal}의 마디점 특성화에서 바로 따른다.
\end{proof}

\begin{lemma}
\label{curves-lemma-nodal-lci}
$k$를 체라 하자. $X$를 차원이 $1$인 등차원 국소 대수적
$k$-스킴이라 하자. 다음은 동치이다.
\begin{enumerate}
\item $X$의 특이점은 기껏해야 마디점이다.
\item $X$는 $k$ 위의 국소 완전교차이고, $\Omega_{X/k}$의 첫째 Fitting
아이디얼로 잘라낸 닫힌부분스킴 $Z\subset X$는 $k$ 위에서 비분기이다.
\end{enumerate}
\end{lemma}

\begin{proof}
독자가 이 보조정리의 증명을 직접 찾아보기를 권한다. 다음 논증은 앞선
결과들을 한데 모은 것일 뿐이어서 실제로 일어나는 일을 가릴 수도 있다.

\medskip\noindent
(2)를 가정하자. $Z\to\Spec(k)$는 준유한이므로
(스킴의 사상, 보조정리 \ref{morphisms-lemma-unramified-quasi-finite}),
스킴의 사상, 보조정리 \ref{morphisms-lemma-residue-field-quasi-finite}에 의해
$x\in Z$인 점들의 잉여체는 $k$ 위에서 유한이고 (또한 분리이다).
따라서 스킴의 사상, 보조정리
\ref{morphisms-lemma-algebraic-residue-field-extension-closed-point-fibre}에
의해 각 $x\in Z$는 $X$의 닫힌점이다. 대수학, 보조정리
\ref{algebra-lemma-lci-CM}에 의해 국소환 $\mathcal{O}_{X,x}$는
Cohen--Macaulay 환이다. 차원이론에 의해
$\dim(\mathcal{O}_{X,x})=1$이므로
(다양체, 절 \ref{varieties-section-algebraic-schemes}),
% P07-ERR-0353: preserve the frozen unmatched parenthesis in the formula.
$\text{depth}(\mathcal{O}_{X,x}))=1$이다. 따라서 보조정리
\ref{curves-lemma-nodal}에 의해 $x$는 마디점이다. $x\in X$이고 $x\not \in Z$이면
인자, 보조정리 \ref{divisors-lemma-d-fitting-ideal-omega-smooth}에 의해
$X\to\Spec(k)$는 $x$에서 매끄럽다.

\medskip\noindent
(1)을 가정하자. 이 가정 아래 $X$는 모든 닫힌점에서 기하적으로 환원이다
(다양체, 보조정리 \ref{varieties-lemma-geometrically-reduced-upstairs} 참조).
따라서 다양체, 보조정리
\ref{varieties-lemma-geometrically-reduced-dense-smooth-open}에 의해
$X\to\Spec(k)$는 조밀한 열린집합 위에서 매끄럽다. 그러므로 $Z$는 닫혀
있고 닫힌점들로 이루어진다. 인자, 보조정리
\ref{divisors-lemma-d-fitting-ideal-omega-smooth}에 의해 사상
$X\setminus Z\to\Spec(k)$는 매끄럽다. 따라서 스킴의 사상, 보조정리
\ref{morphisms-lemma-smooth-syntomic}와 스킴의 사상, 정의
\ref{morphisms-definition-syntomic}의 국소 완전교차 정의에 의해
$X\setminus Z$는 국소 완전교차이다. 모든 $x\in Z$에 대하여 보조정리
\ref{curves-lemma-nodal}에 의해 국소환 $\mathcal{O}_{Z,x}$는 $\kappa(x)$와 같고
$\kappa(x)$는 $k$ 위에서 분리이다. 따라서
$Z\to\Spec(k)$는 비분기이다
(스킴의 사상, 보조정리 \ref{morphisms-lemma-unramified-over-field}).
마지막으로 보조정리 \ref{curves-lemma-nodal-algebraic}의 (3)을 통한 보조정리
\ref{curves-lemma-nodal}에 의해 $\mathcal{O}_{X,x}$는 분할멱 대수,
정의 \ref{dpa-definition-lci}의 의미에서 완전교차이다. 그런데
분할멱 대수, 보조정리 \ref{dpa-lemma-check-lci-agrees}와 스킴의 사상,
보조정리 \ref{morphisms-lemma-local-complete-intersection}은 이것이 체 위의
국소 완전교차 스킴을 정의하는 데 사용한 개념과 일치함을 보여 준다.
이로써 증명이 끝난다.
\end{proof}

\begin{lemma}
\label{curves-lemma-facts-about-nodal-curves}
$k$를 체라 하자. $X$를 차원이 $1$인 등차원 국소 대수적 $k$-스킴이라
하고, 그 특이점이 기껏해야 마디점이라고 하자. 그러면 $X$는 Gorenstein
스킴이며 기하적으로 환원이다.
\end{lemma}

\begin{proof}
Gorenstein이라는 명제는 보조정리 \ref{curves-lemma-nodal-lci}와 스킴의 쌍대성, 보조정리 \ref{duality-lemma-gorenstein-lci}에서 따른다. 또는 대수적
폐포로 옮긴 뒤 확인하면 충분하다는 사실(스킴의 쌍대성, 보조정리
\ref{duality-lemma-gorenstein-base-change})과 Noether 국소환이
Gorenstein일 필요충분조건은 그 완비화가 Gorenstein이라는 사실
(쌍대화 복합체, 보조정리
\ref{dualizing-lemma-flat-under-gorenstein})을 이용한 다음, 국소환
$k[[t]]$와 $k[[x,y]]/(xy)$가 Gorenstein임을 직접 증명할 수도 있다.

\medskip\noindent
$X$가 기하적으로 환원임을 보이려면 $X_{\overline{k}}$가 환원임을 보이면
충분하다(다양체, 보조정리 \ref{varieties-lemma-perfect-reduced} 및
\ref{varieties-lemma-geometrically-reduced}). 그런데 $X_{\overline{k}}$는
대수적으로 닫힌 체 위의 마디곡선이다. 따라서 $X_{\overline{k}}$의 완비
국소환들은 $\overline{k}[[t]]$ 또는 $\overline{k}[[x,y]]/(xy)$와 동형이며,
이들은 원하는 대로 환원환이다.
\end{proof}

\begin{lemma}
\label{curves-lemma-closed-subscheme-nodal-curve}
$k$를 체라 하자. $X$를 차원이 $1$인 등차원 국소 대수적 $k$-스킴이라
하고, 그 특이점이 기껏해야 마디점이라고 하자. $Y\subset X$가 차원이
$1$인 등차원 환원 닫힌부분스킴이면 다음이 성립한다.
\begin{enumerate}
\item $Y$의 특이점은 기껏해야 마디점이다.
\item $Z\subset X$가 $X\setminus Y$의 스킴론적 폐포이면 다음이 성립한다.
\begin{enumerate}
\item 스킴론적 교집합 $Y\cap Z$는 $k$의 유한 분리확장들의 스펙트럼의
서로소 합이다.
\item $Y\cap Z$의 각 점은 $X$의 마디점이다.
\item $Y\to\Spec(k)$는 $Y\cap Z$의 모든 점에서 매끄럽다.
\end{enumerate}
\end{enumerate}
\end{lemma}

\begin{proof}
$X$와 $Y$는 차원이 $1$인 환원 등차원 스킴이므로, $Y$는 $X$의 기약성분
중 일부의 스킴론적 합이다(환원환에서 $(0)$은 극소 소 아이디얼들의
교집합이다). $y\in Y$를 닫힌점이라 하자. $y$가
$X\to\Spec(k)$의 매끄러운 자리에 있으면 $y$는 $X$의 유일한 기약성분
위에 놓이며, $y$의 어떤 열린근방에서 $Y$와 $X$가 일치함을 알 수 있다.
따라서 $Y\to\Spec(k)$는 $y$에서 매끄럽다. $y$가 $X$의 마디점이지만
여전히 $X$의 유일한 기약성분 위에 놓이면, 같은 논증으로 $y$는 $Y$의
마디점이다. 이제 $y$가 $X$의 $1$개보다 많은 기약성분 위에 놓인다고 하자.
보조정리 \ref{curves-lemma-nodal}에 의해 $y$에서 $X$의 기하적 가지 수는 $2$이므로,
스킴의 성질, 보조정리
\ref{properties-lemma-number-of-branches-irreducible-components}에 의해
$y$를 지나는 기약성분은 기껏해야 $2$개이다. $Y$가 이 둘을 모두 포함하면
다시 $y$의 어떤 열린근방에서 $Y=X$이고 $y$는 $Y$의 마디점이다.
마지막으로 $Y$가 그 기약성분 중 하나만 포함한다고 하자.
% P07-ERR-0354: frozen source says neighbourhood of undefined x; preserve x.
$X$를 $x$의 열린근방으로 바꾼 뒤 $Y$는 두 기약성분 중 하나이고 $Z$는
다른 하나라고 가정할 수 있다.
% P07-ERR-0355: frozen source typo irreducble; translated as intended noun.
스킴의 성질, 보조정리
\ref{properties-lemma-number-of-branches-irreducible-components}를 다시 쓰면
$X$가 $y$에서 두 가지를 가짐을 알 수 있다. 즉, 국소환
$\mathcal{O}_{X,y}$는 두 가지를 가지며 이 가지들은
$\mathcal{O}_{Y,y}$와 $\mathcal{O}_{Z,y}$에서 온다. 주의
\ref{curves-remark-trivial-quadratic-extension}에서와 같이
$\mathcal{O}_{X,y}^\wedge\cong\kappa(y)[[u,v]]/(uv)$로 쓰자. 예를 들어
보조정리 \ref{curves-lemma-nodal}에 의해 체 $\kappa(y)$는 $k$ 위에서 유한 분리이다.
따라서 필요하면 $u$와 $v$의 역할을 바꾸어, 사상
% P07-ERR-0356: preserve frozen wrong point O_{Y,Y}.
$\mathcal{O}_{X,y}\to\mathcal{O}_{Y,Y}$의 완비화가
$\kappa(y)[[u,v]]/(uv)\to\kappa(y)[[u]]$에 대응하고, 사상
% P07-ERR-0357: preserve frozen second O_{Y,Y} token.
$\mathcal{O}_{X,y}\to\mathcal{O}_{Y,Y}$의 완비화가
$\kappa(y)[[u,v]]/(uv)\to\kappa(y)[[v]]$에 대응한다고 할 수 있다.
% P07-ERR-0358: frozen source says ideas; translated as mathematical ideals.
스킴론적 교집합 $Y\cap Z$는 두 아이디얼의 합으로 잘라내어지며,
완비화에서는 이것이 $(u,v)$, 즉 극대 아이디얼이다. 그러므로 (2)(a)와
(2)(b)는 분명하다. 마지막으로 (2)(c)가 성립한다. $\mathcal{O}_{Y,y}$의
완비화가 정칙이므로 $\mathcal{O}_{Y,y}$도 정칙이고
(대수학 심화, 보조정리 \ref{more-algebra-lemma-completion-regular}),
$\kappa(y)/k$는 분리이므로 대수학, 보조정리
% P07-ERR-0359: source ends with a sentence fragment; translated grammatically.
\ref{algebra-lemma-separable-smooth}에 의해 어떤 열린근방에서 매끄럽다.
\end{proof}




\section{마디곡선의 족}
\label{curves-section-families-nodal}

\noindent
Stacks 프로젝트에서 곡선은 차원이 $1$인 기약다양체이지만, 문헌에서
``반안정 곡선'' 또는 ``마디곡선''은 보통 기약이 아니며, 특히 ``반안정
곡선의 족'', ``마디곡선의 족'', 또는 ``마디족'' 같은 표현에서는 흔히
고유라고 가정한다. 따라서 문헌과 양립하면서 내부적으로도 일관된 용어를
택하기가 다소 어렵다. 더구나 우리는 먼저 다음 보조정리에 도입하는
개념(원천에서 국소적인 개념)을 연구하고자 한다.

\begin{lemma}
\label{curves-lemma-nodal-family}
$f:X\to S$를 스킴의 사상이라 하자. 다음은 동치이다.
\begin{enumerate}
\item $f$는 평탄하고 국소 유한표현이며, 공집합이 아닌 모든 올
$X_s$는 차원이 $1$인 등차원 스킴이고 $X_s$는 기껏해야 마디점인 특이점만 갖는다.
\item $f$는 상대차원 $1$인 syntomic 사상이고, $\Omega_{X/S}$의 첫째
Fitting 아이디얼로 정의되는 닫힌부분스킴
$\text{Sing}(f)\subset X$는 $S$ 위에서 비분기이다.
\end{enumerate}
\end{lemma}

\begin{proof}
$\text{Sing}(f)$의 형성은 기저변환과 가환한다는 것을 상기하자.
인자, 보조정리
\ref{divisors-lemma-base-change-and-fitting-ideal-omega}를 보라.
따라서 이 보조정리는 보조정리 \ref{curves-lemma-nodal-lci}, 스킴의 사상, 보조정리
\ref{morphisms-lemma-syntomic-flat-fibres}, 그리고 스킴의 사상, 보조정리
\ref{morphisms-lemma-unramified-etale-fibres}에서 따른다.
(자명한 스킴의 사상, 보조정리
\ref{morphisms-lemma-syntomic-locally-finite-presentation} 및
\ref{morphisms-lemma-syntomic-flat}도 사용한다.)
\end{proof}

\begin{definition}
\label{curves-definition-nodal-family}
$f:X\to S$를 스킴의 사상이라 하자. $f$가 보조정리
\ref{curves-lemma-nodal-family}의 동치인 조건들을 만족하면 $f$를
{\it 상대차원 $1$인 기껏해야 마디인 사상}이라 한다.
\end{definition}

\noindent
이처럼 번거로운 용어를 쓰는 데에는 몇 가지 이유가 있다\footnote{더 나은
제안이 있다면 Stacks 프로젝트 관리자에게 전자우편을 보내 주기 바란다.}.
첫째, 이 개념을 문헌의 다른 개념들과 혼동하지 않기를 원한다(이 절의
서론을 보라). 둘째, 이 개념을 더 큰 상대차원의 사상으로 일반화하는 여러
방식을 생각할 수 있다. 예를 들어 \'etale 국소적으로 기껏해야 마디인
사상들의 합성인 사상을 요구할 수도 있고, 올의 차원은 더 크지만 기껏해야
보통 이중점인 특이점만 갖는 사상을 요구할 수도 있다.

\begin{lemma}
\label{curves-lemma-smooth-relative-dimension-1}
상대차원 $1$인 매끄러운 사상은 상대차원 $1$인 기껏해야 마디인
사상이다.
\end{lemma}

\begin{proof}
생략한다.
\end{proof}

\begin{lemma}
\label{curves-lemma-base-change-nodal-family}
$f:X\to S$가 상대차원 $1$인 기껏해야 마디인 사상이면, $f$의 임의의
기저변환도 그러하다.
\end{lemma}

\begin{proof}
syntomic 사상의 기저변환은 syntomic이고
(스킴의 사상, 보조정리 \ref{morphisms-lemma-base-change-syntomic}),
상대차원 $1$인 사상의 기저변환은 상대차원 $1$이며
(스킴의 사상, 보조정리
\ref{morphisms-lemma-base-change-relative-dimension-d}),
$\text{Sing}(f)$의 형성은 기저변환과 가환하고
(인자, 보조정리
\ref{divisors-lemma-base-change-and-fitting-ideal-omega}),
비분기 사상의 기저변환은 비분기이기 때문이다
(스킴의 사상, 보조정리 \ref{morphisms-lemma-base-change-unramified}).
\end{proof}

\noindent
다음 보조정리는 사상이 상대차원 $1$인 기껏해야 마디인 사상인지를
그 올들 위에서 확인할 수 있음을 말한다.

\begin{lemma}
\label{curves-lemma-locus-where-nodal}
$f:X\to S$를 평탄하고 국소 유한표현인 스킴의 사상이라 하자. 제한
$f|_U:U\to S$가 상대차원 $1$인 기껏해야 마디인 사상이 되는 최대
열린부분스킴 $U\subset X$가 존재한다. 더욱이 $U$의 형성은 임의의
기저변환과 가환한다.
\end{lemma}

\begin{proof}
스킴의 사상, 보조정리 \ref{morphisms-lemma-set-points-where-fibres-lci}에 의해
$f$가 syntomic인 그러한 열린집합이 존재한다. 따라서 $f$가 syntomic인
사상이라고 가정해도 된다. 특히 $f$는 Cohen--Macaulay 사상이다
(스킴의 쌍대성, 보조정리 \ref{duality-lemma-lci-gorenstein} 및
\ref{duality-lemma-gorenstein-CM-morphism}). 그러므로 $X$는 $f$가 주어진
상대차원을 갖는 열린닫힌부분스킴들의 서로소 합이다. 스킴의 사상, 보조정리
\ref{morphisms-lemma-flat-finite-presentation-CM-fibres-relative-dimension}를
보라. 이 분해는 임의의 기저변환으로 보존된다. 스킴의 사상, 보조정리
\ref{morphisms-lemma-base-change-relative-dimension-d}를 보라. 하나만
남기고 나머지 조각을 버려 $f$가 상대차원 $1$인 syntomic 사상이라고
가정할 수 있다. $\text{Sing}(f)\subset X$를 $\Omega_{X/S}$의 첫째
Fitting 아이디얼로 정의되는 닫힌부분스킴이라 하자.
% P07-ERR-1093: frozen source misspells subscheme as subscheem.
$W\to S$가 비분기이고 그 형성이 기저변환과 가환하는 최대 열린부분스킴
$W\subset\text{Sing}(f)$가 존재한다
(스킴의 사상, 보조정리
\ref{morphisms-lemma-set-points-where-fibres-unramified}). 또한
$\text{Sing}(f)$의 형성도 기저변환과 가환하므로
(인자, 보조정리
\ref{divisors-lemma-base-change-and-fitting-ideal-omega}),
$$
U=(X\setminus\text{Sing}(f))\cup W
$$
가 $f|_U:U\to S$가 상대차원 $1$인 기껏해야 마디인 사상이 되는 $X$의
최대 열린부분스킴이며, $U$의 형성이 기저변환과 가환함을 알 수 있다.
\end{proof}

\begin{lemma}
\label{curves-lemma-nodal-family-precompose-etale}
$f:X\to S$가 상대차원 $1$인 기껏해야 마디인 사상이라고 하자.
$Y\to X$가 \'etale 사상이면 합성 $g:Y\to S$도 상대차원 $1$인
기껏해야 마디인 사상이다.
\end{lemma}

\begin{proof}
$g$는 평탄하고 국소 유한표현인 사상들의 합성이므로 평탄하고 국소
유한표현이다(스킴의 사상, 보조정리
\ref{morphisms-lemma-etale-locally-finite-presentation},
\ref{morphisms-lemma-etale-flat},
\ref{morphisms-lemma-composition-finite-presentation}, 그리고
\ref{morphisms-lemma-composition-flat}을 사용하라). 따라서 올들이 기껏해야
마디점인 특이점만 가짐을 보이면 충분하다. 이는 보조정리
\ref{curves-lemma-node-etale-local}에서 따른다. 또한 \'etale 사상과 매끄러운
사상의 합성은 스킴의 사상, 보조정리
\ref{morphisms-lemma-etale-smooth-unramified} 및
\ref{morphisms-lemma-composition-smooth}에 의해 매끄럽다는 사실을 쓴다.
\end{proof}

\begin{lemma}
\label{curves-lemma-nodal-family-postcompose-etale}
$S'\to S$를 스킴의 \'etale 사상이라 하자. $f:X\to S'$가 상대차원
$1$인 기껏해야 마디인 사상이면, 합성 $g:X\to S$도 상대차원 $1$인
기껏해야 마디인 사상이다.
\end{lemma}

\begin{proof}
$g$는 평탄하고 국소 유한표현인 사상들의 합성이므로 평탄하고 국소
유한표현이다(스킴의 사상, 보조정리
\ref{morphisms-lemma-etale-locally-finite-presentation},
\ref{morphisms-lemma-etale-flat},
\ref{morphisms-lemma-composition-finite-presentation}, 그리고
\ref{morphisms-lemma-composition-flat}을 사용하라). 따라서 $g$의 올들이
기껏해야 마디점인 특이점만 가짐을 보이면 충분하다. 이는 보조정리
\ref{curves-lemma-node-over-separable-extension}와 매끄러운 점에 관한 유사한
결과에서 따른다.
\end{proof}

\begin{lemma}
\label{curves-lemma-nodal-family-etale-local-source}
$f:X\to S$를 스킴의 사상이라 하고, $\{U_i\to X\}$를 \'etale
피복이라 하자. 다음은 동치이다.
\begin{enumerate}
\item $f$는 상대차원 $1$인 기껏해야 마디인 사상이다.
\item 각 $U_i\to S$는 상대차원 $1$인 기껏해야 마디인 사상이다.
\end{enumerate}
달리 말해, 상대차원 $1$인 기껏해야 마디인 사상이라는 성질은 원천에서
\'etale 국소적이다.
\end{lemma}

\begin{proof}
한 방향은 보조정리 \ref{curves-lemma-nodal-family-precompose-etale}에서 보았다.
다른 방향을 위해, 국소 유한표현, 평탄성, 상대차원 $1$이라는 성질이
원천에서 \'etale 국소적임을 보자(내림, 보조정리
\ref{descent-lemma-locally-finite-presentation-fppf-local-source},
\ref{descent-lemma-flat-fpqc-local-source}, 그리고
\ref{descent-lemma-dimension-at-point}). 올을 취하면 $S$가 한 체의
스펙트럼인 경우로 환원된다. 이 경우 결과는 보조정리
\ref{curves-lemma-node-etale-local}에서 따른다. 매끄러움이 원천에서 \'etale
국소적이라는 사실은 내림, 보조정리
\ref{descent-lemma-smooth-smooth-local-source}를 보라.
\end{proof}

\begin{lemma}
\label{curves-lemma-nodal-family-fpqc-local-target}
$f:X\to S$를 스킴의 사상이라 하고, $\{U_i\to S\}$를 fpqc 피복이라
하자. 다음은 동치이다.
\begin{enumerate}
\item $f$는 상대차원 $1$인 기껏해야 마디인 사상이다.
\item 각 $X\times_SU_i\to U_i$는 상대차원 $1$인 기껏해야 마디인
사상이다.
\end{enumerate}
달리 말해, 상대차원 $1$인 기껏해야 마디인 사상이라는 성질은 목표에서
fpqc 국소적이다.
\end{lemma}

\begin{proof}
한 방향은 보조정리 \ref{curves-lemma-base-change-nodal-family}에서 보았다. 다른
방향을 위해, 국소 유한표현, 평탄성, 상대차원 $1$이라는 성질이 목표에서
fpqc 국소적임을 보자(내림, 보조정리
\ref{descent-lemma-descending-property-locally-finite-presentation},
\ref{descent-lemma-descending-property-flat}, 그리고 스킴의 사상, 보조정리
\ref{morphisms-lemma-dimension-fibre-after-base-change}). 올을 취하면 $S$가
한 체의 스펙트럼인 경우로 환원된다. 이 경우 결과는 보조정리
\ref{curves-lemma-node-field-extension}에서 따른다. 매끄러움이 목표에서 fpqc
국소적이라는 사실은 내림, 보조정리
\ref{descent-lemma-descending-property-smooth}를 보라.
\end{proof}

\begin{lemma}
\label{curves-lemma-descend-nodal-family}
$S=\lim S_i$를 아핀 전이사상을 갖는 유향계의 스킴들의 극한이라 하자.
$0\in I$이고 $f_0:X_0\to Y_0$를 $S_0$ 위 스킴의 사상이라 하자.
$S_0,X_0,Y_0$가 준콤팩트이고 준분리라고 가정하자. $f_i:X_i\to Y_i$를
$f_0$의 $S_i$로의 기저변환이라 하고, $f:X\to Y$를 $f_0$의 $S$로의
기저변환이라 하자. 다음을 가정하자.
\begin{enumerate}
\item $f$는 상대차원 $1$인 기껏해야 마디인 사상이다.
\item $f_0$는 국소 유한표현이다.
\end{enumerate}
그러면 어떤 $i\geq0$에 대하여 $f_i$는 상대차원 $1$인 기껏해야 마디인
사상이다.
\end{lemma}

\begin{proof}
극한, 보조정리 \ref{limits-lemma-descend-syntomic}에 의해 어떤 $i$에
대하여 $f_i$는 syntomic이다. 그러면
$X_i=\coprod_{d\geq0}X_{i,d}$는 $X_{i,d}\to Y_i$가 상대차원 $d$인
열린닫힌부분스킴들의 서로소 합이다. 스킴의 사상, 보조정리
\ref{morphisms-lemma-syntomic-relative-dimension}을 보라. 스킴의 사상,
보조정리 \ref{morphisms-lemma-dimension-fibre-after-base-change}이 주는
기저변환 아래 올 차원의 거동에 의해 $X\to X_i$의 상은 $X_{i,1}$에
포함된다. 그러면 극한, 보조정리
\ref{limits-lemma-limit-contained-in-constructible}에 의해 어떤
$i'\geq i$에 대하여 $X_{i'}\to X_i$의 상도 $X_{i,1}$에 포함된다.
따라서 $f_{i'}:X_{i'}\to Y_{i'}$는 상대차원 $1$인 syntomic 사상이다
(스킴의 사상, 보조정리
\ref{morphisms-lemma-dimension-fibre-after-base-change}을 다시 쓴다).
사상 $\text{Sing}(f_{i'})\to Y_{i'}$를 생각하자. $Y$로의 기저변환은
비분기 사상임을 알고 있다. 따라서 극한, 보조정리
\ref{limits-lemma-descend-unramified}에 의해 $i'$을 충분히 크게 하면
사상 $\text{Sing}(f_{i'})\to Y_{i'}$가 비분기가 된다. 이로써 증명이
끝난다.
\end{proof}

\begin{lemma}
\label{curves-lemma-formal-local-structure-nodal-family}
$f:T\to S$를 스킴의 사상이라 하고, $t\in T$의 상을 $s\in S$라 하자.
다음을 가정하자.
\begin{enumerate}
\item $f$는 $t$에서 평탄하다.
\item $\mathcal{O}_{S,s}$는 Noether 환이다.
\item $f$는 국소 유한형이다.
\item $t$는 올 $T_s$의 분할 마디점이다.
\end{enumerate}
그러면 어떤 $h\in\mathfrak m_s^\wedge$와
$\mathcal{O}_{S,s}^\wedge$-대수의 동형
$$
\mathcal{O}_{T,t}^\wedge\cong
\mathcal{O}_{S,s}^\wedge[[x,y]]/(xy-h)
$$
이 존재한다.
\end{lemma}

\begin{proof}
$S$를 $\Spec(\mathcal{O}_{S,s})$로, $T$를
$\Spec(\mathcal{O}_{S,s})$로의 기저변환으로 바꾼다. 그러면 $T$는 국소
Noether이고 따라서 $\mathcal{O}_{T,t}$는 Noether 환이다.
$A=\mathcal{O}_{S,s}^\wedge$, $\mathfrak m=\mathfrak m_A$, 그리고
$B=\mathcal{O}_{T,t}^\wedge$로 두자. 대수학 심화, 보조정리
\ref{more-algebra-lemma-flat-completion}에 의해 $A\to B$는 평탄하다.
$$
\mathcal{O}_{T,t}/\mathfrak m_s\mathcal{O}_{T,t}=\mathcal{O}_{T_s,t}
$$
이므로 $B/\mathfrak mB=\mathcal{O}_{T_s,t}^\wedge$이다. 가정 (4)와
보조정리 \ref{curves-lemma-split-node}에 의해 어떤
$\overline{u},\overline{v}\in B/\mathfrak mB$가 존재하여 사상
% P07-ERR-0361: frozen source maps x to both generators; preserve both x tokens.
$$
(A/\mathfrak m)[[x,y]]\longrightarrow B/\mathfrak mB,\quad
x\longmapsto\overline{u},
x\longmapsto\overline{v}
$$
이 전사이고 핵이 $(xy)$이다.

\medskip\noindent
$n\geq1$이고 $\overline{u},\overline{v}$로 가는 $u,v\in B$가 있어
$$
uv=h+\delta
$$
라고 하자. 여기서 $h\in A$이고 $\delta\in\mathfrak m^nB$이다.
$u-u',v-v'\in\mathfrak m^nB$이고
$$
u'v'=h'+\delta'
$$
이 되는 $u',v'\in B$, $h'\in A$, 그리고
$\delta'\in\mathfrak m^{n+1}B$가 존재한다고 주장한다. 이를 보이기 위해
$f_i\in\mathfrak m^n$, $b_i\in B$인
$\delta=\sum f_ib_i$로 쓰자. 이어서
$$
b_i=a_i+ub_{i,1}+vb_{i,2}+\delta_i
$$
로 쓴다. 여기서 $a_i\in A$, $b_{i,1},b_{i,2}\in B$, 그리고
$\delta_i\in\mathfrak mB$이다. $B$의 잉여체가 $A$의 잉여체와 같고,
$u,v$의 $B/\mathfrak mB$에서의 상이 극대 아이디얼을 생성하므로 이렇게
쓸 수 있다. 이제
$$
u'=u-\sum b_{i,2}f_i,\quad
v'=v-\sum b_{i,1}f_i
$$
로 두면, 어떤 $c_{i,j}\in B$에 대하여
$$
u'v'=h+\delta-\sum(b_{i,1}u+b_{i,2}v)f_i+\sum
c_{ij}f_if_j=
h+\sum a_if_i+\sum f_i\delta_i+\sum c_{ij}f_if_j
$$
을 얻는다. 따라서 $h'=h+\sum a_if_i$와
$\delta'=\sum f_i\delta_i+\sum c_{ij}f_if_j$로 두면 위와 같은 식을
얻는다.

\medskip\noindent
귀납적으로 논증하자. $\overline{u},\overline{v}$의 임의의 올림
$u_1,v_1\in B$에서 시작하면 앞 단락의 결과로
$u_n,v_n\in B$와 $h_n\in A$의 열을 얻되,
$u_n-u_{n+1}\in\mathfrak m^nB$,
$v_n-v_{n+1}\in\mathfrak m^nB$,
$h_n-h_{n+1}\in\mathfrak m^n$, 그리고
$u_nv_n-h_n\in\mathfrak m^nB$이다. $A$와 $B$가 완비이므로
$u_\infty=\lim u_n$, $v_\infty=\lim v_n$, 그리고
$h_\infty=\lim h_n$으로 둘 수 있고, $B$에서
$u_\infty v_\infty=h_\infty$를 얻는다. 따라서 $x$를 $u_\infty$로,
% P07-ERR-0362: frozen source says v maps to v_infinity; preserve the v token.
$v$를 $v_\infty$로 보내는 $A$-대수 사상
$$
A[[x,y]]/(xy-h_\infty)\longrightarrow B
$$
을 얻는다. 이는 $\mathfrak m$으로 나눈 뒤 동형이 되는 평탄한
$A$-대수들의 사상이다. $\mathfrak m$을 법으로 전사이므로 완비성과
대수학, 보조정리 \ref{algebra-lemma-completion-generalities}에 의해 전사이다.
이제 대수학, 보조정리 \ref{algebra-lemma-mod-injective}를 적용하면 이것이
동형임을 알 수 있다.
\end{proof}

\begin{lemma}
\label{curves-lemma-genus-in-nodal-family-of-curves}
$f:X\to S$를 스킴의 사상이라 하자. 다음을 가정하자.
\begin{enumerate}
\item $f$는 고유이다.
\item $f$는 상대차원 $1$인 기껏해야 마디인 사상이다.
\item $f$의 기하적 올들은 연결이다.
\end{enumerate}
그러면 (a) $f_*\mathcal{O}_X=\mathcal{O}_S$이고 이는 임의의 기저변환
뒤에도 성립하며, (b) $R^1f_*\mathcal{O}_X$는 유한 국소자유
$\mathcal{O}_S$-가군이고 그 형성은 임의의 기저변환과 가환하며, (c)
$q\geq2$이면 $R^qf_*\mathcal{O}_X=0$이다.
\end{lemma}

\begin{proof}
(a)는 스킴의 유도 범주, 보조정리
\ref{perfect-lemma-proper-flat-geom-red-connected}에서 따른다.
스킴의 유도 범주, 보조정리
\ref{perfect-lemma-proper-flat-h0}에 의해 $S$ 위에서 국소적으로
$$
Rf_*\mathcal{O}_X=\mathcal{O}_S\oplus P
$$
로 쓸 수 있다. 여기서 $P$는 $[1,\infty)$ 안의 Tor 진폭을 갖는
완전 복합체이다. $Rf_*\mathcal{O}_X$의 형성은 임의의 기저변환과
가환함을 상기하자(스킴의 유도 범주, 보조정리
\ref{perfect-lemma-flat-proper-perfect-direct-image-general}). 따라서
$s\in S$에 대하여
$$
H^i(P\otimes_{\mathcal{O}_S}^{\mathbf{L}}\kappa(s))=
H^i(X_s,\mathcal{O}_{X_s})
\text{ for }i\geq1
$$
이다. $X_s$는 $1$차원 Noether 스킴이므로 이는 $i=1$이 아니면 영이다.
코호몰로지, 명제
\ref{cohomology-proposition-vanishing-Noetherian}을 보라. 그러면
$P=H^1(P)[-1]$이고 $H^1(P)$는 유한 국소자유이다. 예를 들어 대수학 심화, 보조정리 \ref{more-algebra-lemma-lift-perfect-from-residue-field}를
보라. 모든 것이 기저변환과 양립하므로 결론을 얻는다.
\end{proof}





\section{마디족의 \'Etale 국소구조}
\label{curves-section-etale-local-nodal}


\noindent
스킴의 사상
$$
\Spec(\mathbf{Z}[u,v,a]/(uv-a))
\longrightarrow
\Spec(\mathbf{Z}[a])
$$
을 생각하자. 다음 보조정리는 이 사상이 마디곡선 족의 \'etale
국소구조에 대한 모형임을 보여 준다.

\begin{lemma}
\label{curves-lemma-etale-local-structure-nodal-family}
$f:X\to S$를 스킴의 사상이라 하자. $f$가 상대차원 $1$인 기껏해야
마디인 사상이라고 가정하자. $x\in X$를 올 $X_s$의 특이점이라 하자.
그러면 다음과 같은 스킴의 가환도표가 존재한다.
$$
\xymatrix{
X \ar[d] &
U \ar[rr] \ar[l] \ar[rd] & &
W \ar[r] \ar[ld] &
\Spec(\mathbf{Z}[u, v, a]/(uv - a)) \ar[d] \\
S & &
V \ar[ll] \ar[rr] & & \Spec(\mathbf{Z}[a])
}
$$
여기서 $X\leftarrow U$, $S\leftarrow V$, $U\to W$는 \'etale 사상이고,
오른쪽 사각형은 데카르트이며, $X$에서 $x$로 가는 점 $u\in U$가
존재한다.
\end{lemma}

\begin{proof}[Artin 근사를 사용하는 첫째 증명]
먼저 절대 Noether 근사를 써서 $\mathbf{Z}$ 위 유한형인 스킴들의 경우로
환원한다. 문제는 $X$와 $S$에서 국소적이다. 따라서 $X$와 $S$가 아핀이라고
가정해도 된다. 그러면 $S=\Spec(R)$로 쓰고 $R$을 유한형
$\mathbf{Z}$-대수들의 여과쌍극한 $R=\colim R_i$로 쓸 수 있다.
극한, 보조정리 \ref{limits-lemma-descend-finite-presentation}을 사용하면
어떤 $i$와, $S$로의 기저변환이 $f$인 사상
$f_i:X_i\to\Spec(R_i)$를 찾을 수 있다. $i$를 충분히 크게 하면 보조정리
\ref{curves-lemma-descend-nodal-family}에 의해 $f_i$가 상대차원 $1$인 기껏해야
마디인 사상이라고 가정할 수 있다. $x$의 상 $x_i\in X_i$는 그 올의
특이점이다. 예를 들어 $\text{Sing}(f)$의 형성이 기저변환과 가환하기
때문이다(인자, 보조정리
\ref{divisors-lemma-base-change-and-fitting-ideal-omega}).
$f_i:X_i\to S_i$와 $x_i$에 대하여 보조정리를 증명할 수 있으면 기저변환을
통해 $f:X\to S$에 대한 결과가 따른다. 따라서 다음 단락에서 다루는
경우로 환원된다.

\medskip\noindent
$S$가 $\mathbf{Z}$ 위 유한형이라고 가정하고, $s\in S$를 $x$의 상이라
하자. $\kappa(x)$는 $\kappa(s)$의 유한 분리확장임을 상기하자. 예를 들어
$\text{Sing}(f)\to S$가 비분기이기 때문이거나, $x$가 올 $X_s$의
마디점이므로 보조정리 \ref{curves-lemma-nodal}을 적용할 수 있기 때문이다.
또 $\kappa'/\kappa(x)$를 주의 \ref{curves-remark-quadratic-extension}에서
$\mathcal{O}_{X_s,x}$에 결부된 차수 $2$의 분리대수라 하자. 스킴의 사상 심화, 보조정리
\ref{more-morphisms-lemma-realize-prescribed-residue-field-extension-etale}에
의해 \'etale 근방 $(V,v)\to(S,s)$를 택하여, 확장
$\kappa(v)/\kappa(s)$가 $\kappa'\cong\kappa(x)\times\kappa(x)$인
경우에는 확장 $\kappa(x)/\kappa(s)$을, $\kappa'$가 체인 경우에는
확장 $\kappa'/\kappa(s)$을 실현하게 할 수 있다. $X$를
$X\times_SV$로, $S$를 $V$로 바꾸면 다음 단락의 상황으로 환원된다.

\medskip\noindent
$S$가 $\mathbf{Z}$ 위 유한형이고 $x\in X_s$가 분할 마디점이라고
가정하자. 정의 \ref{curves-definition-split-node}을 보라. 보조정리
\ref{curves-lemma-formal-local-structure-nodal-family}에 의해 어떤
$h\in\mathfrak m_s^\wedge\subset\mathcal{O}_{S,s}^\wedge$에 대하여
$\mathcal{O}_{S,s}$-대수의 동형
$$
\mathcal{O}_{X,x}^\wedge\cong
\mathcal{O}_{S,s}^\wedge[[s,t]]/(st-h)
$$
이 존재한다. 달리 말해, $a$를 $h$로 보내는 준동형
$$
\sigma:\mathbf{Z}[a]\longrightarrow\mathcal{O}_{S,s}^\wedge
$$
을 생각하면 $\mathcal{O}_{S,s}$-대수의 동형
$$
\mathcal{O}_{X,x}^\wedge
\longrightarrow
\mathcal{O}_{Y_\sigma,y_\sigma}^\wedge
$$
이 존재한다. 여기서
% P07-ERR-0360: preserve the frozen unused variable t in this model.
$$
Y_\sigma=\Spec(\mathbf{Z}[u,v,t]/(uv-a))
\times_{\Spec(\mathbf{Z}[a]),\sigma}\Spec(\mathcal{O}_{S,s}^\wedge)
$$
이고, $y_\sigma$는 $\Spec(\mathcal{O}_{S,s}^\wedge)$의 닫힌점 위에 놓이며
좌표 $u,v$가 영인 $Y_\sigma$의 점이다. 대수학 심화, 명제
\ref{more-algebra-proposition-ubiquity-G-ring}에 의해
$\mathcal{O}_{S,s}$는 G-환이므로 스킴의 사상 심화, 보조정리
\ref{more-morphisms-lemma-relative-map-approximation-pre}를 적용하여 결론을
얻는다.
\end{proof}

\begin{proof}[Artin 근사를 사용하지 않는 증명]
이 증명의 개요만 제시한다. 이 증명은 Mohan Swaminathan이 기여했으며
\cite[명제 X.2.1]{ACG}의 논증을 모방한다.
% P07-ERR-0363: frozen source misspells mimics as mimicks.
앞의 증명과 달리 이 증명은 실제 방정식을 사용한다.

\medskip\noindent
첫째 증명에서와 정확히 같은 방식으로, $S$가 $\mathbf{Z}$ 위 유한형이고
$x\in X_s$가 분할 마디점인 경우로 환원한다. 정의
\ref{curves-definition-split-node}을 보라. 보조정리 \ref{curves-lemma-split-node}에 의해
$\kappa(s)$-대수의 동형
$$
\mathcal{O}_{X_s,x}^\wedge\cong\kappa(s)[[u,v]]/(uv)
$$
이 존재한다. $x$에서 $X_s$의 접공간은 차원이 $2$이다. 따라서
다양체, 보조정리 \ref{varieties-lemma-immersion-into-affine}에 의해 Zariski
축소를 한 뒤 $S$ 위 스킴의 닫힌몰입 $X\to Y$가 존재하고
$Y\to S$가 상대차원 $2$인 매끄러운 사상이라고 가정할 수 있다.
$X$는 (정의에 의해) 상대차원 $1$인 syntomic 사상이므로 $X$는 $Y$의
유효 Cartier 인자이다(인자, 보조정리
\ref{divisors-lemma-immersion-lci-into-smooth-regular-immersion}에서
따른다). 따라서 $Y$와 $X$를 축소한 뒤 $Y=\Spec(B)$가 아핀이고,
$X=\Spec(B/gB)$가 되는 영인자가 아닌 원소 $g\in B$가 있다고 가정할
수 있다.

\medskip\noindent
$S=\Spec(R)$로 쓰고, $s$에 대응하는 소 아이디얼을
% P07-ERR-0364: frozen source has malformed “denote ... be”; translated naturally.
$\mathfrak p\subset R$라 하자. $\mathfrak q\subset B$를 ($Y$의 점으로 본)
$x$에 대응하는 소 아이디얼이라 하자. 사상
$$
B_\mathfrak q/\mathfrak pB_\mathfrak q
\longrightarrow
B_\mathfrak q/gB_\mathfrak q+\mathfrak pB_\mathfrak q
\longrightarrow
\kappa(s)[[u,v]]/(uv)
$$
이 있으며 둘째 화살표는 완비화 뒤 동형이 된다. $B$를 적당한
% P07-ERR-0365: frozen source misspells principal; translated as intended.
주원소 국소화로 바꾼 뒤, $(u,v)^2$을 법으로 각각 $u,v$로 가는
$U,V\in B$가 존재한다고 가정할 수 있다. 그러면
$B_\mathfrak q/\mathfrak pB_\mathfrak q$의 완비화는
$\kappa(s)[[U,V]]$이다. 또한 $Y$를 축소하고 $g$에 단원을 곱한 뒤,
$g$가 $\kappa(s)[[U,V]]$에서 $UV$와 고차항들의 합으로 간다고 가정할
수 있다. $Y$를 더 축소하면 다음을 가정할 수 있다.
\begin{enumerate}
\item $\mathfrak q=\mathfrak pB+(U,V)$,
\item $\Omega_{B/R}$은 $\text{d}U$와 $\text{d}V$를 기저로 하는 자유가군이다.
\end{enumerate}
$g_1,g_2\in B$에 대하여
$\text{d}(g)=g_2\text{d}U+g_1\text{d}V$로 쓰자. 아이디얼
$J=(g_1,g_2)\subset B$는 (고정된 $g$에 대해) ``좌표'' $U,V$의 선택과
무관하다. $\mathfrak pB_\mathfrak q+(U,V)^3B_\mathfrak q$을 법으로
보면 $g_1,g_2$는 $\mathfrak pB_\mathfrak q+(U,V)^2B_\mathfrak q$을
법으로 각각 $U,V$와 합동이다. 따라서 위의 논증에서 $U$를 $g_1$로,
$V$를 $g_2$로 바꾸어 다시 적용할 수 있다. 그러면 $J=(U,V)$를 얻고,
따라서 $g_1,g_2\in(U,V)=J$이다(물론 이들은 앞의 좌표 선택에서의
$g_1$과 $g_2$와 다르다).

\medskip\noindent
$\Spec(B/J)\to S=\Spec(R)$는 $x$에서 \'etale이다. 따라서 $S$를
\'etale 근방으로 바꾸고 $Y$를 다시 축소하여 $R\to B\to B/J$가
동형이라고 가정할 수 있다. $B/J$에서의 상이 $-g$의 상인 원소를
$r\in R$라 하자. 그러면 어떤 $a,b\in R$에 대하여
$$
g+r\equiv aU+bV\bmod(U,V)^2
$$
로 쓸 수 있다. $\text{d}g\in J\Omega_{B/R}$임을 이용하여 미분하면
$a=b=0$을 얻는다. 따라서 어떤 $\alpha,\beta,\gamma\in B$에 대하여
$$
g+r=\alpha U^2+\beta UV+\gamma V^2
$$
로 쓸 수 있다. $\beta$의 $\kappa(x)$에서의 상은 $1$이고
$\alpha,\gamma$의 $\kappa(x)$에서의 상은 $0$이다. 그러므로 $\beta$와
$\alpha+\beta+\gamma$가 $B$에서 가역이라고 가정할 수 있다.
$Y$의 피복 $Y'$을
$$
Y'=\Spec(B[t]/(t(1-t)-\alpha\gamma/\beta^2))
$$
로 두고, $x'$을 $t=1$인 $x$ 위의 유일한 점이라 하자. 사상
$Y'\to Y$는 $x'$에서 \'etale이다. $\Gamma(Y',\mathcal{O}_{Y'})$의 원소
$$
T_1=(\alpha+t\beta)U+((1-t)\beta+\gamma)V
\quad\text{and}\quad
T_2=
\frac{(\alpha+(1-t)\beta)U+(t\beta+\gamma)V}
{\alpha+\beta+\gamma}
$$
를 생각하자. 그러면
$$
g=T_1T_2-r
$$
임을 알 수 있다. $T_1,T_2$도 위의 의미에서 좌표이므로, $T_1,T_2$가
정의하는 사상 $Y'\to\mathbf A^2_S$는 $x'$의 어떤 근방에서 \'etale이고,
$X$의 역상 $X'\subset Y'$은 방정식 $x_1x_2-r=0$으로 주어진
$\mathbf A^2_S=\Spec(R[x_1,x_2])$의 닫힌부분스킴의 역상이다(필요하면
$Y'$을 다시 축소한다). 원하는 도표를 얻기 위해 이 모든 내용을 조합하는
일은 독자에게 맡긴다.
\end{proof}




\section{추가 소멸 결과}
\label{curves-section-more-vanishing}

\noindent
절 \ref{curves-section-vanishing}의 계속이다.

\begin{lemma}
\label{curves-lemma-h1-nonzero-degree-leq-2g-2}
상황 \ref{curves-situation-Cohen-Macaulay-curve}에서 $X$가 종수 $g$인
정역 스킴이라고 가정하자. $\mathcal{L}$을 가역
$\mathcal{O}_X$-가군이라 하자. $Z\subset X$를 아이디얼층
$\mathcal{I}\subset\mathcal{O}_X$를 갖는 $0$차원 닫힌부분스킴이라 하자.
$H^1(X,\mathcal{I}\mathcal{L})$이 영이 아니면
$$
\deg(\mathcal{L})\leq2g-2+\deg(Z)
$$
이고, $\mathcal{I}\mathcal{L}\cong\omega_X$가 아니면 부등식은 엄격하다.
\end{lemma}

\begin{proof}
임의의 곡선, 예를 들어 $X$는 Cohen--Macaulay이다.
$H^1(X,\mathcal{I}\mathcal{L})$이 영이 아니면 보조정리
\ref{curves-lemma-duality-dim-1-CM}에 의해 영이 아닌 사상
$\mathcal{I}\mathcal{L}\to\omega_X$가 존재한다.
$\mathcal{I}\mathcal{L}$이 꼬임 없는 가군이므로 이 사상은 단사이다.
체는 Gorenstein이고 $X$는 환원이므로 $X$의 Gorenstein 자리는 공집합이
아닌 열린집합 $U\subset X$이다. 스킴의 쌍대성, 보조정리
\ref{duality-lemma-gorenstein}을 보라. 같은 보조정리에 의해
$\omega_X|_U$도 가역이다. 따라서 짧은 완전열
$$
0\to\mathcal{I}\mathcal{L}\to\omega_X\to\mathcal{Q}\to0
$$
을 얻으며, $\mathcal{Q}$의 지지는 영차원이다. 그러므로
\begin{align*}
0&\leq\dim\Gamma(X,\mathcal{Q})\\
&=\chi(\mathcal{Q})\\
&=\chi(\omega_X)-\chi(\mathcal{I}\mathcal{L})\\
&=\chi(\omega_X)-\deg(\mathcal{L})-\chi(\mathcal{I})\\
&=2g-2-\deg(\mathcal{L})+\deg(Z)
\end{align*}
이다. 여기서 보조정리 \ref{curves-lemma-euler}와 \ref{curves-lemma-rr}, 식
(\ref{curves-equation-genus}), 그리고 다양체, 보조정리
\ref{varieties-lemma-chi-tensor-finite}와
\ref{varieties-lemma-degree-on-proper-curve}를 썼다. 또한
$\deg(Z)=\dim_k\Gamma(Z,\mathcal{O}_Z)=\chi(\mathcal{O}_Z)$와 짧은 완전열
$0\to\mathcal{I}\to\mathcal{O}_X\to\mathcal{O}_Z\to0$도 사용했다.
이로부터 보조정리가 따른다.
\end{proof}

\begin{lemma}
\label{curves-lemma-degree-more-than-2g-2}
\begin{reference}
\cite[보조정리 2]{Jongmin}
\end{reference}
상황 \ref{curves-situation-Cohen-Macaulay-curve}에서 $X$가 종수 $g$인
정역 스킴이라고 가정하자. $\mathcal{L}$을 가역
$\mathcal{O}_X$-가군이라 하자. $Z\subset X$를 아이디얼층
$\mathcal{I}\subset\mathcal{O}_X$를 갖는 $0$차원 닫힌부분스킴이라 하자.
$\deg(\mathcal{L})>2g-2+\deg(Z)$이면
$H^1(X,\mathcal{I}\mathcal{L})=0$이고, 다음 두 경우 중 하나가 성립한다.
\begin{enumerate}
\item $H^0(X,\mathcal{I}\mathcal{L})\not =0$이다.
\item $g=0$이고 $\deg(\mathcal{L})=\deg(Z)-1$이다.
\end{enumerate}
(2)에서 $Z=\emptyset$이면 $X\cong\mathbf{P}^1_k$이고 $\mathcal{L}$은
$\mathcal{O}_{\mathbf{P}^1}(-1)$에 대응한다.
\end{lemma}

\begin{proof}
$H^1(X,\mathcal{I}\mathcal{L})$의 소멸은 보조정리
\ref{curves-lemma-h1-nonzero-degree-leq-2g-2}에서 따른다.
$H^0(X,\mathcal{I}\mathcal{L})=0$이면
$\chi(\mathcal{I}\mathcal{L})=0$이다. 짧은 완전열
$0\to\mathcal{I}\mathcal{L}\to\mathcal{L}\to\mathcal{O}_Z\to0$에서
$\deg(\mathcal{L})=g-1+\deg(Z)$를 얻는다. 따라서
$g-1+\deg(Z)>2g-2+\deg(Z)$이므로 $g=0$이고 (2)가 성립한다.
(2)에서 $Z=\emptyset$이면 $\mathcal{L}^{-1}$은 차수 $1$인 가역층이다.
보조정리 \ref{curves-lemma-genus-zero-positive-degree}에 의해 이는 동형
$X\to\mathbf{P}^1_k$가 존재하고 $\mathcal{L}^{-1}$이
$\mathcal{O}_{\mathbf{P}^1}(1)$의 역상임을 뜻한다.
\end{proof}

\begin{lemma}
\label{curves-lemma-degree-more-than-2g}
\begin{reference}
\cite[보조정리 3]{Jongmin}
\end{reference}
상황 \ref{curves-situation-Cohen-Macaulay-curve}에서 $X$가 종수 $g$인
정역 스킴이라고 가정하자. $\mathcal{L}$을 가역
$\mathcal{O}_X$-가군이라 하자. $\deg(\mathcal{L})\geq2g$이면
$\mathcal{L}$은 전역 생성된다.
\end{lemma}

\begin{proof}
$Z\subset X$를 $\mathcal{L}$의 전역단면들로 잘라낸 닫힌부분스킴이라
하자. 보조정리 \ref{curves-lemma-degree-more-than-2g-2}에 의해 $Z\not = X$이다.
$\mathcal{I}\subset\mathcal{O}_X$를 $Z$를 잘라내는 아이디얼층이라 하자.
짧은 완전열
$$
0\to\mathcal{I}\mathcal{L}
\to\mathcal{L}\to\mathcal{O}_Z\to0
$$
을 생각하자. $Z\not =\emptyset$이면 코호몰로지 긴 완전열에 의해
$H^1(X,\mathcal{I}\mathcal{L})$은 영이 아니다. 보조정리
\ref{curves-lemma-duality-dim-1-CM}에 의해 영이 아니며 따라서 단사인 사상
$$
\mathcal{I}\mathcal{L}\longrightarrow\omega_X
$$
을 얻는다. 특히 단사 사상
$H^0(X,\mathcal{L})=H^0(X,\mathcal{I}\mathcal{L})
\to H^0(X,\omega_X)$를 얻는다. 그러나
$$
\dim_kH^0(X,\mathcal{L})=\dim_kH^1(X,\mathcal{L})+
\deg(\mathcal{L})+1-g\geq g+1
$$
이고 식 (\ref{curves-equation-genus})에 의해
$\dim H^0(X,\omega_X)=g$이므로 이는 불가능하다.
\end{proof}

\begin{lemma}
\label{curves-lemma-degree-more-than-2g-1-and-Z}
상황 \ref{curves-situation-Cohen-Macaulay-curve}에서 $X$가 종수 $g$인
정역 스킴이라고 가정하자. $\mathcal{L}$을 가역
$\mathcal{O}_X$-가군이라 하자. $Z\subset X$를 공집합이 아닌 $0$차원
닫힌부분스킴이라 하자. $\deg(\mathcal{L})\geq2g-1+\deg(Z)$이면
$\mathcal{L}$은 전역 생성되고 사상
$H^0(X,\mathcal{L})\to H^0(X,\mathcal{L}|_Z)$는 전사이다.
\end{lemma}

\begin{proof}
전역 생성은 보조정리 \ref{curves-lemma-degree-more-than-2g}에서 따른다.
$\mathcal{I}\subset\mathcal{O}_X$가 $Z$의 아이디얼층이면 보조정리
\ref{curves-lemma-h1-nonzero-degree-leq-2g-2}에 의해
$H^1(X,\mathcal{I}\mathcal{L})=0$이다. 따라서 전사성도 따른다.
\end{proof}

\begin{lemma}
\label{curves-lemma-degree-at-least-2g+1}
상황 \ref{curves-situation-Cohen-Macaulay-curve}에서 $X$가 $k$ 위에서
기하적으로 정역이고 종수가 $g$라고 가정하자. $\mathcal{L}$을 가역
$\mathcal{O}_X$-가군이라 하자. $\deg(\mathcal{L})\geq2g+1$이면
$\mathcal{L}$은 매우 풍부하다.
\end{lemma}

\begin{proof}
보조정리 \ref{curves-lemma-degree-more-than-2g}에 의해 $\mathcal{L}$은 전역 생성되고,
따라서 $n=h^0(X,\mathcal{L})-1$일 때 사상
$f:X\to\mathbf{P}^n_k$를 정한다. $\mathcal{L}$이 매우 풍부함을 보이는
것은 $f$가 닫힌몰입임을 보이는 것과 같다. $k$의 대수적 폐포
$\overline{k}$로 $f$를 기저변환한 사상이 닫힌몰입인지 확인하면 충분하다
(내림, 보조정리 \ref{descent-lemma-descending-property-closed-immersion}).
따라서 $k$가 대수적으로 닫혀 있다고 가정해도 된다. 가정에 의해 $X$는
여전히 정역 스킴이다. 보조정리 \ref{curves-lemma-degree-more-than-2g-1-and-Z}에
의해 차수가 2인 모든 $0$차원 닫힌부분스킴 $Z\subset X$에 대하여 제한
사상 $H^0(X,\mathcal{L})\to H^0(X,\mathcal{L}|_Z)$가 전사이다.
다양체, 보조정리
\ref{varieties-lemma-variant-separate-points-tangent-vectors}에 의해
$\mathcal{L}$은 매우 풍부하다.
\end{proof}


\begin{lemma}
\label{curves-lemma-vanishing-on-gorenstein}
\begin{reference}
\cite[보조정리 4]{Jongmin}의 약한 형태
\end{reference}
$k$를 체라 하자. $X$를 $k$ 위의 고유 스킴으로서 환원이고 연결이며
차원이 $1$인 스킴이라 하자. $\mathcal{L}$을 가역
$\mathcal{O}_X$-가군이라 하자. $Z\subset X$를 아이디얼층
$\mathcal{I}\subset\mathcal{O}_X$를 갖는 $0$차원 닫힌부분스킴이라 하자.
$H^1(X,\mathcal{I}\mathcal{L})\not =0$이면, 차원이 $1$인 환원 연결
닫힌부분스킴 $Y\subset X$가 존재하여
$$
\deg(\mathcal{L}|_Y)\leq-2\chi(Y,\mathcal{O}_Y)+\deg(Z\cap Y)
$$
이다. 여기서 $Z\cap Y$는 스킴론적 교집합이다.
\end{lemma}

\begin{proof}
$H^1(X,\mathcal{I}\mathcal{L})$이 영이 아니면 보조정리
\ref{curves-lemma-duality-dim-1-CM}에 의해 영이 아닌 사상
$\varphi:\mathcal{I}\mathcal{L}\to\omega_X$가 존재한다.
$Y\subset X$를 $\varphi$가 $C$의 일반점에서 영이 아닌 $X$의 기약성분 $C$들의
합이라 하자. 그러면 $Y$는 환원 닫힌부분스킴이다.
$\mathcal{J}\subset\mathcal{O}_X$를 $Y$의 아이디얼층이라 하자.
$\mathcal{J}\mathcal{I}\mathcal{L}$은 ($\mathcal{L}$의 부분가군이므로)
매장된 결합점을 갖지 않고, $Y$의 선택과 $X$의 환원성에 의해
$\varphi$는 $\mathcal{J}$의 지지의 일반점들에서 영이다. 따라서
$\varphi$는 다음과 같이 분해된다.
$$
\mathcal{I}\mathcal{L}\to
\mathcal{I}\mathcal{L}/\mathcal{J}\mathcal{I}\mathcal{L}\to\omega_X.
$$
$\mathcal{I}\mathcal{L}/\mathcal{J}\mathcal{I}\mathcal{L}$을 $Y$ 위의
연접층의 직상으로 볼 수 있으며, 표기를 남용하여 그 층도 같은 기호로
나타낸다. 보조정리 \ref{curves-lemma-closed-immersion-dim-1-CM}에 의해
$\omega_Y=\SheafHom(\mathcal{O}_Y,\omega_X)$이므로 $Y$의 일반점들에서 단사인
$\mathcal{O}_Y$-가군의 사상
$$
\mathcal{I}\mathcal{L}/\mathcal{J}\mathcal{I}\mathcal{L}
\to\omega_Y
$$
을 얻는다. $\mathcal{I}'\subset\mathcal{O}_Y$를 $Z\cap Y$의 아이디얼층이라
하자. 사상
$\mathcal{I}\mathcal{L}/\mathcal{J}\mathcal{I}\mathcal{L}
\to\mathcal{I}'\mathcal{L}|_Y$가 있고 그 핵은 닫힌점들에 지지된다.
$\omega_Y$는 Cohen--Macaulay 가군이므로 위 사상은 단사 사상
$\mathcal{I}'\mathcal{L}|_Y\to\omega_Y$를 통해 분해된다. 따라서 $Y$ 위의
연접층들의 완전열
$$
0\to\mathcal{I}'\mathcal{L}|_Y\to\omega_Y\to\mathcal{Q}\to0
$$
을 얻으며, $\mathcal{Q}$는 차원 $0$에 지지된다(여기서는 $\omega_Y$가
$Y$의 일반점들에서 가역가군임을 쓴다). 그러므로
$$
0\leq\dim\Gamma(Y,\mathcal{Q})=
\chi(\mathcal{Q})=\chi(\omega_Y)-\chi(\mathcal{I}'\mathcal{L})=
-2\chi(\mathcal{O}_Y)-\deg(\mathcal{L}|_Y)+\deg(Z\cap Y)
$$
이다. 보조정리 \ref{curves-lemma-euler}와 다양체, 보조정리
\ref{varieties-lemma-chi-tensor-finite}를 썼다. $Y$가 연결이면 보조정리가
증명되었다. 그렇지 않으면 $Y$의 한 연결성분에 대하여 논증의 마지막
부분을 반복한다.
\end{proof}

\begin{lemma}
\label{curves-lemma-global-generation}
$k$를 체라 하자. $X$를 $k$ 위의 고유 스킴으로서 환원이고 연결이며
차원이 $1$인 스킴이라 하자. $\mathcal{L}$을 가역
$\mathcal{O}_X$-가군이라 하자. 차원이 $1$인 모든 환원 연결 닫힌부분스킴
$Y\subset X$에 대하여
$$
\deg(\mathcal{L}|_Y)\geq2\dim_kH^1(Y,\mathcal{O}_Y)
$$
라고 가정하자. 그러면 $\mathcal{L}$은 전역 생성된다.
\end{lemma}

\begin{proof}
$X$의 기약성분 수에 대한 귀납법을 쓴다. $X$가 기약이면 $X$를 체
$k'=H^0(X,\mathcal{O}_X)$ 위의 스킴으로 보아 보조정리
\ref{curves-lemma-degree-more-than-2g}를 적용하면 결론을 얻는다. $X$가 기약이
아니라고 하자. 계속하기 전에 $k$가 유한체이면 $k$를 순수 초월확장
$K$로 바꾼다. 이는 다양체, 보조정리
\ref{varieties-lemma-globally-generated-base-change},
\ref{varieties-lemma-degree-base-change},
\ref{varieties-lemma-geometrically-reduced-any-base-change}, 그리고
\ref{varieties-lemma-bijection-irreducible-components}, 스킴의 코호몰로지, 보조정리 \ref{coherent-lemma-flat-base-change-cohomology}, 보조정리
\ref{curves-lemma-sanity-check-duality}, 그리고 $K$가 $k$ 위에서 기하적으로
정역이라는 초등적 사실에 의해 허용된다.

\medskip\noindent
모순을 얻기 위해 $\mathcal{L}$이 전역 생성되지 않는다고 가정하자.
그러면
$H^0(X,\mathcal{I}\mathcal{L})=H^0(X,\mathcal{L})$이고
$\mathcal{O}_X/\mathcal{I}$가 영이 아니며 차원 $0$인 지지를 갖는 연접
아이디얼층 $\mathcal{I}\subset\mathcal{O}_X$를 택할 수 있다. 예를 들어
$\mathcal{L}$의 전역단면들의 공통 영점자리 안의 임의의 닫힌점의
아이디얼층 $\mathcal{I}$을 택한다. 짧은 완전열
$$
0\to\mathcal{I}\mathcal{L}\to\mathcal{L}\to
\mathcal{L}/\mathcal{I}\mathcal{L}\to0
$$
을 생각하자. $\mathcal{L}/\mathcal{I}\mathcal{L}$의 지지는 차원 $0$이므로
$\mathcal{L}/\mathcal{I}\mathcal{L}$은 전역단면들로 생성된다
(다양체, 보조정리 \ref{varieties-lemma-chi-tensor-finite}). 위 짧은
완전열과 $H^0(X,\mathcal{I}\mathcal{L})=H^0(X,\mathcal{L})$에서 단사 사상
$H^0(X,\mathcal{L}/\mathcal{I}\mathcal{L})
\to H^1(X,\mathcal{I}\mathcal{L})$을 얻는다.

\medskip\noindent
$k$-벡터공간 $H^1(X,\mathcal{I}\mathcal{L})$은
$\Hom(\mathcal{I}\mathcal{L},\omega_X)$의 쌍대임을 상기하자.
$\varphi:\mathcal{I}\mathcal{L}\to\omega_X$를 택하자. 보조정리
\ref{curves-lemma-vanishing-on-gorenstein}에 의해 $H^1(X,\mathcal{L})=0$이다.
따라서
$$
\dim_kH^0(X,\mathcal{I}\mathcal{L})=\dim_kH^0(X,\mathcal{L})=
\deg(\mathcal{L})+\chi(\mathcal{O}_X)>
\dim_kH^1(X,\mathcal{O}_X)=\dim_kH^0(X,\omega_X)
$$
이다. 그러므로 $\varphi$는 전역단면들 위에서 단사가 아니며, 특히
$\varphi$ 자체가 단사가 아니다.
% P07-ERR-0366: source's generic-point definition is grammatically malformed.
$X$의 기약성분의 각 일반점 $\eta\in X$에 대하여,
$V_\eta\subset\Hom(\mathcal{I}\mathcal{L},\omega_X)$를 $\eta$에서 영인
$\varphi$들로 이루어진 $k$-부분벡터공간이라 하자.
$\mathcal{I}\mathcal{L}$의 모든 결합점은 $X$의 일반점이므로 위 논증은
$$
\Hom(\mathcal{I}\mathcal{L},\omega_X)=\bigcup V_\eta
$$
임을 보여 준다. $X$는 유한 개의 일반점을 갖고 $k$는 무한체이므로,
어떤 $\eta$에 대하여
$\Hom(\mathcal{I}\mathcal{L},\omega_X)=V_\eta$이다. $\eta\in C\subset X$를
대응하는 기약성분이라 하자. $Y\subset X$를 $X$의 나머지 기약성분들의
합이라 하자. 그러면 $Y$는 공집합이 아니고 $X$와 다른 환원
닫힌부분스킴이다. $\mathcal{J}\subset\mathcal{O}_X$를 $Y$의 아이디얼층이라
하자. $\mathcal{J}$의 지지가 $C$임을 유의하라.

\medskip\noindent
$\varphi:\mathcal{I}\mathcal{L}\to\omega_X$를 임의로 택하자.
$\mathcal{J}\mathcal{I}\mathcal{L}$은 ($\mathcal{L}$의 부분가군이므로)
매장된 결합점을 갖지 않고, $\varphi$는 $\mathcal{J}$의 지지의 일반점
$\eta$에서 영이므로 $\varphi$는 다음과 같이 분해된다.
$$
\mathcal{I}\mathcal{L}\to
\mathcal{I}\mathcal{L}/\mathcal{J}\mathcal{I}\mathcal{L}\to\omega_X.
$$
$\mathcal{I}\mathcal{L}/\mathcal{J}\mathcal{I}\mathcal{L}$을 $Y$ 위의
연접층의 직상으로 볼 수 있으며, 표기를 남용하여 그 층도 같은 기호로
나타낸다. 보조정리 \ref{curves-lemma-closed-immersion-dim-1-CM}에 의해
$\omega_Y=\SheafHom(\mathcal{O}_Y,\omega_X)$이므로 $\varphi$의 분해
$$
\mathcal{I}\mathcal{L}\to
\mathcal{I}\mathcal{L}/\mathcal{J}\mathcal{I}\mathcal{L}
\xrightarrow{\varphi'}\omega_Y\to\omega_X
$$
를 얻는다. $\mathcal{I}'\subset\mathcal{O}_Y$를
$\mathcal{I}\subset\mathcal{O}_X$의 상이라 하자. 전사 사상
$\mathcal{I}\mathcal{L}/\mathcal{J}\mathcal{I}\mathcal{L}
\to\mathcal{I}'\mathcal{L}|_Y$가 있고 그 핵은 닫힌점들에 지지된다.
$\omega_Y$는 $Y$ 위의 Cohen--Macaulay 가군이므로 $\varphi'$은 사상
$\varphi'':\mathcal{I}'\mathcal{L}|_Y\to\omega_Y$를 통해 분해된다.
따라서 가환도표
$$
\vcenter{
\xymatrix{
0 \ar[r] &
\mathcal{I}\mathcal{L} \ar[r] \ar[d] &
\mathcal{L} \ar[r] \ar[d] &
\mathcal{L}/\mathcal{I}\mathcal{L} \ar[r] \ar[d] &
0 \\
0 \ar[r] &
\mathcal{I}'\mathcal{L}|_Y \ar[r] &
\mathcal{L}|_Y \ar[r] &
\mathcal{L}|_Y/\mathcal{I}'\mathcal{L}|_Y \ar[r] &
0
}
}
\quad\text{and}\quad
\vcenter{
\xymatrix{
\mathcal{I}\mathcal{L} \ar[r]_\varphi \ar[d] & \omega_X \\
\mathcal{I}'\mathcal{L}|_Y \ar[r]^{\varphi''} & \omega_Y \ar[u]
}
}
$$
를 얻는다. 이제 다음과 같이 증명을 끝낼 수 있다. 모든 $\varphi$에
대하여 $\varphi''$가 있고
$\omega_X\in\textit{Coh}(\mathcal{O}_X)$가 함자
$\mathcal{F}\mapsto\Hom_k(H^1(X,\mathcal{F}),k)$를 표현하므로, 사상
$H^1(X,\mathcal{I}\mathcal{L})
\to H^1(Y,\mathcal{I}'\mathcal{L}|_Y)$는 단사이다. 경계사상
$H^0(X,\mathcal{L}/\mathcal{I}\mathcal{L})
\to H^1(X,\mathcal{I}\mathcal{L})$도 단사이므로 합성
$$
H^0(X,\mathcal{L}/\mathcal{I}\mathcal{L})\to
H^0(X,\mathcal{L}|_Y/\mathcal{I}'\mathcal{L}|_Y)\to
H^1(X,\mathcal{I}'\mathcal{L}|_Y)
$$
은 단사이다. 사상
$\mathcal{L}/\mathcal{I}\mathcal{L}\to
\mathcal{L}|_Y/\mathcal{I}'\mathcal{L}|_Y$는 지지가 차원 $0$인 연접가군들의
전사 사상이므로, 첫째 사상
$H^0(X,\mathcal{L}/\mathcal{I}\mathcal{L})
\to H^0(X,\mathcal{L}|_Y/\mathcal{I}'\mathcal{L}|_Y)$는 전사이다(따라서
전단사이다). 그러나 귀납가정에 의해 $\mathcal{L}|_Y$는 전역 생성된다
($Y$가 비연결이어도 물론 성립한다). 따라서 경계사상
$$
H^0(X,\mathcal{L}|_Y/\mathcal{I}'\mathcal{L}|_Y)
\to H^1(X,\mathcal{I}'\mathcal{L}|_Y)
$$
은 단사일 수 없다. 이 모순으로 증명이 끝난다.
\end{proof}





\section{유리 꼬리의 축약}
\label{curves-section-contracting-rational-tails}

\noindent
이 절에서는 표준층의 양성 성질을 개선하기 위해 스킴을 축약하는 가장
간단한 경우를 논한다.

\begin{example}[유리 꼬리의 축약]
\label{curves-example-rational-tail}
$k$를 체라 하자. $X$를 $k$ 위의 차원 $1$인 고유 스킴으로서
$H^0(X,\mathcal{O}_X)=k$인 스킴이라 하자. $X$의 특이점은 기껏해야
마디점이라고 가정하자. {\it 유리 꼬리}란 다음 성질들을 갖는 기약성분
$C\subset X$를 말한다(이를 정역 닫힌부분스킴으로 본다).
\begin{enumerate}
\item $X'\subset X$가 $X\setminus C$의 스킴론적 폐포일 때
$X'\not =\emptyset$이다.
\item 스킴론적 교집합 $C\cap X'$은 하나의 환원점 $x$이다.
\item $H^0(C,\mathcal{O}_C)$는 $x$의 잉여체로 동형사상된다.
\item $C$의 종수는 영이다.
\end{enumerate}
$x$를 지나는 $X$의 기약성분이 적어도 두 개이므로 $x$는 마디점이다.
$k'=H^0(C,\mathcal{O}_C)=\kappa(x)$로 두자. 그러면 $k'/k$는 유한 분리
체확장이다(보조정리 \ref{curves-lemma-nodal}). $X'$ 위에서 항등을 유도하고
표준 사상 $C\to\Spec(k')=x$를 통해 $C$를 $x\in X'$로 보내는 표준 사상
$$
c:X\longrightarrow X'
$$
이 존재한다. $X$는 $C$와 $X'$의 스킴론적 합이고 $X$는 환원이므로, 따라서
이는 스킴의 사상, 보조정리 \ref{morphisms-lemma-scheme-theoretic-union}에서
따른다. 또한
$$
c_*\mathcal{O}_X=\mathcal{O}_{X'}
\quad\text{and}\quad
R^1c_*\mathcal{O}_X=0
$$
이라고 주장한다. 이를 보이기 위해 매장을
$i_C:C\to X$, $i_{X'}:X'\to X$, $i_x:x\to X$라 하고 스킴의 사상,
보조정리 \ref{morphisms-lemma-scheme-theoretic-union}의 완전열
$$
0\to\mathcal{O}_X\to
i_{C,*}\mathcal{O}_C\oplus i_{X',*}\mathcal{O}_{X'}\to
i_{x,*}\kappa(x)\to0
$$
을 사용하자. 고차 직상의 긴 완전열을 보면
$H^0(C,\mathcal{O}_C)=k'$와 $H^1(C,\mathcal{O}_C)=0$을 보이면 충분하며,
이는 가정에서 따른다. $X'$도 $k$ 위의 차원 $1$인 고유 스킴이고,
특이점은 기껏해야 마디점이며(보조정리
\ref{curves-lemma-closed-subscheme-nodal-curve}),
$H^0(X',\mathcal{O}_{X'})=k$이고, $X'$의 종수는 $X$의 종수와 같다.
$c:X\to X'$를 {\it 유리 꼬리의 축약}이라 한다.
\end{example}

\begin{lemma}
\label{curves-lemma-rational-tail-negative}
$k$를 체라 하자. $X$를 $k$ 위의 차원 $1$인 고유 스킴으로서
$H^0(X,\mathcal{O}_X)=k$인 스킴이라 하자. $X$의 특이점은 기껏해야
마디점이라고 가정하자. $C\subset X$가 유리 꼬리이면
(예 \ref{curves-example-rational-tail}), $\deg(\omega_X|_C)<0$이다.
\end{lemma}

\begin{proof}
$X'\subset X$를 위 예에서와 같이 두자. 그러면 짧은 완전열
$$
0\to\omega_C\to\omega_X|_C\to\mathcal{O}_{C\cap X'}\to0
$$
이 있다. 보조정리 \ref{curves-lemma-closed-subscheme-reduced-gorenstein},
\ref{curves-lemma-facts-about-nodal-curves}, 그리고
\ref{curves-lemma-closed-subscheme-nodal-curve}를 보라. 위 예의 $k'$을 쓰면,
명제 \ref{curves-proposition-projective-line}에 의해
$C\cong\mathbf{P}^1_{k'}$이므로 $\deg(\omega_C)=-2[k':k]$이고
$\deg(C\cap X')=[k':k]$이다. 따라서
$\deg(\omega_X|_C)=-[k':k]$이며 이는 음수이다.
\end{proof}

\begin{lemma}
\label{curves-lemma-rational-tail-field-extension}
$k$를 체라 하자. $X$를 $k$ 위의 차원 $1$인 고유 스킴으로서
$H^0(X,\mathcal{O}_X)=k$인 스킴이라 하자. $X$의 특이점은 기껏해야
마디점이라고 가정하자. $C\subset X$가 유리 꼬리라고 하자
(예 \ref{curves-example-rational-tail}). 임의의 체확장 $K/k$에 대하여
기저변환 $C_K\subset X_K$는 유한 개 유리 꼬리의 서로소 합이다.
\end{lemma}

\begin{proof}
$x\in C$와 $k'=\kappa(x)$를 위 예에서와 같이 두자. 명제
\ref{curves-proposition-projective-line}에 의해
$C\cong\mathbf{P}^1_{k'}$임을 유의하라. $k'/k$가 유한 분리이므로
$$
k'\otimes_kK=K'_1\times\ldots\times K'_n
$$
은 유한 분리확장 $K'_i/K$들의 유한 곱이다.
$C_i=\mathbf{P}^1_{K'_i}$로 두고 $x_i\in C_i$를 $x$의 역상이라 하자.
그러면 원하는 대로 $C_K=\coprod C_i$이고 $X'_K\cap C_i=x_i$이다.
\end{proof}

\begin{lemma}
\label{curves-lemma-no-rational-tail}
$k$를 체라 하자. $X$를 $k$ 위의 차원 $1$인 고유 스킴으로서
$H^0(X,\mathcal{O}_X)=k$인 스킴이라 하자. $X$의 특이점은 기껏해야
마디점이라고 가정하자. $X$가 유리 꼬리를 갖지 않으면
(예 \ref{curves-example-rational-tail}), 차원이 $1$인 모든 환원 연결
닫힌부분스킴 $Y\subset X$, $Y\not = X$에 대하여
$\deg(\omega_X|_Y)\geq\dim_kH^1(Y,\mathcal{O}_Y)$이다.
\end{lemma}

\begin{proof}
$Y\subset X$를 보조정리의 명제에서와 같이 두자. 그러면
$k'=H^0(Y,\mathcal{O}_Y)$는 체이고 $k$의 유한확장이며, $[k':k]$는 아래에서
$Y$와 $Y$ 위의 연접층에 결부된 모든 수치 불변량을 나눈다. 다양체,
보조정리 \ref{varieties-lemma-divisible}를 보라.
$Z\subset X$를 보조정리 \ref{curves-lemma-closed-subscheme-reduced-gorenstein}에서와
같이 두자. 이 보조정리와 보조정리 \ref{curves-lemma-facts-about-nodal-curves} 및
\ref{curves-lemma-closed-subscheme-nodal-curve}의 결과를 더 언급하지 않고
사용하겠다. 그러면 짧은 완전열
$$
0\to\omega_Y\to\omega_X|_Y\to\mathcal{O}_{Y\cap Z}\to0
$$
을 얻는다. 보조정리 \ref{curves-lemma-closed-subscheme-reduced-gorenstein}을 보라.
따라서
$$
\deg(\omega_X|_Y)=\deg(Y\cap Z)+\deg(\omega_Y)=
\deg(Y\cap Z)-2\chi(Y,\mathcal{O}_Y)
$$
이다. 그러므로 보조정리가 거짓이면
$$
2[k':k]>\deg(Y\cap Z)+\dim_kH^1(Y,\mathcal{O}_Y)
$$
이다. $Y\cap Z$는 공집합이 아니고, 위에서 언급한 나눗셈 성질에 의해
% P07-ERR-0368: frozen source misspells divisibility; translated as intended.
이는 $Y\cap Z$가 $Y$의 매끄러운 자리의 하나의 $k'$-유리점이고
$H^1(Y,\mathcal{O}_Y)=0$일 때만 가능하다. $Y$가 기약이면 이는 $Y$가
유리 꼬리임을 뜻한다. $Y$가 가약이면
$\deg(\omega_X|_Y)=-[k':k]$이므로, 다양체, 보조정리
\ref{varieties-lemma-degree-in-terms-of-components}에 의해
$\deg(\omega_X|_C)<0$인 $Y$의 어떤 기약성분 $C$가 존재한다.
위 분석을 $C$에 적용하면 $C$가 유리 꼬리임을 알 수 있다.
\end{proof}

\begin{lemma}
\label{curves-lemma-no-rational-tail-semiample-genus-geq-2}
$k$를 체라 하자. $X$를 $k$ 위의 차원 $1$인 고유 스킴으로서
$H^0(X,\mathcal{O}_X)=k$인 스킴이라 하자. $X$의 특이점은 기껏해야
마디점이라고 가정하자. $X$가 유리 꼬리를 갖지 않는다고 하자
(예 \ref{curves-example-rational-tail}). 다음이 성립한다.
\begin{enumerate}
\item $X$의 종수가 $0$이면 $X$는 기약 평면 원뿔곡선과 동형이고
$\omega_X^{\otimes-1}$은 매우 풍부하다.
\item $X$의 종수가 $1$이면 $\omega_X\cong\mathcal{O}_X$이다.
\item $X$의 종수가 $\geq2$이면 $m\geq2$에 대하여
$\omega_X^{\otimes m}$은 전역 생성된다.
\end{enumerate}
\end{lemma}

\begin{proof}
보조정리 \ref{curves-lemma-facts-about-nodal-curves}에 의해 $X$는 Gorenstein이다.
즉, $\omega_X$는 가역 $\mathcal{O}_X$-가군이다.

\medskip\noindent
$X$의 종수가 영이면 $\deg(\omega_X)<0$이다. 따라서 $X$가 둘 이상의
기약성분을 가지면 보조정리 \ref{curves-lemma-no-rational-tail}과 모순이다.
기약인 경우 보조정리 \ref{curves-lemma-genus-zero}에 의해 $X$는 기약 평면
원뿔곡선과 동형이고 $\omega_X^{\otimes-1}$은 매우 풍부하다.

\medskip\noindent
$X$의 종수가 $1$이면 $\omega_X$는 전역단면을 가지며, 모든 기약성분에
대하여 $\deg(\omega_X|_C)=0$이다. 실제로 보조정리
\ref{curves-lemma-no-rational-tail}에 의해 모든 기약성분 $C$에 대하여
$\deg(\omega_X|_C)\geq0$이고, 보조정리 \ref{curves-lemma-genus-gorenstein}에 의해
이 수들의 합은 $0$이며, 다양체, 보조정리
\ref{varieties-lemma-degree-in-terms-of-components}를 적용할 수 있다.
그러면 다양체, 보조정리
\ref{varieties-lemma-no-sections-dual-nef}에 의해
$\omega_X\cong\mathcal{O}_X$이다.

\medskip\noindent
$X$의 종수 $g$가 $2$ 이상이라고 가정하자. $X$가 기약이면 보조정리
\ref{curves-lemma-degree-more-than-2g}에 의해 끝난다. $X$가 가약이라고 하자.
보조정리 \ref{curves-lemma-no-rational-tail}에 의해 보조정리
\ref{curves-lemma-global-generation}의 부등식은 명제에서와 같은 모든
$Y\subset X$에 대하여 성립하며, 예외는 $Y=X$뿐이다. 보조정리
\ref{curves-lemma-global-generation}의 증명을 분석하면, 가약인 경우
% P07-ERR-0369: source has singular/plural agreement; two inequalities follow.
$Y=X$에 대하여 사용되는 부등식은 다음 두 개뿐이다.
$$
\deg(\omega_X^{\otimes m})>-2\chi(\mathcal{O}_X)
\quad\text{and}\quad
\deg(\omega_X^{\otimes m})+\chi(\mathcal{O}_X)>
\dim_kH^1(X,\mathcal{O}_X).
$$
$g\geq2$이고 $m\geq2$이면 둘 다 성립하므로 결론을 얻는다.
\end{proof}

\begin{lemma}
\label{curves-lemma-contracting-rational-tails}
$k$를 체라 하자. $X$를 $k$ 위의 차원 $1$인 고유 스킴으로서
$H^0(X,\mathcal{O}_X)=k$인 스킴이라 하자. $X$의 특이점은 기껏해야
마디점이라고 가정하자. 유리 꼬리가 없어질 때까지 유리 꼬리를 축약한
사상들의 열
$$
X=X_0\to X_1\to\ldots\to X_n=X'
$$
을 생각하자(예 \ref{curves-example-rational-tail}). 그러면 다음이 성립한다.
\begin{enumerate}
\item $X$의 종수가 $0$이면 $X'$은 기약 평면 원뿔곡선이다.
% P07-ERR-0367: preserve frozen O_X; expected target is O_{X'}.
\item $X$의 종수가 $1$이면 $\omega_{X'}\cong\mathcal{O}_X$이다.
\item $X$의 종수가 $>1$이면 $m\geq2$에 대하여
$\omega_{X'}^{\otimes m}$은 전역 생성된다.
\end{enumerate}
$X$의 종수가 $\geq1$이면 사상 $X\to X'$은 선택과 무관하고, 이 사상의
형성은 기저체 확장과 가환한다.
\end{lemma}

\begin{proof}
유리 꼬리가 없어질 때까지 이를 축약한다. 그러면 보조정리
\ref{curves-lemma-no-rational-tail-semiample-genus-geq-2}에 의해 (1), (2), (3)이
성립한다.

\medskip\noindent
유일성. $f:X\to X'$이 한 선택과 무관함을 보이려면, 임의의 유리 꼬리
$C\subset X$가 $X\to X'$에 의해 한 점으로 간다는 것을 보이면 충분하다.
일부 세부사항은 생략한다. 그렇지 않으면 기약성분 $f(C)$의 일반점에서
영이 되지 않는 단면 $s\in\Gamma(X',\omega_{X'}^{\otimes2})$를 찾을 수
있다. 각 축약 $X_i\to X_{i+1}$에 단면 $X_{i+1}\to X_i$가 있으므로
$f$의 단면 $X'\to X$가 있다. 그러면 완전열
$$
0\to\omega_{X'}\to\omega_X\to\omega_X|_{X''}\to0
$$
을 얻는다. 여기서 $X''\subset X$는 $f$에 의해 축약되는 기약성분들의
합이다. 보조정리 \ref{curves-lemma-closed-subscheme-reduced-gorenstein}을 보라.
따라서 사상 $\omega_{X'}^{\otimes2}\to\omega_X^{\otimes2}$를 얻으며,
$s$의 상을 취하면 $C$의 일반점에서 영이 되지 않는
$\omega_X^{\otimes2}$의 단면을 얻는다. 이는 $\omega_X$의 유리 꼬리로의
제한이 음의 차수를 갖는다는 사실과 모순이다
(보조정리 \ref{curves-lemma-rational-tail-negative}).

\medskip\noindent
기저체 확장에 관한 명제는 보조정리
\ref{curves-lemma-rational-tail-field-extension}에서 따른다. 일부 세부사항은
생략한다.
\end{proof}





\section{유리 다리의 축약}
\label{curves-section-contracting-rational-bridges}

\noindent
이 절에서는 표준층의 양성 성질을 개선하기 위해 스킴을 축약하는 그 다음
간단한 경우(절 \ref{curves-section-contracting-rational-tails}에서 논한
경우 다음)를 논한다.

\begin{example}[유리 다리의 축약]
\label{curves-example-rational-bridge}
$k$를 체라 하자. $X$를 $k$ 위의 차원 $1$인 고유 스킴으로서
$H^0(X,\mathcal{O}_X)=k$인 스킴이라 하자. $X$의 특이점은 기껏해야
마디점이라고 가정하자. {\it 유리 다리}란 다음 성질들을 갖는 기약성분
$C\subset X$를 말한다(이를 정역 닫힌부분스킴으로 본다).
\begin{enumerate}
\item $X'\subset X$가 $X\setminus C$의 스킴론적 폐포일 때
$X'\not =\emptyset$이다.
\item 스킴론적 교집합 $C\cap X'$의
$H^0(C,\mathcal{O}_C)$ 위의 차수는 $2$이다.
\item $C$의 종수는 영이다.
\end{enumerate}
$k'=H^0(C,\mathcal{O}_C)$ 및
$k''=H^0(C\cap X',\mathcal{O}_{C\cap X'})$로 두자. 그러면 $k'$은 체이고
(다양체, 보조정리
\ref{varieties-lemma-proper-geometrically-reduced-global-sections}),
$\dim_{k'}(k'')=2$이다. $C\cap X'$의 각 점을 지나는 $X$의 기약성분이
적어도 두 개이므로, 이 점들은 $X$의 마디점이고 $C$와 $X'$ 모두의
매끄러운 점이다(보조정리 \ref{curves-lemma-closed-subscheme-nodal-curve}).
따라서 $k'/k$는 유한 분리 체확장이고, $k''/k'$은 차수 $2$인 분리
체확장이거나 $k''=k'\times k'$이다(보조정리 \ref{curves-lemma-nodal}).
절 \ref{curves-section-pushouts}에 의해 다음 밀어붙임이 존재한다.
$$
\xymatrix{
C \cap X' \ar[r] \ar[d] &
X' \ar[d]^a \\
\Spec(k') \ar[r] &
Y
}
$$
이 밀어붙임은 여러 좋은 성질을 가지며, 아래에서 이를 더 언급하지 않고
모두 사용한다. $y\in Y$를 $\Spec(k')\to Y$의 상이라 하자. 그러면
$C\cap X'$이 $2$점인지 $1$점인지에 따라 각각
$$
\mathcal{O}_{Y,y}^\wedge\cong k'[[s,t]]/(st)
\quad\text{or}\quad
\mathcal{O}_{Y,y}^\wedge\cong
\{f\in k''[[s]]:f(0)\in k'\}
$$
이다. 이는 보조정리 \ref{curves-lemma-complete-local-ring-pushout}와, 대수학 심화, 보조정리 \ref{more-algebra-lemma-power-series-over-residue-field}에
의해 $p\in C\cap X'$에 대하여
$\mathcal{O}_{X',p}\cong\kappa(p)[[t]]$라는 사실에서 따른다. 따라서
$y\in Y$는 마디점이다. 보조정리 \ref{curves-lemma-nodal}와
\ref{curves-lemma-2-branches-delta-1}, 특히 보조정리
\ref{curves-lemma-2-branches-delta-1}의 (2) $\Rightarrow$ (1)의 증명에서 경우 II의
논의를 보라. 그러므로 $Y$의 특이점은 기껏해야 마디점이다.

\medskip\noindent
위 가환도표를 다음 도표로 확장할 수 있다.
$$
\xymatrix{
C \cap X' \ar[r] \ar[d] &
X' \ar[d]^a \ar[r] & X \ar[ld]^c &
C \ar[ld] \ar[l] \\
\Spec(k') \ar[r] &
Y &
\Spec(k') \ar[l]
}
$$
여기서 아래의 두 수평 화살표는 같다. 실제로 $X$는 $X'$과 $C$의
스킴론적 합이므로 밀어붙임이고(스킴의 사상, 보조정리
\ref{morphisms-lemma-scheme-theoretic-union}), 사상 $C\to Y$와
$X'\to Y$는 $C\cap X'$에서 일치한다. 마지막으로
$$
c_*\mathcal{O}_X=\mathcal{O}_Y
\quad\text{and}\quad
R^1c_*\mathcal{O}_X=0
$$
이라고 주장한다. 이를 보이기 위해 스킴의 사상, 보조정리
\ref{morphisms-lemma-scheme-theoretic-union}의 완전열
$$
0\to\mathcal{O}_X\to
\mathcal{O}_C\oplus\mathcal{O}_{X'}\to
\mathcal{O}_{C\cap X'}\to0
$$
을 사용하자. 고차 직상의 긴 완전열은
$$
0\to c_*\mathcal{O}_X\to c_*\mathcal{O}_C\oplus c_*\mathcal{O}_{X'}\to
c_*\mathcal{O}_{C\cap X'}\to
R^1c_*\mathcal{O}_X\to
R^1c_*\mathcal{O}_C\oplus R^1c_*\mathcal{O}_{X'}
$$
이다. $c|_{X'}=a$는 아핀이므로 $R^1c_*\mathcal{O}_{X'}=0$이다.
$c|_C$는 $C\to\Spec(k')\to X$로 분해되고 $C$의 종수는 영이므로
$R^1c_*\mathcal{O}_C=0$이다. $\mathcal{O}_{X'}\to\mathcal{O}_{C\cap X'}$가
전사이고 $c|_{X'}$가 아핀이므로
$c_*\mathcal{O}_{X'}\to c_*\mathcal{O}_{C\cap X'}$는 전사이다. 이로부터
$R^1c_*\mathcal{O}_X=0$을 얻는다. 마지막으로 위 완전열과 스킴의 사상 심화, 명제
\ref{more-morphisms-proposition-pushout-along-closed-immersion-and-integral}의
밀어붙임 구조층에 대한 기술로부터
$\mathcal{O}_Y=c_*\mathcal{O}_X$를 얻는다.

\medskip\noindent
이 모든 사실은 $Y$도 $k$ 위의 차원 $1$인 고유 스킴이고
$H^0(Y,\mathcal{O}_Y)=k$이며, 그 특이점은 기껏해야 마디점이고
(보조정리 \ref{curves-lemma-closed-subscheme-nodal-curve}), $Y$의 종수는 $X$의 종수와
같음을 뜻한다. $c:X\to Y$를 {\it 유리 다리의 축약}이라 한다.
\end{example}

\begin{lemma}
\label{curves-lemma-rational-bridge-zero}
$k$를 체라 하자. $X$를 $k$ 위의 차원 $1$인 고유 스킴으로서
$H^0(X,\mathcal{O}_X)=k$인 스킴이라 하자. $X$의 특이점은 기껏해야
마디점이라고 가정하자. $C\subset X$가 유리 다리이면
(예 \ref{curves-example-rational-bridge}), $\deg(\omega_X|_C)=0$이다.
\end{lemma}

\begin{proof}
$X'\subset X$를 위 예에서와 같이 두자. 그러면 짧은 완전열
$$
0\to\omega_C\to\omega_X|_C\to\mathcal{O}_{C\cap X'}\to0
$$
이 있다. 보조정리 \ref{curves-lemma-closed-subscheme-reduced-gorenstein},
\ref{curves-lemma-facts-about-nodal-curves}, 그리고
\ref{curves-lemma-closed-subscheme-nodal-curve}를 보라. 위 예의 $k''/k'/k$를
쓰면, $C$의 종수가 $0$이므로(보조정리 \ref{curves-lemma-rr})
$\deg(\omega_C)=-2[k':k]$이고
$\deg(C\cap X')=[k'':k]=2[k':k]$이다. 따라서
$\deg(\omega_X|_C)=0$이다.
\end{proof}

\begin{lemma}
\label{curves-lemma-rational-bridge-field-extension}
$k$를 체라 하자. $X$를 $k$ 위의 차원 $1$인 고유 스킴으로서
$H^0(X,\mathcal{O}_X)=k$인 스킴이라 하자. $X$의 특이점은 기껏해야
마디점이라고 가정하자. $C\subset X$가 유리 다리라고 하자
(예 \ref{curves-example-rational-bridge}). 임의의 체확장 $K/k$에 대하여
기저변환 $C_K\subset X_K$는 유한 개 유리 다리의 서로소 합이다.
\end{lemma}

\begin{proof}
$k''/k'/k$를 위 예에서와 같이 두자. $k'/k$가 유한 분리이므로
$$
k'\otimes_kK=K'_1\times\ldots\times K'_n
$$
은 유한 분리확장 $K'_i/K$들의 유한 곱이다. 이에 대응하는 곱분해
$k''\otimes_kK=\prod K''_i$는 차수 $2$인 분리대수 확장
$K''_i/K'_i$를 준다. $C_i=C_{K'_i}$로 두자. 그러면
$C_K=\coprod C_i$이고, 평탄 기저변환에 의해
$H^1(C_K,\mathcal{O}_{C_K})=0$이므로 각 $C_i$는 ($K'_i$ 위의 곡선으로
보아) 종수가 $0$이다. 마지막으로 원하는 대로
$X'_K\cap C_i=\Spec(K''_i)$는 $K'_i$ 위 차수 $2$이다.
\end{proof}

\begin{lemma}
\label{curves-lemma-rational-bridge-canonical}
$c:X\to Y$를 유리 다리의 축약이라 하자
(예 \ref{curves-example-rational-bridge}). 그러면
$c^*\omega_Y\cong\omega_X$이다.
\end{lemma}

\begin{proof}
직접 계산으로 증명할 수 있지만, 우리는 $\omega_X$가 함자
$\textit{Coh}(\mathcal{O}_X)\to\textit{Sets}$,
$$
\mathcal{F}\mapsto\Hom_k(H^1(X,\mathcal{F}),k)
=H^1(X,\mathcal{F})^\vee
$$
를 표현하는 연접 $\mathcal{O}_X$-가군이라는 특성화를 사용한다.
보조정리 \ref{curves-lemma-duality-dim-1-CM} 또는 스킴의 쌍대성, 보조정리
\ref{duality-lemma-dualizing-module-proper-over-A}를 보라.

\medskip\noindent
정확히 말해, $\mathcal{C}_Y$를 대상이 가역 $\mathcal{O}_Y$-가군이고
사상이 $\mathcal{O}_Y$-가군 준동형인 범주라 하자. $\mathcal{C}_X$를
대상이 $\mathcal{L}|_C\cong\mathcal{O}_C$인 가역
$\mathcal{O}_X$-가군 $\mathcal{L}$이고 사상이 $\mathcal{O}_Y$-가군
준동형인 범주라 하자. 함자
$$
c^*:\mathcal{C}_Y\to\mathcal{C}_X
$$
가 범주의 동치라고 주장한다. 실제로 스킴의 사상 심화, 보조정리
\ref{more-morphisms-lemma-bijection-on-Pic}에 의해 본질적으로 전사이다.
그러면 사영공식(코호몰로지, 보조정리
\ref{cohomology-lemma-projection-formula})에 의해
$c_*c^*\mathcal N=\mathcal N$이고, 따라서 $c^*$은 $c_*$를 준역으로 갖는
동치이다.

\medskip\noindent
$\omega_X$가 $\mathcal{C}_X$의 대상이라고 주장한다. 실제로 짧은 완전열
$$
0\to\omega_C\to\omega_X|_C\to\mathcal{O}_{C\cap X'}\to0
$$
이 있다. 보조정리 \ref{curves-lemma-closed-subscheme-reduced-gorenstein}을 보라.
차수를 취하면 $\deg(\omega_X|_C)=0$을 얻는다(작은 세부사항은 생략한다).
따라서 보조정리 \ref{curves-lemma-genus-zero-pic}에 의해 $\omega_X|_C$는 자명하고
$\omega_X$는 $\mathcal{C}_X$의 대상이다.

\medskip\noindent
$R^1c_*\mathcal{O}_X=0$이므로 사영공식은
$\mathcal N\in\Ob(\mathcal{C}_Y)$에 대하여
$R^1c_*c^*\mathcal N=0$임을 보여 준다. 따라서 Leray 스펙트럼열
(코호몰로지, 보조정리 \ref{cohomology-lemma-apply-Leray})에 의해
% P07-ERR-0372: frozen source omits the governing verb before this diagram.
다음 범주와 함자의 도표는 가환한다.
$$
\xymatrix{
\mathcal{C}_Y \ar[rr]_{c^*} \ar[dr]_{H^1(Y, -)^\vee} & &
\mathcal{C}_X \ar[ld]^{H^1(X, -)^\vee} \\
& \textit{Sets}
}
$$
$\omega_Y\in\Ob(\mathcal{C}_Y)$는 남동쪽 화살표를 표현하고
$\omega_X\in\Ob(\mathcal{C}_X)$도 남동쪽 화살표를 표현하므로 Yoneda
보조정리(범주, 보조정리 \ref{categories-lemma-yoneda})에 의해 결론을
얻는다.
\end{proof}

\begin{lemma}
\label{curves-lemma-no-rational-bridge-ample-genus-geq-2}
$k$를 체라 하자. $X$를 $k$ 위의 차원 $1$인 고유 스킴으로서
$H^0(X,\mathcal{O}_X)=k$인 스킴이라 하자. 다음을 가정하자.
\begin{enumerate}
\item $X$의 특이점은 기껏해야 마디점이다.
\item $X$는 유리 꼬리를 갖지 않는다
(예 \ref{curves-example-rational-tail}).
\item $X$는 유리 다리를 갖지 않는다
(예 \ref{curves-example-rational-bridge}).
\item $X$의 종수 $g$는 $\geq2$이다.
\end{enumerate}
그러면 $\omega_X$는 풍부하다.
\end{lemma}

\begin{proof}
$X$의 모든 기약성분 $C$에 대하여 $\deg(\omega_X|_C)>0$임을 보이면
충분하다. 다양체, 보조정리
\ref{varieties-lemma-ampleness-in-terms-of-degrees-components}를 보라.
$X=C$가 기약이면 $g\geq2$와 보조정리 \ref{curves-lemma-genus-gorenstein}에서
따른다. 그렇지 않으면 $k'=H^0(C,\mathcal{O}_C)$로 두자. 이는 체이고
$k$의 유한확장이며, $[k':k]$는 아래에서 $C$와 $C$ 위의 연접층에
결부된 모든 수치 불변량을 나눈다. 다양체, 보조정리
\ref{varieties-lemma-divisible}를 보라. $X'\subset X$를 보조정리
\ref{curves-lemma-closed-subscheme-reduced-gorenstein}에서와 같이
$X\setminus C$의 폐포라 하자. 이 보조정리와 보조정리
\ref{curves-lemma-facts-about-nodal-curves} 및
\ref{curves-lemma-closed-subscheme-nodal-curve}의 결과를 더 언급하지 않고
사용하겠다. 그러면 짧은 완전열
$$
0\to\omega_C\to\omega_X|_C\to\mathcal{O}_{C\cap X'}\to0
$$
을 얻는다. 보조정리 \ref{curves-lemma-closed-subscheme-reduced-gorenstein}을 보라.
따라서
$$
\deg(\omega_X|_C)=\deg(C\cap X')+\deg(\omega_C)=
\deg(C\cap X')-2\chi(C,\mathcal{O}_C)
$$
이다. 그러므로 보조정리가 거짓이면
$$
2[k':k]\geq\deg(C\cap X')+2\dim_kH^1(C,\mathcal{O}_C)
$$
이다. $C\cap X'$은 공집합이 아니고, 위에서 언급한 나눗셈 성질에 의해
% P07-ERR-0370: frozen source misspells divisibility; translated as intended.
이는 다음 두 경우 중 하나일 때만 가능하다.
\begin{enumerate}
\item[(a)] $C\cap X'$은 $C$의 하나의 $k'$-유리점이고
$H^1(C,\mathcal{O}_C)=0$이다.
\item[(b)] $C\cap X'$은 $k'$ 위 차수 $2$이고
$H^1(C,\mathcal{O}_C)=0$이다.
\end{enumerate}
첫째 가능성은 $C$가 유리 꼬리임을 뜻하고, 둘째는 $C$가 유리 다리임을
뜻한다. 둘 다 배제했으므로 증명이 끝난다.
\end{proof}

\begin{lemma}
\label{curves-lemma-contracting-rational-bridges}
$k$를 체라 하자. $X$를 $k$ 위의 차원 $1$인 고유 스킴으로서
$H^0(X,\mathcal{O}_X)=k$이고 종수 $g\geq2$인 스킴이라 하자.
$X$의 특이점은 기껏해야 마디점이고 $X$는 유리 꼬리를 갖지 않는다고
가정하자. 유리 다리가 없어질 때까지 유리 다리를 축약한 사상들의 열
$$
X=X_0\to X_1\to\ldots\to X_n=X'
$$
을 생각하자(예 \ref{curves-example-rational-bridge}).
% P07-ERR-0371: frozen source omits the copula; translated as an assertion.
그러면 $\omega_{X'}$은 풍부하다. 사상 $X\to X'$은 선택과 무관하고,
이 사상의 형성은 기저체 확장과 가환한다.
\end{lemma}

\begin{proof}
유리 다리가 없어질 때까지 이를 축약한다. 그러면 보조정리
\ref{curves-lemma-no-rational-bridge-ample-genus-geq-2}에 의해
$\omega_{X'}$은 풍부하다.

\medskip\noindent
$f:X\to X'$을 그 합성이라 하자. 보조정리
\ref{curves-lemma-rational-bridge-canonical}과 귀납법에 의해
$f^*\omega_{X'}=\omega_X$이다. 유리 다리의 축약에 대하여 성립하므로
$f_*\mathcal{O}_X=\mathcal{O}_{X'}$이다. 따라서 사영공식은 모든 가역
$\mathcal{O}_{X'}$-가군 $\mathcal{L}$에 대하여
$f_*f^*\mathcal{L}=\mathcal{L}$임을 말한다. 그러므로 모든 $m$에 대하여
$$
\Gamma(X',\omega_{X'}^{\otimes m})=
\Gamma(X,\omega_X^{\otimes m})
$$
이다. $X'$은 이 공간들의 직합의 Proj이므로, 스킴의 사상, 보조정리
% P07-ERR-0373: citation attachment in frozen source is malformed.
\ref{morphisms-lemma-proper-ample-is-proj}에 의해 사상 $X\to X'$은 완전히
표준적임을 알 수 있다.

\medskip\noindent
% P07-ERR-0374: frozen source has a comma splice; translated as two clauses.
$K/k$를 체의 확장이라 하자. 그러면 $\omega_{X_K}$는 $\omega_X$의
역상이고(보조정리 \ref{curves-lemma-sanity-check-duality}), 스킴의 코호몰로지,
보조정리 \ref{coherent-lemma-flat-base-change-cohomology}에 의해
$\Gamma(X,\omega_X^{\otimes m})\otimes_kK$는
$\Gamma(X_K,\omega_{X_K}^{\otimes m})$와 같다. 따라서 위 논증에 의해
$f:X\to X'$의 형성은 $K/k$에 의한 기저변환과 가환한다.
일부 세부사항은 생략한다.
\end{proof}






\section{안정곡선으로의 축약}
\label{curves-section-contracting-to-stable}

\noindent
이 절에서는 절 \ref{curves-section-contracting-rational-tails}와
\ref{curves-section-contracting-rational-bridges}에서 찾은 축약 사상들을
결합한다. 즉, $k$를 체라 하고 $X$를 $k$ 위의 차원 $1$인 고유
스킴으로서 $H^0(X,\mathcal{O}_X)=k$이고 종수 $g\geq2$인 스킴이라 하자.
$X$의 특이점은 기껏해야 마디점이라고 가정하자. 보조정리
\ref{curves-lemma-contracting-rational-tails}의 사상과 보조정리
\ref{curves-lemma-contracting-rational-bridges}의 사상을 합성하면 사상
$$
c:X\longrightarrow Y
$$
를 얻는다. 여기서 $Y$ 역시 $k$ 위의 차원 $1$인 고유 스킴이고,
그 특이점은 기껏해야 마디점이며, $k=H^0(Y,\mathcal{O}_Y)$이고 종수는
$g$이며, $\mathcal{O}_Y=c_*\mathcal{O}_X$와
$R^1c_*\mathcal{O}_X=0$이 성립하고, $\omega_Y$는 $Y$ 위에서 풍부하다.
보조정리 \ref{curves-lemma-characterize-contraction-to-stable}는 실제로 이 조건들이
이 사상을 특성화함을 보여 준다.

\begin{lemma}
\label{curves-lemma-contract-gorenstein-canonical}
$k$를 체라 하고 $c:X\to Y$를 $k$ 위 고유 스킴의 사상이라 하자.
다음을 가정하자.
\begin{enumerate}
\item $\mathcal{O}_Y=c_*\mathcal{O}_X$이고 $R^1c_*\mathcal{O}_X=0$이다.
\item $X$와 $Y$는 환원이고 Gorenstein이며 차원이 $1$이다.
\item $\exists\ m \in \mathbf{Z}$에 대하여
$H^1(X,\omega_X^{\otimes m})=0$이고
$\omega_X^{\otimes m}$은 전역단면들로 생성된다.
\end{enumerate}
그러면 $c^*\omega_Y\cong\omega_X$이다.
\end{lemma}

\begin{proof}
스킴의 사상 심화, 정리
\ref{more-morphisms-theorem-stein-factorization-Noetherian}에 의해 $c$의
올들은 기하적으로 연결이다. 특히 $c$는 전사이다. $X_y$의 차원이
$1$인 $Y$의 닫힌점은 유한 개 $y=y_1,\ldots,y_r$이고,
$Y\setminus\{y_1,\ldots,y_r\}$ 위에서 $c$는 동형이다. 일부 세부사항은
생략한다. 힌트: $\{y_1,\ldots,y_r\}$ 밖에서 $c$는 유한사상이다.
스킴의 코호몰로지, 보조정리
\ref{coherent-lemma-characterize-finite}를 보라.

\medskip\noindent
사상 $b:c^*\omega_Y\to\omega_X$를 주의 깊게 구성하자.
$f:X\to\Spec(k)$와 $g:Y\to\Spec(k)$를 구조사상이라 하자. 보조정리
\ref{curves-lemma-duality-dim-1}과 그 증명에 의해
$f^!k=\omega_X[1]$이고 $g^!k=\omega_Y[1]$이다. 그러면
$f^!=c^!\circ g^!$이고 따라서 $c^!\omega_Y=\omega_X$이다.
그러므로 $c^!$의 정의에 의해 연접 $\mathcal{O}_X$-가군 $\mathcal{F}$에
대하여 함자적 동형
$$
\Hom_{D(\mathcal{O}_X)}(\mathcal{F},\omega_X)
\longrightarrow
\Hom_{D(\mathcal{O}_Y)}(Rc_*\mathcal{F},\omega_Y)
$$
이 있다\footnote{스킴의 쌍대성, 보조정리
\ref{duality-lemma-twisted-inverse-image}의 오른쪽 수반함자를
$D^+_\QCoh(\mathcal{O}_Y)$에 제한한 것이다.}. 이 동형은 대각사상의
공단위인 자취사상 $t:Rc_*\omega_X\to\omega_Y$가 유도한다.
사영공식(코호몰로지, 보조정리 \ref{cohomology-lemma-projection-formula})에
의해 표준 사상 $a:\omega_Y\to Rc_*c^*\omega_Y$는 동형이다.
위 사실들을 결합하면
$$
t\circ Rc_*(b)=a^{-1}
$$
을 만족하는 표준 사상 $b:c^*\omega_Y\to\omega_X$가 존재한다. 특히
$b$를 $c^{-1}(Y\setminus\{y_1,\ldots,y_r\})$로 제한하면 동형이다
(가역가군 사이의 사상이고 다른 사상과의 합성이 동형 $a^{-1}$이기
때문이다).

\medskip\noindent
% P07-ERR-0375: frozen source joins the choice and map without punctuation.
(3)에서와 같은 $m\in\mathbf{Z}$를 택하고 사상
$$
b^{\otimes m}:
\Gamma(Y,\omega_Y^{\otimes m})
\longrightarrow
\Gamma(X,\omega_X^{\otimes m})
$$
을 생각하자. $Y$가 환원이고 $b$가 구성 과정에서 언급한 마지막 성질을
가지므로 이 사상은 단사이다. Riemann--Roch (보조정리 \ref{curves-lemma-rr})에
의해
$\chi(X,\omega_X^{\otimes m})=\chi(Y,\omega_Y^{\otimes m})$이다.
따라서
$$
\dim_k\Gamma(Y,\omega_Y^{\otimes m})\geq
\dim_k\Gamma(X,\omega_X^{\otimes m})=
\chi(X,\omega_X^{\otimes m})
$$
이고, $b^{\otimes m}$이 전역단면들 위에서 동형을 유도함을 알 수 있다.
그러므로 $\omega_X^{\otimes m}$의 생성원들이 상에 들어가므로
$b^{\otimes m}:c^*\omega_Y^{\otimes m}\to\omega_X^{\otimes m}$은
전사이다. 따라서 $b^{\otimes m}$은 동형이고, 이에 따라 $b$도 동형이다.
\end{proof}

\begin{lemma}
\label{curves-lemma-characterize-contraction-to-stable}
$k$를 체라 하자. $X$를 $k$ 위의 차원 $1$인 고유 스킴으로서
$H^0(X,\mathcal{O}_X)=k$이고 종수 $g\geq2$인 스킴이라 하자.
$X$의 특이점은 기껏해야 마디점이라고 가정하자. 다음 성질을 갖는
$k$ 위 스킴의 사상
$$
c:X\longrightarrow Y
$$
은 유일한 동형을 제외하고 유일하게 존재한다.
\begin{enumerate}
\item $Y$는 $k$ 위에서 고유이고 $\dim(Y)=1$이며, $Y$의 특이점은
기껏해야 마디점이다.
\item $\mathcal{O}_Y=c_*\mathcal{O}_X$이고 $R^1c_*\mathcal{O}_X=0$이다.
\item $\omega_Y$는 $Y$ 위에서 풍부하다.
\end{enumerate}
\end{lemma}

\begin{proof}
존재성: 이 절의 서론에서 논한 대로 보조정리
\ref{curves-lemma-contracting-rational-tails}와
\ref{curves-lemma-contracting-rational-bridges}를 결합하면 열거된 모든 성질을
갖는 사상이 존재한다. 더구나 이 사상은 합성
$$
X\to X_1\to X_2\ldots\to X_n\to X_{n+1}\to\ldots\to X_{n+n'}
$$
으로 쓸 수 있다. 여기서 처음 $n$개의 사상은 유리 꼬리의 축약이고
마지막 $n'$개의 사상은 유리 다리의 축약이다. 성질 (2)는 유리 꼬리의
각 축약(예 \ref{curves-example-rational-tail})과 유리 다리의 각 축약
(예 \ref{curves-example-rational-bridge})에 대하여 성립한다. 이 성질이
사상의 합성으로 이어짐은 쉽게 알 수 있다.

\medskip\noindent
유일성: $c:X\to Y$를 조건 (1), (2), (3)을 만족하는 사상이라 하자.
사상 $X\to X_{n+n'}$와 $c$에 양립하는 동형
$X_{n+n'}\to Y$가 유일하게 존재함을 보이겠다.

\medskip\noindent
증명을 시작하기 전에 $c$에 관한 몇 가지 사실을 살펴보자. 먼저 스킴의 사상 심화, 정리
\ref{more-morphisms-theorem-stein-factorization-Noetherian}에 의해 $c$의
올들은 기하적으로 연결이다. 특히 $c$는 전사이다. 닫힌점 $y\in Y$에
대하여 올 $X_y$는
$$
H^1(X_y,\mathcal{O}_{X_y})=0
\quad\text{and}\quad
H^0(X_y,\mathcal{O}_{X_y})=\kappa(y)
$$
을 만족한다. 첫째 등식은 스킴의 사상 심화, 보조정리
\ref{more-morphisms-lemma-check-h1-fibre-zero}에서, 둘째 등식은 스킴의 사상 심화, 보조정리
\ref{more-morphisms-lemma-h1-fibre-zero-check-h0-kappa}에서 따른다.
따라서 $X_y=x$이고 $x$가 $y$로 가는 $X$의 유일한 점이면서 $y$와 같은
잉여체를 갖거나, $X_y$가 $\kappa(y)$ 위의 $1$차원 고유 스킴이다.
둘째 경우 $X_y$는 Cohen--Macaulay임을 유의하라
(보조정리 \ref{curves-lemma-automatic}). 그러나 $X$가 환원이므로 $X_y$는 모든
일반점에서 환원이어야 하고(세부사항은 생략한다), 따라서 스킴의 성질,
보조정리 \ref{properties-lemma-criterion-reduced}에 의해 $X_y$는 환원이다.
이에 따라 $X_y$의 특이점은 기껏해야 마디점이다
(보조정리 \ref{curves-lemma-closed-subscheme-nodal-curve}). $X_y$의 종수는 영임을
유의하라(위를 보라). 마지막으로 $X_y$의 차원이 $1$인 점 $y$는 유한
개뿐이다. 이를 $\{y_1,\ldots,y_r\}$라 하자. 그러면 $c$는
$c^{-1}(Y\setminus\{y_1,\ldots,y_r\})$를
$Y\setminus\{y_1,\ldots,y_r\}$로 동형사상한다. 일부 세부사항은
생략한다. 힌트: $\{y_1,\ldots,y_r\}$ 밖에서 $c$는 유한사상이다.
스킴의 코호몰로지, 보조정리
\ref{coherent-lemma-characterize-finite}를 보라.

\medskip\noindent
$C\subset X$를 유리 꼬리라 하자. $c$가 $C$를 한 점으로 보낸다고
주장한다. 모순을 얻기 위해 그렇지 않다고 가정하자. 그러면 $C$의 상은
기약성분 $D\subset Y$이다. $H^0(C,\mathcal{O}_C)=k'$은 $k$의 유한
분리확장이고, $C$는 $k'$-유리점 $x$를 가지며 이 점은 $C$와 $X$의
``나머지'' 부분의 유일한 교점임을 상기하자. 위 일반적 논의로부터
$C\setminus\{x\}\subset c^{-1}(Y\setminus\{y_1,\ldots,y_r\})$가 $D$의
어떤 열린집합 $V$로 동형사상됨을 알 수 있다. $y=c(x)\in D$로 두자.
$y$는 $D$와 $Y$의 ``나머지'' 부분이 만나는 유일한 점이다.
$y\not \in\{y_1,\ldots,y_r\}$이면 $C\cong D$이고, $D$가 $Y$의 유리
꼬리임이 분명하다. 이는 $\omega_Y$의 풍부성과 모순이다
(보조정리 \ref{curves-lemma-rational-tail-negative}). 따라서
$y\in\{y_1,\ldots,y_r\}$이고 $\dim(X_y)=1$이다. 그러면
$x\in X_y\cap C$이고 $x$는 $X_y$와 $C$의 매끄러운 점이다
(보조정리 \ref{curves-lemma-closed-subscheme-nodal-curve}). $y\in D$가 $D$의
특이점이면 $y$는 마디점이고, 그러면 $Y=D$이다(보조정리
\ref{curves-lemma-closed-subscheme-nodal-curve}에 의해 $y$를 지나는 $Y$의 다른
성분은 있을 수 없다). 따라서 $X=X_y\cup C$이고, 이는 예
\ref{curves-example-rational-tail}의 논의에 의해 $X_y$의 종수와 같아서
$g=0$임을 뜻한다. 이는 모순이다. $y\in D$가 $D$의 매끄러운 점이면
$C\to D$는 동형이다(비특이 사영모형은 유일하고 $C,D$가 쌍유리이기
때문이다. 절 \ref{curves-section-curves-function-fields}를 보라). 그러면
$D$는 $Y$의 유리 꼬리이고, 이는 $\omega_Y$의 풍부성과 모순이다.

\medskip\noindent
$n\geq1$이라고 하자. $C\subset X$가 $X\to X_1$에 의해 축약되는
유리 꼬리이면 앞 단락에 의해 $C$는 $Y$의 한 점으로 간다. 따라서
$c:X\to Y$는 $X\to X_1$을 통해 분해된다($X$가 $C$와 $X_1$의
밀어붙임이기 때문이다. 예 \ref{curves-example-rational-tail}의 논의를
보라). $X$를 $X_1$로 바꾸면 $n$이 감소한다. 귀납법에 의해 $n=0$이라고
가정할 수 있다. 즉, $X$는 유리 꼬리를 갖지 않는다.

\medskip\noindent
$n=0$, 즉 $X$가 유리 꼬리를 갖지 않는다고 하자. 보조정리
\ref{curves-lemma-no-rational-tail-semiample-genus-geq-2}에 의해
$\omega_X^{\otimes2}$와 $\omega_X^{\otimes3}$은 전역 생성된다. 보조정리
\ref{curves-lemma-vanishing-twist}에 의해
$H^1(X,\omega_X^{\otimes3})=0$이 따른다. $m=3$으로 보조정리
\ref{curves-lemma-contract-gorenstein-canonical}을 적용하면
$c^*\omega_Y\cong\omega_X$를 얻는다. 또한 보조정리
\ref{curves-lemma-rational-bridge-canonical}과 귀납법에 의해
$\omega_X=(X\to X_{n'})^*\omega_{X_{n'}}$이다. $c$와
$X\to X_{n'}$ 모두에 사영공식을 적용하면 모든 $m$에 대하여
$$
\Gamma(X_{n'},\omega_{X_{n'}}^{\otimes m})=
\Gamma(X,\omega_X^{\otimes m})=
\Gamma(Y,\omega_Y^{\otimes m})
$$
이다. $X_{n'}$와 $Y$는 이 공간들의 직합의 Proj이므로 스킴의 사상, 보조정리
\ref{morphisms-lemma-proper-ample-is-proj}에 의해 원하는 표준 동형
$X_{n'}=Y$를 얻는다. 이것이 도표를 가환하게 하는 유일한 동형임을
확인하는 일은 생략한다.
\end{proof}

\begin{lemma}
\label{curves-lemma-tricanonical}
$k$를 체라 하자. $X$를 $k$ 위의 차원 $1$인 고유 스킴으로서
$H^0(X,\mathcal{O}_X)=k$이고 종수 $g\geq2$인 스킴이라 하자.
$X$의 특이점은 기껏해야 마디점이고 $\omega_X$는 풍부하다고 가정하자.
그러면 $\omega_X^{\otimes3}$은 매우 풍부하고
$H^1(X,\omega_X^{\otimes3})=0$이다.
\end{lemma}

\begin{proof}
다양체, 보조정리
\ref{varieties-lemma-ampleness-in-terms-of-degrees-components}와 보조정리
\ref{curves-lemma-rational-tail-negative} 및 \ref{curves-lemma-rational-bridge-zero}를
결합하면 $X$가 유리 꼬리와 유리 다리를 갖지 않음을 알 수 있다.
그러면 보조정리 \ref{curves-lemma-contracting-rational-tails}에 의해
$\omega_X^{\otimes3}$은 전역 생성된다.
$H^0(X,\omega_X^{\otimes3})$의 $k$-기저 $s_0,\ldots,s_n$을 택하자.
그러면 사상
$$
\varphi_{\omega_X^{\otimes3},(s_0,\ldots,s_n)}:
X\longrightarrow\mathbf{P}^n_k
$$
을 얻는다. 구성, 절
\ref{constructions-section-projective-space}을 보라. 보조정리의 주장은 이
사상이 닫힌몰입이라는 것이다. 이를 확인하기 위해 $k$를 그 대수적
폐포로 바꾸어도 된다. 내림, 보조정리
\ref{descent-lemma-descending-property-closed-immersion}을 보라. 따라서
$k$가 대수적으로 닫혀 있다고 가정해도 된다.

\medskip\noindent
$k$가 대수적으로 닫혀 있다고 가정하자. 다양체, 보조정리
\ref{varieties-lemma-variant-separate-points-tangent-vectors}를 사용하여
보조정리를 증명하겠다.
% P07-ERR-0376: preserve frozen self-referential degree “over Z”.
$Z\subset X$를 $Z$ 위 차수 $2$인 닫힌부분스킴으로서 아이디얼층
$\mathcal{I}\subset\mathcal{O}_X$를 갖는 부분스킴이라 하자. 사상
$$
H^0(X,\mathcal{L})\to H^0(Z,\mathcal{L}|_Z)
$$
가 전사임을 보여야 한다. 따라서 $H^1(X,\mathcal{I}\mathcal{L})=0$임을
보이면 충분하다. 이를 위해 보조정리 \ref{curves-lemma-vanishing-on-gorenstein}을
사용할 것이다. 그러므로 모든 환원 연결 닫힌부분스킴 $Y\subset X$에
대하여
$$
3\deg(\omega_X|_Y)>-2\chi(Y,\mathcal{O}_Y)+\deg(Z\cap Y)
$$
임을 보이면 충분하다. $k$가 대수적으로 닫혀 있고 $Y$가 연결이고
환원이므로
$H^0(Y,\mathcal{O}_Y)=k$이다(다양체, 보조정리
\ref{varieties-lemma-proper-geometrically-reduced-global-sections}). 따라서
$\chi(Y,\mathcal{O}_Y)=1-\dim H^1(Y,\mathcal{O}_Y)$이다. 그러므로
$$
3\deg(\omega_X|_Y)>-2+2\dim H^1(Y,\mathcal{O}_Y)+\deg(Z\cap Y)
$$
임을 보여야 한다. 이는 보조정리 \ref{curves-lemma-no-rational-tail}에 의해
$Y=X$이거나 $\deg(\omega_X|_Y)=0$인 경우를 제외하면 성립한다.
$\omega_X$가 풍부하므로 둘째 가능성은 일어나지 않는다(이 증명에서
처음 인용한 보조정리를 보라). 마지막으로 $Y=X$이면 Riemann--Roch
(보조정리 \ref{curves-lemma-rr})와 $g\geq2$라는 사실을 사용하면
% P07-ERR-0377: frozen source misspells inequality; translated as intended.
부등식이 성립함을 알 수 있다. $Z=\emptyset$으로 같은 논증을 적용하면
$H^1(X,\omega_X^{\otimes3})=0$을 얻는다.
\end{proof}






\section{벡터장}
\label{curves-section-vector-fields}

\noindent
이 절에서는 곡선 위의 벡터장 공간을 연구한다. 벡터장은 무한소
자기동형과 대응하므로(스킴의 사상 심화, 절
\ref{more-morphisms-section-action-by-derivations}), 모듈라이 이론에서
중요한 역할을 한다.

\medskip\noindent
$k$를 대수적으로 닫힌 체라 하자. $X$를 $k$ 위 유한형 스킴이라 하고,
$x\in X$를 닫힌점이라 하자. $\mathcal{I}\subset\mathcal{O}_X$가 $x$의
아이디얼층일 때 $D(\mathcal{I})\subset\mathcal{I}$이면
$D\in\text{Der}_k(\mathcal{O}_X,\mathcal{O}_X)$가 {\it $x$를 고정한다}고
한다.

\begin{lemma}
\label{curves-lemma-smooth-vector-fields}
$k$를 대수적으로 닫힌 체라 하자. $X$를 $k$ 위의 매끄럽고 고유이며
연결인 곡선이라 하고, $g$를 $X$의 종수라 하자.
\begin{enumerate}
\item $g\geq2$이면 $\text{Der}_k(\mathcal{O}_X,\mathcal{O}_X)$는 영이다.
\item $g=1$이고 $D\in\text{Der}_k(\mathcal{O}_X,\mathcal{O}_X)$가 영이
아니면 $D$는 $X$의 어떤 닫힌점도 고정하지 않는다.
\item $g=0$이고 $D\in\text{Der}_k(\mathcal{O}_X,\mathcal{O}_X)$가 영이
아니면 $D$는 $X$의 닫힌점을 기껏해야 $2$개 고정한다.
\end{enumerate}
\end{lemma}

\begin{proof}
보편 $k$-미분 $d:\mathcal{O}_X\to\Omega_{X/k}$가 있고, 따라서 어떤
$\mathcal{O}_X$-선형사상 $\theta:\Omega_{X/k}\to\mathcal{O}_X$에 대하여
$D=\theta\circ d$임을 상기하자. 보조정리 \ref{curves-lemma-duality-dim-1}에 의해
$\Omega_{X/k}\cong\omega_X$이다. Riemann--Roch에 의해
$\deg(\omega_X)=2g-2$이다(보조정리 \ref{curves-lemma-rr}). 따라서 $g>1$이면
다양체, 보조정리 \ref{varieties-lemma-check-invertible-sheaf-trivial}에
의해 $\theta$는 영일 수밖에 없다. 이로써 (1)을 증명했다. $g=1$이면
영이 아닌 $\theta$는 어느 곳에서도 영이 되지 않고, $g=0$이면 영이 아닌
$\theta$는 차수 $2$인 인자에서 영이 된다. 그러므로 닫힌점 $x\in X$에서
$\theta$가 영인 것과 위 정의의 의미에서 $D$가 $x$를 고정하는 것이
동치임을 보이면 (2)와 (3)이 따른다. $z\in\mathcal{O}_{X,x}$를 균등화원이라
하자. 그러면 보조정리 \ref{curves-lemma-uniformizer-works}에 의해 $dz$는
$\Omega_{X,x}$의 기저원소이다. $D(z)=\theta(dz)$이므로 결론을 얻는다.
\end{proof}

\begin{lemma}
\label{curves-lemma-nodal-vector-fields}
$k$를 대수적으로 닫힌 체라 하자. $X$를 기껏해야 마디인, 고유이고 연결인
$1$차원 $k$-스킴이라 하자. $\nu:X^\nu\to X$를 정규화라 하고,
$S\subset X^\nu$를 $\nu$가 동형이 아닌 점들의 집합이라 하자. 그러면
$$
\text{Der}_k(\mathcal{O}_X,\mathcal{O}_X)=
\{D'\in\text{Der}_k(\mathcal{O}_{X^\nu},\mathcal{O}_{X^\nu})\mid
D'\text{ fixes every }x^\nu\in S\}
$$
이다.
\end{lemma}

\begin{proof}
$x\in X$를 마디점이라 하자. $x',x''\in X^\nu$를 $x$의 역상이라 하자.
($k$가 대수적으로 닫혀 있으므로 모든 마디점은 분할 마디점이다.
% P07-ERR-0378: frozen source misspells algebraically; translated as intended.
정의 \ref{curves-definition-split-node}과 보조정리 \ref{curves-lemma-split-node}을
보라.) $u\in\mathcal{O}_{X^\nu,x'}$와 $v\in\mathcal{O}_{X^\nu,x''}$를
균등화원이라 하자. 완전열
$$
0\to\mathcal{O}_{X,x}\to
\mathcal{O}_{X^\nu,x'}\times\mathcal{O}_{X^\nu,x''}\to k\to0
$$
이 있음을 유의하라. 이는 보조정리 \ref{curves-lemma-multicross}에서 따른다.
따라서 $u,v$를 $uv=0$인 $\mathcal{O}_{X,x}$의 원소로 볼 수 있다.

\medskip\noindent
$D\in\text{Der}_k(\mathcal{O}_X,\mathcal{O}_X)$라 하자. 그러면
$0=D(uv)=vD(u)+uD(v)$이다. $\mathcal{O}_{X,x}$에서 $(u)$는 $v$의
소멸자이고 그 역도 성립하므로 $D(u)\in(u)$이고 $D(v)\in(v)$이다.
$\mathcal{O}_{X^\nu,x'}=k+(u)$이므로 $D$를
$\mathcal{O}_{X^\nu,x'}$로 확장할 수 있고, 더구나 그 확장은 $x'$을
고정한다. 이로부터 등식 오른쪽의 $D'$을 얻는다. 역으로 $x'$과 $x''$을
고정하는 $D'$이 주어지면 $D'$은 부분환
$\mathcal{O}_{X,x}\subset
\mathcal{O}_{X^\nu,x'}\times\mathcal{O}_{X^\nu,x''}$을 보존하며, 이것이
등식의 오른쪽에서 왼쪽으로 가는 방법이다.
\end{proof}

\begin{lemma}
\label{curves-lemma-stable-vector-fields}
$k$를 대수적으로 닫힌 체라 하자. $X$를 기껏해야 마디인, 고유이고 연결인
$1$차원 $k$-스킴이라 하자. $X$의 종수는 적어도 $2$이고 $X$는 유리
꼬리와 유리 다리를 갖지 않는다고 가정하자. 그러면
$\text{Der}_k(\mathcal{O}_X,\mathcal{O}_X)=0$이다.
\end{lemma}

\begin{proof}
$D\in\text{Der}_k(\mathcal{O}_X,\mathcal{O}_X)$라 하자. $X^\nu$를 $X$의
정규화라 하고, $D'\in\text{Der}_k(\mathcal{O}_{X^\nu},
\mathcal{O}_{X^\nu})$를 보조정리 \ref{curves-lemma-nodal-vector-fields}를 통해 $D$에
대응하는 원소라 하자. $C\subset X^\nu$를 기약성분이라 하자. $C$의
종수가 $>1$이면 보조정리 \ref{curves-lemma-smooth-vector-fields}의 (1)에 의해
$D'|_{\mathcal{O}_C}=0$이다. $C$의 종수가 $1$이면 $X$의 마디점으로 가는
$C$의 닫힌점 $c$가 적어도 하나 있다(그렇지 않으면 $X\cong C$의 종수가
$1$이다). 위 대응에 의해 이는 $D'|_{\mathcal{O}_C}$가 $c$를 고정함을
뜻하므로 보조정리 \ref{curves-lemma-smooth-vector-fields}의 (2)에 의해 영이다.
마지막으로 $C$의 종수가 영이면 $X$의 마디점으로 가는 서로 다른 닫힌점
$c_1,c_2,c_3\in C$가 적어도 $3$개 있다. 그렇지 않으면 $X$가 두 점을
붙인 $C$($C$의 그 두 점이 같은 마디점으로 가는 경우)이거나, $C$가 유리
다리(두 점이 $X$의 서로 다른 마디점으로 가는 경우)이거나, $C$가 유리
꼬리(한 점이 $X$의 마디점으로 가는 경우)이다.
% P07-ERR-0379: preserve frozen subject C in the inconsistent genus claim.
이 세 가능성은 $C$의 종수가 $\geq2$이고 유리 다리나 유리 꼬리를 갖지
않으므로 허용되지 않는다. 따라서 $D'|_{\mathcal{O}_C}$는
$c_1,c_2,c_3$을 고정하고, 보조정리 \ref{curves-lemma-smooth-vector-fields}의 (3)에
의해 영이다.
\end{proof}
